%! Author = Omar Iskandarani
%! Date = 4/10/2026
%! Affiliation = Independent Researcher, Groningen, The Netherlands
%! ORCID = 0009-0006-1686-3961
%! DOI = 10.5281/zenodo.21249056
%! v0.8.20 edition: chronos kelvin + contact pressure + rank9 + SSDL + particle/clock ansatz (been_processed).

\newcommand{\paperdoi}{10.5281/zenodo.21249056}
\newcommand{\papertitle}{Swirl-String-Theory Canon-v0.8.20}
\newcommand{\paperkeywords}{Swirl String Theory, SST, Canon, Topological Hydrodynamics, Knot Theory, Relational Time, Delay Dynamics, Epistemic Classification, chronos kelvin hitting conditions, contact pressure saturation, rank nine contact channels, SSDL monopole DtN}

\documentclass[11pt]{article}
% ---------------------------
% Packages
% ---------------------------
\usepackage{amsmath,amssymb,amsfonts,bm,amsthm}
\usepackage{mathrsfs}
\usepackage{siunitx}
\usepackage[a4paper,margin=1in]{geometry}
\usepackage[absolute,overlay]{textpos}
\usepackage[T1]{fontenc}
\usepackage[utf8]{inputenc}
\usepackage{graphicx}
\usepackage{physics}
\usepackage{mathtools}
\usepackage{array}
\usepackage{tabularx}
\usepackage{import}
\usepackage{subfiles}
\usepackage{lmodern}
\usepackage{booktabs}
\usepackage{tikz}
\usetikzlibrary{arrows.meta,positioning,calc,fit,decorations.pathmorphing}
% Tables and Figures
\usepackage{float}

% ---------------------------
% Theorem environments
% ---------------------------
\theoremstyle{definition}
\newtheorem{definition}{Definition}[section]

\theoremstyle{plain}
\newtheorem{axiom}{Axiom}[section]
\newtheorem{theorem}{Theorem}[section]
\newtheorem{lemma}{Lemma}[section]

\theoremstyle{remark}
\newtheorem{remark}{Remark}[section]

% hyperref: load after float/tikz/amsthm; *before* tocloft (required by tocloft.pdf — otherwise TOC can be empty).
\usepackage{hyperref}
\hypersetup{colorlinks=true,linkcolor=blue,citecolor=blue,urlcolor=blue}

% TOC (tocloft must follow hyperref)
\usepackage{tocloft}
\setcounter{tocdepth}{2}
\renewcommand{\cftsecfont}{\footnotesize}
\renewcommand{\cftsubsecfont}{\footnotesize\itshape}
\renewcommand{\cftsecleader}{\cftdotfill{5}}
\newcommand{\compacttoc}{%
    \begingroup
    \setlength{\cftparskip}{0pt}%
    \setlength{\cftbeforesecskip}{0.5pt}%
    \setlength{\cftbeforesubsecskip}{0.5pt}%
    \renewcommand{\cftsecfont}{\scriptsize}%
    \renewcommand{\cftsubsecfont}{\scriptsize\itshape}%
    \renewcommand{\cftsecpagefont}{\scriptsize}%
    \renewcommand{\cftsubsecpagefont}{\scriptsize}%
    \setlength{\cftsecnumwidth}{2.8em}%
    \setlength{\cftsubsecindent}{1.4em}%
    \renewcommand{\cftdotsep}{2}%
    \setlength{\parskip}{0pt}%
    \linespread{0.80}\selectfont
    \tableofcontents
    \endgroup
}

% ---------------------------
% SST macro prelude
% ---------------------------
\newcommand{\swirlarrow}{\mathchoice{\mkern-2mu\scriptstyle\boldsymbol{\circlearrowleft}}{\mkern-2mu\scriptstyle\boldsymbol{\circlearrowleft}}{\mkern-2mu\scriptscriptstyle\boldsymbol{\circlearrowleft}}{\mkern-2mu\scriptscriptstyle\boldsymbol{\circlearrowleft}}}
\newcommand{\swirlarrowcw}{\mathchoice{\mkern-2mu\scriptstyle\boldsymbol{\circlearrowright}}{\mkern-2mu\scriptstyle\boldsymbol{\circlearrowright}}{\mkern-2mu\scriptscriptstyle\boldsymbol{\circlearrowright}}{\mkern-2mu\scriptscriptstyle\boldsymbol{\circlearrowright}}}
\newcommand{\vswirl}{\mathbf{v}_{\swirlarrow}}
\newcommand{\vswirlcw}{\mathbf{v}_{\swirlarrowcw}}
\newcommand{\rhoF}{{\rho_{\!f}}}
\newcommand{\rhoE}{\rho_{\!E}}
\newcommand{\rhoM}{\rho_{\!m}}
\newcommand{\rhocore}{\rho_{\mathrm{core}}}
\newcommand{\rhohorn}{\rho_{\mathrm{horn}}^{\mathrm{eff}}}
\newcommand{\SwirlClock}{{S_{(t)}^{\swirlarrow}}}
\newcommand{\SwirlClockcw}{{S_{(t)}^{\swirlarrowcw}}}
\newcommand{\rc}{{r_c}}
\newcommand{\vnorm}{\lVert \mathbf{v}_{\!\boldsymbol{\circlearrowleft}} \rVert}
\newcommand{\vchar}{v_{\!\boldsymbol{\circlearrowleft}}^{\ast}}
\newcommand{\uswirl}{\mathbf{u}_{\!\boldsymbol{\circlearrowleft}}}
\newcommand{\omegaswirl}{\boldsymbol{\omega}_{\!\boldsymbol{\circlearrowleft}}}
\newcommand{\tauCirc}{\tau_{\rm circ}}
\newcommand{\tauCircC}{\tau_{\rm circ,c}}


\begin{document}
    
    \thispagestyle{empty}

    % -------------------------------------------------
    % Title page / front matter
    % -------------------------------------------------
    \begin{center}
        \vspace*{0.06\textheight}

        {\LARGE\bfseries \papertitle\par}
        \vspace{0.45cm}

        {\large Canonical Reference and Research Framework\par}
        \vspace{1.0cm}

        {\Large Omar Iskandarani\textsuperscript{*}\par}
        \vspace{0.35cm}

        {\normalsize Independent Researcher\par}
        {\normalsize Groningen, The Netherlands\par}
        \vspace{0.6cm}

        {\normalsize Version v0.8.20\par}
        {\normalsize \today\par}
        \vspace{0.9cm}

        \rule{0.82\textwidth}{0.5pt}
        \vspace{0.7cm}
    \end{center}

    \begin{abstract}
        \noindent
        This document presents Swirl-String Theory (SST) as a structured canonical synthesis rather than as a claim that every phenomenological sector has already been derived to theorem level. The text consolidates the primitive constant chain linking the effective fluid density, characteristic swirl speed, Compton anchoring, and horn/circulation scale; records delay-induced mode selection as the preferred non-operator route to spectral discreteness; and organizes the atomic bridge, spectroscopic filter, relational-time sector, and topology-based mass program within one explicitly stratified framework. Throughout, orthodox inputs, SST-internal derivations, bridge-level constructions, and conjectural extensions are kept distinct so that the canon remains internally coherent, externally comparable, and open to falsification.
    \end{abstract}

    \vspace{0.2cm}
    \noindent\textbf{Keywords:} \paperkeywords

    \vspace{0.5cm}
    \noindent\rule{\textwidth}{0.4pt}
    \vspace{0.35cm}
    
    \noindent\begin{quote}
        \small
        \textbf{Canon reading guide.}
        
        \vspace{0.3em}
        \begin{tabular}{@{}>{\raggedleft\arraybackslash}p{0.21\textwidth}@{\hspace{0.32em}}>{\raggedright\arraybackslash}p{0.75\textwidth}@{}}
            \textbf{[ORTHODOX]} & established physics or mathematics. \\[0.25em]
            \textbf{[DERIVED]} & consequences obtained from definitions and assumptions of this SST layer. \\[0.25em]
            \textbf{[SPECULATIVE]} & conjectural extensions or identifications not yet closed as theorems. \\
        \end{tabular}
        
        \vspace{0.4em}
        Optional \emph{role} labels (\emph{e.g.}, \emph{Definition}, \emph{Benchmark}, \emph{Constraint}) qualify how to read a statement but do not replace the primary tag.
        Orthodox constraints---including quantum predictions, atomic spectroscopy, and relativistic consistency---are not overridden: bridge constructions, benchmark sections, and research-track material remain subordinate to the consistency stratification introduced later in the canon.
    \end{quote}

    \medskip

    \begin{textblock*}{1\textwidth}(0.15\textwidth,1.1\textheight)\raggedright\footnotesize\renewcommand{\arraystretch}{1.0} \noindent\rule{\textwidth}{0.4pt}\\[0.5em]
        \textsuperscript{\textbf{*}} Email: \texttt{info@omariskandarani.com}\\
        ORCID: \texttt{\href{https://orcid.org/0009-0006-1686-3961}{0009-0006-1686-3961}}\\
        DOI: \href{https://doi.org/\paperdoi}{\paperdoi}
    \end{textblock*}
    
    \newpage
    % Original: \tableofcontents
    % Use compact version to save space
    \compacttoc
    \newpage


% ===========================
% CANON STRUCTURE (v0.8.1)
% ===========================

% ---------- PRELUDE ----------
    \section{Philosophical Prelude}
        \label{sec:philosophy}

    \subsection{Why a philosophical prelude is necessary}
        The SST canon is not only a collection of equations but also a statement about what counts as fundamental. Standard quantum theory and general relativity are highly successful formalisms, yet they begin from different ontological primitives: operator-valued states in one case and spacetime geometry in the other. SST instead adopts a single material substrate as the underlying level of description and treats particles, radiation, and clock structure as organized states of that substrate.

    \subsection{The ontological shift}
        The canon therefore makes the following conceptual replacement:
        \begin{align}
            \text{point particles} &\longrightarrow \text{stable swirl-string configurations}, \\
            \text{vacuum as background} &\longrightarrow \text{vacuum as structured medium}, \\
            \text{quantization as postulate} &\longrightarrow \text{discreteness from dynamics and topology}.
        \end{align}
        This is not a denial of effective particle language. Rather, it is a claim that particle language is secondary and emerges from a more primitive hydrodynamic and topological layer.

    \subsection{Relation to orthodox physics}
        SST is not introduced as a replacement for the empirical content of orthodox physics. Its canonical role is instead threefold:
        \begin{enumerate}
            \item to reproduce known low-energy structures in the appropriate limits,
            \item to identify which apparently fundamental constants or rules may be structurally redundant,
            \item to supply a dynamical mechanism for phenomena that are otherwise postulated axiomatically.
        \end{enumerate}
        In this sense the canon follows a \emph{Rosetta principle}: whenever a standard description already works, SST must map to it before claiming any deeper reinterpretation.

    \subsection{Why delay and topology matter}
        Two themes recur across the development of the canon.
        First, finite circulation delay on closed carriers provides a natural source of temporal nonlocality. Second, topological invariants constrain admissible configurations and stabilize structures that would otherwise collapse or disperse. Together they motivate the canonical SST slogan
        \begin{center}
            \textit{delay selects modes, topology protects them.}
        \end{center}
        This is the philosophical reason why the canon treats quantization neither as a primitive operator rule nor as a purely statistical statement.

    \subsection{Scope of the present canon}
        The present version of the canon should be read as a structured reference document rather than a claim that every sector has already been derived to theorem level. Some components are mature and reusable, such as the primitive constant chain, circulation quantization, and the Swirl-Clock. Other components are deliberately retained as research-track layers, such as full gauge emergence, topology-to-charge assignments, and complete many-body atomic closure. The point of the canon is therefore not maximal rhetorical closure, but explicit logical organization.

    \subsection{Interpretive summary}
        The philosophical stance of SST can be summarized in one sentence:
        \begin{center}
            \textit{Physics is interpreted as the dynamics of a continuous medium whose stable, delayed, and topologically organized motions appear effectively as particles, clocks, fields, and spectra.}
        \end{center}

% ---------- FOUNDATIONS ----------
\section{Formal Foundations}
    \label{sec:foundations}
    
    \subsection{Primitive Structure of the Theory}
        
        We define the Swirl-String Theory (SST) as a continuum-based framework built on a minimal set of primitive quantities.
        
        \textbf{Primitive calibration set:}
        \begin{align}
            \mathcal{P}_{\mathrm{cal}} = \{ \rhoF, \vchar, \omega_c \}
        \end{align}
        
        where:
        \begin{itemize}
            \item $\rhoF$ is the effective background fluid density,
            \item $\vchar=\lVert\mathbf{v}_{\!\boldsymbol{\circlearrowleft}}\rVert$ is the canonical characteristic swirl-speed magnitude,
            \item $\omega_c$ is the Compton angular frequency.
        \end{itemize}

        \textbf{[CANON / NOTATION]}: $\vchar$ is a calibrated scalar speed. The local swirl-velocity field is denoted
        \begin{align}
            \uswirl(\mathbf{x},t),
            \qquad
            \omegaswirl(\mathbf{x},t)=\nabla\times\uswirl(\mathbf{x},t).
        \end{align}
        The symbol $\mathbf{v}_{\!\boldsymbol{\circlearrowleft}}$ without explicit arguments is reserved for the canonical characteristic carrier velocity; it must not be used as the unrestricted local field in bridge equations.
        
        \textbf{[ORTHODOX]}: $\omega_c = \frac{m_e c^2}{\hbar}$ is a standard physical constant.
        
        \textbf{[SPECULATIVE]}: The identification of $\omega_c$ as the primitive Compton phase frequency of the SST carrier is an SST postulate; it is not a local vorticity variable.
    
    \subsection{Derived Horn/Circulation Radius}
        \label{subsec:derived_horn_circulation_radius}

        The quantity previously denoted as the characteristic ``core radius'' is fixed more precisely as a
        Compton-locked horn/circulation radius:
        \begin{align}
            \rc \equiv R_{\mathrm{horn}}
            = \frac{\vchar}{\omega_c}
            = \frac{\hbar\vchar}{m_e c^2}
            = \frac{\alpha\hbar}{2m_e c}
            = \frac{r_e}{2}.
            \label{eq:rc_definition}
        \end{align}
        Here
        \begin{align}
            \omega_c \equiv \frac{m_e c^2}{\hbar}
            \label{eq:omega_c_compton}
        \end{align}
        is the electron Compton angular frequency. Thus $\rc$ is the radius at which the canonical
        tangential swirl speed $\vchar$ is obtained at the Compton angular frequency.

        \textbf{[CANON / NOTATION]}: $\omega_c$ denotes the Compton angular frequency. It must not be confused with the local vorticity scale. If $\omega_{\mathrm{vort},c}$ denotes the local vorticity of a rigid local rotation, then
        \begin{align}
            \omega_{\mathrm{vort},c}=2\omega_c,
            \qquad
            R_{\mathrm{horn}}=\frac{2\vchar}{\omega_{\mathrm{vort},c}}.
        \end{align}

        \textbf{[CANON / SEMANTIC RULE]}: $\rc$ is not the resolved physical vortex-tube radius. The local tube/core thickness is denoted by $a_{\mathrm{core}}$ and remains a separate geometric or variational quantity:
        \begin{align}
            a_{\mathrm{core}} \neq \rc.
        \end{align}

        \textbf{[DERIVED] Consequence}: Eq.~\eqref{eq:rc_definition} removes $\rc$ as an independent parameter while avoiding the earlier ambiguity between horn-envelope scale, tube radius, and vorticity.

        \textbf{[CALIBRATED CHAIN GUARD]}: Eq.~\eqref{eq:rc_definition} is a Compton--electron-radius closure inside the primitive calibration chain. It must not be cited as an independent hydrodynamic derivation of either $\alpha$, $\vchar$, or $\rc$. A future fluid-dynamical derivation would have to obtain the same radius from a resolved vortex-tube variational problem using $a_{\mathrm{core}}$ or an explicit profile, without inserting the electron-radius identity.
    
    \subsection{Circulation Quantization}
        
        The circulation is defined as:
        
        \begin{align}
            \Gamma = \oint \mathbf{v} \cdot d\ell
        \end{align}
        
        We introduce the fundamental circulation quantum:
        
        \begin{align}
            \Gamma_0 = 2\pi \rc \, \vchar
            \label{eq:gamma0}
        \end{align}
        
        and impose:
        
        \begin{align}
            \Gamma = n \Gamma_0, \quad n \in \mathbb{Z}
        \end{align}
        
        \textbf{[ORTHODOX]}: Circulation quantization appears in superfluid systems.
        
        \textbf{[DERIVED]}: In SST, $\Gamma_0$ is no longer an input but follows from Eq.~\ref{eq:rc_definition}.
    
    \subsection{Energy and Mass Relation}
        
        Mass is interpreted as stored rotational energy:
        
        \begin{align}
            E = m c^2
        \end{align}
        
        \textbf{[ORTHODOX]}: Standard relativistic energy relation.
        
        \textbf{[SPECULATIVE]}: SST interprets this energy as localized vortex kinetic energy.
    
    \subsection{Swirl Velocity Constraint}
        
        We define the dimensionless ratio:
        
        \begin{align}
            \eta_0 = \frac{\vchar}{c}
        \end{align}
        
        Within the present calibrated constant chain:
        
        \begin{align}
            \eta_0 = \frac{\alpha}{2}
        \end{align}
        
        \textbf{[CALIBRATED ALGEBRAIC IDENTITY]}: Once the Compton anchoring relation $\rc=\vchar/\omega_c$ is combined with the electron-radius closure $\rc=\alpha\hbar/(2m_ec)$, the identity $\vchar=\alpha c/2$ follows by direct substitution.

        \textbf{[CRITICAL NOTE]} This is not an independent fluid-dynamical derivation of $\alpha$ or of the characteristic speed. The fine-structure constant enters through the adopted electron-radius / horn-radius calibration, so the result is a calibrated algebraic identity rather than an emergent theorem.
    
    \subsection{Swirl-Clock Definition}
        
        We introduce the Swirl-Clock as a local scalar field encoding relativistic time dilation:
        
        \begin{align}
            \SwirlClock = \sqrt{1 - \frac{\vchar{}^2}{c^2}}
            \label{eq:swirl_clock}
        \end{align}
        
        \textbf{[ORTHODOX] Functional form}: This matches the standard Lorentz time-dilation factor.

        \textbf{[SPECULATIVE] Identification}: In SST, the relevant velocity is interpreted as a local medium-kinematic state rather than only a standard worldline speed.

        \textbf{[DERIVED] Consequence}: Once this identification is adopted, the Swirl-Clock provides a local scalar measure of kinematic time modulation within the medium description.

    \subsection{Two-Speed Discipline for Clock Emergence}
        \label{subsec:two_speed_clock_discipline}

        The canon distinguishes the canonical horn/circulation speed \(\vchar\) from the transverse causal speed \(c\):
        \begin{align}
            \vchar = 1.09384563\times 10^{6}\ \mathrm{m\,s^{-1}},
            \qquad
            c = 2.99792458\times 10^{8}\ \mathrm{m\,s^{-1}},
            \qquad
            \frac{\vchar}{c}=3.64867628\times10^{-3}.
        \end{align}
        Thus \(\vchar\) is a calibrated circulation/horn speed, not the universal causal speed.

        \textbf{[CANON / SEMANTIC RULE]} The calibrated identity \(\vchar=\alpha c/2\) must not be inverted into a claim that the internal core sound speed equals \(c\).  In a Gross--Pitaevskii / NLSE core model with circulation quantum and healing length
        \begin{align}
            \Gamma_q = \frac{h_*}{m_*},
            \qquad
            \xi = \frac{\hbar_*}{\sqrt{2}\,m_* c_s},
        \end{align}
        imposing \(\Gamma_q=\Gamma_0\) and \(\xi=\rc\) gives
        \begin{align}
            c_s
            =
            \frac{\Gamma_0}{2\sqrt{2}\pi\rc}
            =
            \frac{\vchar}{\sqrt{2}}
            =
            7.73465663\times10^{5}\ \mathrm{m\,s^{-1}}.
            \label{eq:core_sound_speed_vchar_over_sqrt2}
        \end{align}
        Therefore a calculation that closes an inertia ratio as \(u^2/c_s^2\) has produced an internal core-acoustic closure, not yet the Lorentz-compatible ratio \(u^2/c^2\).

        \textbf{[RESEARCH-TRACK GATE]} A Lorentz-compatible derivation of the Swirl-Clock requires a separate transverse torsion/shear sector with propagation speed \(c_T=c\), plus a core--torsion impedance identity of the form
        \begin{align}
            \mathsf M_{\mathrm{torsion}}[K]
            \stackrel{?}{=}
            \frac{2E_0[K]}{c_T^2}\,\mathsf I .
            \label{eq:core_torsion_impedance_gate_main}
        \end{align}
        Until Eq.~\eqref{eq:core_torsion_impedance_gate_main} is derived from a controlled field calculation, the Swirl-Clock remains an orthodox functional form with an SST interpretation, or a derived-conditional bridge result, rather than a hard theorem from the NLSE core alone.

    \subsection{Density Hierarchy}
        
        We distinguish the following densities:
        
        \begin{align}
            \rhoF &\quad \text{(effective background fluid density)}, \\
            \rhoE &= \tfrac12 \rhoF \, \vchar{}^2 \quad \text{(swirl energy density)}, \\
            \rhoM &= \rhoE / c^2 \quad \text{(mass-equivalent density)}, \\
            \rhocore &\quad \text{(canonical saturated envelope-equivalent density)}.
        \end{align}

        \textbf{[ORTHODOX] Definition}: $\rhoM=\rhoE/c^2$ is the mass-equivalent density associated with the swirl kinetic accounting above; in orthodox limits, total energy $E$ in a region still yields mass density $E/c^2$.

        \textbf{[SPECULATIVE] Interpretation}: The distinction between background, energetic, and saturated envelope-equivalent sectors is elevated in SST to a structural ontology of the medium.
    
    \subsection{Horn-Envelope Density Closure}
        
        The density formerly denoted as a ``core density'' is now interpreted as an
        effective horn-envelope normalization density, not as a measured local
        material density of the resolved vortex tube:
        
        \begin{align}
            \rhohorn
            \equiv
            \frac{m_e c^2}{2\pi \vchar{}^2 R_{\mathrm{horn}}^3}
            =
            \frac{m_e c^2}{2\pi \vchar{}^2 \rc^3}.
            \label{eq:core_density}
        \end{align}
        
        \textbf{[CANON / SEMANTIC RULE]}: Equation~\eqref{eq:core_density}
        defines the canonical horn-envelope density closure used by the present
        SST calibration program. Legacy appearances of $\rhocore$ in expressions
        proportional to $\rhocore\rc^3$ should be read as $\rhohorn$ unless a
        separate resolved-tube model is explicitly introduced.

        \textbf{[DERIVED] Consequence}: Within that closure, $\rhohorn$ is no
        longer an independent parameter. If a local tube mass density is required,
        it must be introduced through the separate tube radius $a_{\rm core}$ and
        a corresponding local tube-volume or profile model.


    \subsection{Chronos--Kelvin Invariant and Clock--Radius Transport}
        \label{subsec:chronos_kelvin}

        For a thin material loop of radius $R(t)$ carried by an incompressible, inviscid flow with no reconnection events, Kelvin's theorem implies conservation of circulation. In the near-solid-body core approximation this gives

        \begin{align}
            \frac{D}{Dt}\!\bigl(R^2 \omega\bigr) = 0,
            \label{eq:chronos_kelvin_basic}
        \end{align}

        where $\omega$ is the characteristic local vorticity on the loop. For rigid local rotation, $v_t=(\omega/2)R$; evaluated at the canonical horn scale this gives $v_t=(\omega/2)\rc$. Using this relation together with Eq.~\eqref{eq:swirl_clock}, one may rewrite

        \begin{align}
            \omega = \frac{2c}{\rc}\sqrt{1-S_{(t)}^2},
        \end{align}

        so that Eq.~\eqref{eq:chronos_kelvin_basic} takes the equivalent form

        \begin{align}
            \frac{D}{Dt}\left(\frac{2c}{\rc} R^2 \sqrt{1-S_{(t)}^2}\right) = 0.
            \label{eq:chronos_kelvin_clock}
        \end{align}

        \textbf{[ORTHODOX]}: Eq.~\eqref{eq:chronos_kelvin_basic} is the thin-loop form of Kelvin circulation conservation in inviscid barotropic flow \cite{Batchelor1967,Saffman1992}.

        \textbf{[CONDITIONAL DERIVED]}: Eq.~\eqref{eq:chronos_kelvin_clock} is the SST rewriting of the same invariant in terms of the Swirl-Clock, conditional on evaluating the local tangential speed at the fixed horn/circulation scale $\rc$. It is not a generic radius-transport theorem for an arbitrary resolved tube radius $a_{\rm core}(t)$.

        Differentiating Eq.~\eqref{eq:chronos_kelvin_clock} gives a useful transport identity:

        \begin{align}
            \frac{dS_{(t)}}{dt}
            =
            \frac{2\bigl(1-S_{(t)}^2\bigr)}{S_{(t)}}
            \frac{1}{R}\frac{dR}{dt}.
            \label{eq:clock_radius_transport}
        \end{align}

        \textbf{[CONDITIONAL DERIVED]}: within the horn-anchored approximation, contraction of a material loop drives stronger clock suppression, while expansion relaxes it.

    \subsection{Reparametrized Zero-Parameter Corollary}
        \label{subsec:zero_parameter_corollary}

        Older SST canon layers often used the triplet $(\Gamma_0,\rhoF,\rc)$ as the primitive parameterization. In the present v0.8.10 architecture, Eq.~\eqref{eq:rc_definition} and Eq.~\eqref{eq:gamma0} show that this is a reparameterization of the present primitive calibration set $(\rhoF,\vchar,\omega_c)$ rather than a distinct axiom.

        Explicitly,

        \begin{align}
            \Gamma_0 = 2\pi \rc \, \vchar
            = 2\pi \frac{\vchar{}^2}{\omega_c}.
            \label{eq:gamma0_reparam}
        \end{align}

        Hence any dimensional SST quantity written in terms of $(\Gamma_0,\rhoF,\rc)$ can be rewritten in terms of $(\rhoF,\vchar,\omega_c)$ with no new calibration input.

        A coarse-grained density law retained from earlier canon layers is

        \begin{align}
            \rhoF = K\,\Omega,
            \qquad
            K = \frac{\rhocore \, \rc}{\vchar},
            \label{eq:coarse_grain_density}
        \end{align}

        where $\Omega$ is a coarse-grained local rotation rate and $K$ has dimensions $\mathrm{kg\,m^{-3}\,s}$. This relation records the canonical bridge between a sparse effective bulk density and a high-density horn-envelope normalization; it is not a resolved local-tube density law unless $a_{\rm core}$ is supplied.

        \textbf{[CRITICAL NOTE]} Eq.~\eqref{eq:coarse_grain_density} is not an additional zero-parameter theorem unless the averaging prescription that defines $\Omega$ is fixed independently of the data being fitted. If $\Omega$ is chosen phenomenologically, the relation is a bridge ansatz and $K$ is a calibration transfer coefficient, not a new prediction.

        \textbf{[DERIVED COROLLARY]}: the older ``zero-parameter'' language is retained only as a statement of reparameterization discipline: once one primitive triplet is fixed, no additional dimensional constants should be introduced without explicit justification.

    \subsection{Canonical Master Equation}
    \label{sec:master-equation}

        The preceding subsections define the four ingredients that later SST papers often treat separately:
        a local mass-equivalent density, a local Swirl-Clock, a geometric gate, and a topological kernel.
        We therefore introduce a single canonical equation that organizes these sectors in one place.

        First define the local swirl energy density and its mass-equivalent form:

        \begin{align}
            \rhoE(x,t) &= \frac{1}{2}\rhoF(x,t)\,\norm{\uswirl(x,t)}^2 \\
            \rhoM(x,t) &= \frac{\rhoE(x,t)}{c^2}
        \end{align}

        Using Eq.~\ref{eq:swirl_clock}, the local Swirl-Clock factor is

        \begin{align}
            \SwirlClock(x,t) = \sqrt{1 - \frac{\norm{\uswirl(x,t)}^2}{c^2}}
        \end{align}

        We then define a binary geometric exposure gate
        \begin{align}
            G(K) \in \{0,1\}
        \end{align}
        together with the canonical geometric factor
        \begin{align}
            \Pi_K = \left( \frac{\lambda_c}{\pi \rc} \right)^{G(K)} = \left( \frac{4}{\alpha} \right)^{G(K)} .
        \end{align}

        \textbf{[CANON / ALPHA-GATE GUARD]} The equality \(\lambda_c/(\pi\rc)=4/\alpha\) is a calibrated geometric rewriting inside the present constant chain. It must not be cited as an independent derivation of the fine-structure constant. The finite-cell trefoil result is retained separately as an \(\alpha\)-free sub-per-mille coincidence plus obstruction result in Subsection~\ref{subsec:finite_cell_alpha_obstruction}.

        Numerically,
        \begin{align}
            \frac{\lambda_c}{\pi\rc}
            =
            \frac{4}{\alpha}
            \simeq 5.48144\times10^{2}.
            \label{eq:alpha_gate_numeric_guard}
        \end{align}
        This number is therefore an exposure or shielding gate inside the calibrated mass functional, not a new measurement or theorem for \(\alpha\).


        The canonical topological kernel is taken to be
        \begin{align}
            \Xi_K = \left[ \alpha_C\,C(K) + \beta_L\,L(K) + \gamma_H\,\mathcal{H}(K) \right] \varphi^{-2k(K)} .
        \end{align}

        \textbf{[CANON / GOLDEN-LAYER GUARD].}
        The golden-layer factor has an exact hyperbolic reparametrization,
        \begin{align}
            \eta_\varphi
            &\equiv
            \operatorname{asinh}\!\left(\frac12\right)
            =\ln\varphi,
            \label{eq:golden_layer_rapidity_def}\\
            \varphi^{-2k(K)}
            &=
            \exp\!\left[-2k(K)\eta_\varphi\right].
            \label{eq:golden_layer_rapidity_guard}
        \end{align}
        This is a parametrization guard for the existing layer factor, not a
        derivation of the golden factor from Euler dynamics, ropelength
        criticality, or knot stability.  It also does not assert that the
        electron trefoil is a hyperbolic knot; the trefoil remains a torus knot
        in the topology dictionary.  Any future promotion of this notation must
        identify \(k(K)\) with a mechanically defined discrete transfer,
        impedance, or relaxation step.

        The canonical SST master equation is therefore
        \begin{align}
            M_K(x,t)
            &= \rhoM(x,t)\,\rc^3\,\Pi_K\,\Xi_K\,\SwirlClock(x,t)^{-2} \\
            &= \left( \frac{\rhoF(x,t)\,\norm{\uswirl(x,t)}^2}{2c^2} \right)
               \rc^3
               \left( \frac{\lambda_c}{\pi \rc} \right)^{G(K)}
               \left[ \alpha_C\,C(K) + \beta_L\,L(K) + \gamma_H\,\mathcal{H}(K) \right]
               \varphi^{-2k(K)}
               \SwirlClock(x,t)^{-2}
            \label{eq:sst_master_equation}
        \end{align}

        \textbf{[SPECULATIVE] Organizing principle}: This equation factorizes SST into four sectors:
        \begin{align}
            \text{medium scale} \times \text{geometric gate} \times \text{topological kernel} \times \text{clock impedance} .
        \end{align}

        \textbf{[DERIVED] Dimensional check}:
        \begin{align}
            \left[ \frac{\rhoF\,\norm{\uswirl}^2}{c^2} \rc^3 \right] = \mathrm{kg}
        \end{align}
        and all remaining factors are dimensionless, so Eq.~\ref{eq:sst_master_equation} has the correct mass dimension.

        \textbf{[CANON / SEMANTIC RULE]}: the factor \(\rc^3\) in Eq.~\eqref{eq:sst_master_equation} is a canonical horn/envelope normalization volume scale. It is not a resolved physical tube volume unless a separate model sets \(a_{\rm core}=\rc\). Resolved tube energies, cutoffs, and local vorticity estimates must use \(a_{\rm core}\) or an explicit profile.

        \textbf{[SPECULATIVE] Interpretation}: The same local kinematic state that slows the Swirl-Clock also enters the effective mass scale; mass and local proper-time modulation are therefore treated as two aspects of one medium configuration.

        \textbf{[CRITICAL NOTE]} Equation~\eqref{eq:sst_master_equation} is presently a canonical synthesis formula, not yet a theorem derived from first-principles filament dynamics alone.

        \textbf{[NORMALIZATION TRANSPARENCY]} With the canonical background-energy reading, $\rhoM=\rhoE/c^2=4.65949\times10^{-12}\,\mathrm{kg\,m^{-3}}$ and $\rhoM\rc^3=1.30330\times10^{-56}\,\mathrm{kg}=1.43072\times10^{-26}m_e$. Therefore an exposed electron-scale benchmark with $\Pi_K=4/\alpha$ requires a dimensionless kernel of order
        \begin{align}
            \Xi_{3_1}^{(e)}
            \sim
            \frac{m_e}{\rhoM\rc^3(4/\alpha)}
            =
            1.28\times10^{23}.
            \label{eq:kernel_normalization_scale_guard}
        \end{align}
        Unless a branch explicitly re-anchors the medium scale to $\rhohorn$, the coefficients inside $\Xi_K$ are therefore calibrated matching coefficients carrying the medium-to-particle hierarchy, not order-unity topological invariants. This note is a scale-accounting guard; it does not alter the dimensional validity of Eq.~\eqref{eq:sst_master_equation}.

        In the stationary weak-swirl limit $\SwirlClock \to 1$, it reduces to the static mass-functional form; in the local-kinematic limit at fixed $K$, it reduces to a pure clock-impedance correction. The hierarchy is therefore carried predominantly by the geometric and topological factors, while the Swirl-Clock supplies a coherent local correction rather than the primary source of mass separation.

    \subsection{Geometric Representation}
        
        A swirl-string is represented as a closed curve:
        
        \begin{align}
            X(t,\sigma) = (x(t,\sigma), y(t,\sigma), z(t,\sigma))
        \end{align}
        
        with $\sigma \in [0,2\pi)$.
    
    \subsection{Summary of Dependency Structure}
        
        The full dependency chain is:
        
        \begin{align}
        (\rhoF, \vchar, \omega_c)
            \rightarrow \rc
            \rightarrow \Gamma_0
            \rightarrow (\rhoE, \rhoM, \rhocore)
            \rightarrow \SwirlClock
            \rightarrow M_K.
        \end{align}
        
        \textbf{[DERIVED] Consistency statement}:
        No quantity is defined circularly. The canonical closure is summarized by Eq.~\ref{eq:sst_master_equation}.
    
    \subsection{Consistency Statement}
        
        \begin{itemize}
            \item All derived quantities depend on the stated calibration set, explicit orthodox constants, and explicitly labeled bridge assumptions.
            \item No \emph{unlabeled} external parameters are introduced.
            \item Quantities fixed by $\alpha$, $m_e$, $\hbar$, $c$, $R_\infty$, or $t_p$ are labeled \textbf{[CALIBRATED]} rather than independent predictions.
        \end{itemize}
    
    \subsection{Interpretation}
        
        The Formal Foundations establish SST as:
        
        \begin{itemize}
            \item a continuum theory with minimal primitives
            \item a geometrically grounded framework
            \item a non-circular parameter system
        \end{itemize}

    \subsection{Maximal Swirl Tension}
    \label{sec:Fmax}

        \subsubsection{Definition and canonical form}

        Within the primitive calibration set $\mathcal{P}_{\mathrm{cal}} = \{\rhoF, \vchar, \omega_c\}$, a characteristic force scale emerges naturally from the Compton energy stored over the canonical circulation cross-section. We define the \emph{maximal swirl tension} as

        \begin{align}
            F_{\rm swirl}^{\max}
            \;=\; \frac{\vchar\,\hbar}{2\rc^2}
            \;=\; \frac{\hbar\omega_c}{2\rc}.
            \label{eq:Fmax_def}
        \end{align}

        \textbf{[CALIBRATED ALGEBRAIC IDENTITY] Canonical form}: substituting $\vchar = \omega_c \rc$ into Eq.~\eqref{eq:Fmax_def} gives $F_{\rm swirl}^{\max} = \hbar\omega_c/2\rc$, the Compton energy scale divided by twice the canonical circulation radius. Numerically, $F_{\rm swirl}^{\max} = 29.053507\;\text{N}$.

        \textbf{[CRITICAL NOTE]} $F_{\rm swirl}^{\max}$ is a mechanical reformulation of the Compton energy scale over the canonical circulation cross-section. It is \emph{not} independent of $\hbar$ and $m_e$; its role in the canon is interpretive and organisational rather than that of an independent primitive.

        \subsubsection{Internal consistency identities}

        The following expressions are algebraically equivalent within the canonical constant chain and serve as coherence verifications rather than independent derivations:

        \begin{align}
            F_{\rm swirl}^{\max}
            &= \frac{e^2}{16\pi\varepsilon_0 \rc^2},
            \label{eq:Fmax_em}\\
            &= \frac{h\alpha c}{8\pi \rc^2},
            \label{eq:Fmax_hac}\\
            &= \pi \rc^2\,\rho_{\rm calc}\,\vchar{}^2,
            \label{eq:Fmax_rho}\\
            &= \frac{c^4}{4G}\,\alpha
            \left(\frac{\rc}{\ell_P}\right)^{-2}.
            \label{eq:Fmax_planck_suppressed}
        \end{align}

        where $\rho_{\rm calc} = 4F_{\rm swirl}^{\max}/(\pi\alpha^2 c^2 \rc^2)$ is the derived bulk density consistent with the canonical neutral-circulation closure.

        \textbf{[CALIBRATED ALGEBRAIC IDENTITY] Consistency check}: Eqs.~\eqref{eq:Fmax_em}--\eqref{eq:Fmax_planck_suppressed} are equivalent given $\vchar = \alpha c/2$, the standard identity $e^2/(4\pi\varepsilon_0) = \alpha\hbar c$, $\rc=\alpha\hbar/(2m_ec)$, and $\ell_P^2=\hbar G/c^3$. The Planck-force form in Eq.~\eqref{eq:Fmax_planck_suppressed} is a hierarchy reparameterization of the same calibrated tension scale, not an independent derivation from Planck physics. Their mutual agreement is an internal algebraic coherence test of the calibrated chain, not an independent dynamical prediction.

        \subsubsection{Rydberg consistency identity}

        \textbf{[CALIBRATED IDENTITY]} The Rydberg constant is recoverable from the three SST calibration variables without $e$, $\varepsilon_0$, or $m_e$ appearing \emph{explicitly}:

        \begin{align}
            R_\infty \;=\; \frac{\vchar{}^3}{\pi\,\rc\,c^3}.
            \label{eq:Rydberg_SST}
        \end{align}

        \textbf{[DIMENSIONAL GUARD]} The direction of Eq.~\eqref{eq:Rydberg_SST} is fixed: the numerator is $\vchar{}^3$ and the denominator is $\pi\rc c^3$. Its units are $\mathrm{m^{-1}}$. The reciprocal expression $\pi\rc c^3/\vchar{}^3$ has units of length and is not a Rydberg constant.

        This is a benchmark for algebraic consistency of the calibration chain, not for independent non-redundancy of the primitive set. \textbf{[CRITICAL NOTE]} Substituting the canonical definitions $\vchar=\alpha c/2$ and $\rc=\vchar/\omega_c=\alpha\hbar/(2m_ec)$ reduces Eq.~\eqref{eq:Rydberg_SST} \emph{exactly} to the standard $R_\infty=\alpha^2 m_e c/(4\pi\hbar)$ (verified to machine precision): $m_e$ is hidden inside $\rc$ and $\alpha$ inside both $\vchar$ and $\rc$. The identity is therefore a consistency check, not an independent non-circular derivation of $R_\infty$. Combining Eq.~\eqref{eq:Rydberg_SST} with Eq.~\eqref{eq:Fmax_def} yields the compact relation

        \begin{align}
            \boxed{
            F_{\rm swirl}^{\max}
            \;=\; \frac{32\pi^2\hbar R_\infty^2 c}{\alpha^5},
            }
            \label{eq:Fmax_Rydberg}
        \end{align}

        which expresses $F_{\rm swirl}^{\max}$ in terms of spectroscopic observables alone, with no reference to $m_e$, $\rc$, or $e$ as separate inputs. A further identity connecting atomic and force scales is

        \begin{align}
            R_\infty \;=\; F_{\rm swirl}^{\max} \cdot \frac{\alpha^3}{4\pi m_e c^2},
            \label{eq:Rydberg_Fmax}
        \end{align}

        which states that the Rydberg constant is the maximal swirl tension scaled by the Compton energy and a purely geometric factor $\alpha^3/4\pi$.

        \subsubsection{Threefold physical interpretation}

        The quantity $F_{\rm swirl}^{\max}$ admits three distinct physical interpretations, each arising from a different sector of physics. \textbf{[CRITICAL NOTE]} These three values are algebraically equivalent expressions built from the same primitive set $\{\alpha,m_e,\hbar,c\}$; their numerical agreement is an \emph{identity}, not an independent coincidence, and is not on its own confirmation of the theory.

        \paragraph{(i) Gravitational analogue: the Planck force ceiling.}

        In general relativity, the combination $c^4/4G$ sets the maximum force that any physical process can transmit without forming a causal horizon~\cite{Gibbons2002}. We denote this scale the maximal gravitational tension:

        \begin{align}
            F_{\rm gr}^{\max} \;=\; \frac{c^4}{4G}.
            \label{eq:Fgr}
        \end{align}

        \textbf{[ORTHODOX]} $F_{\rm gr}^{\max}$ is the Planck force, an invariant upper bound in classical general relativity \cite{Gibbons2002}.

        \textbf{[CALIBRATED ALGEBRAIC IDENTITY]} The ratio of the two tension scales reduces exactly to

        \begin{align}
            \boxed{
            \frac{F_{\rm swirl}^{\max}}{F_{\rm gr}^{\max}}
            \;=\; \frac{4\alpha_g}{\alpha},
            }
            \label{eq:Fmax_ratio}
        \end{align}

        where $\alpha_g = Gm_e^2/\hbar c$ is the gravitational fine-structure constant of the electron and $\alpha$ is the electromagnetic fine-structure constant. Numerically, $4\alpha_g/\alpha \approx 9.6\times10^{-43}$, in agreement with the directly computed ratio $29.053507\;\text{N}\,/\,3.026\times10^{43}\;\text{N}$.

        \textbf{[SPECULATIVE] Interpretation}: the hierarchy between the electromagnetic and gravitational force scales may be organized by the ratio of the two electron coupling constants. This is an interpretive compression of calibrated constants; by itself it does not prove a new medium-dynamical stiffness law.

        \paragraph{(ii) Electrostatic interpretation: the Coulomb ceiling at the canonical circulation scale.}

        When two protons are compressed to a centre-to-centre separation equal to the canonical circulation radius $\rc$, the Coulomb repulsion is

        \begin{align}
            F_{pp}(\rc)
            \;=\; \frac{e^2}{4\pi\varepsilon_0 \rc^2}
            \;=\; 4\,F_{\rm swirl}^{\max}.
            \label{eq:Coulomb_core}
        \end{align}

        \textbf{[CALIBRATED]} The maximal Coulomb repulsion at the canonical circulation scale is exactly four times $F_{\rm swirl}^{\max}$. \textbf{[CRITICAL NOTE]} The factor of four is an arithmetic identity once $\vchar=\alpha c/2$ is adopted (the $4/\alpha$ gate); it is a self-consistency of the calibrated chain, not an independent prediction.

        \textbf{[SPECULATIVE] Interpretation}: $F_{\rm swirl}^{\max}$ is the quarter-Coulomb force at the canonical circulation scale. The medium sustains one quarter of the maximum proton--proton Coulomb repulsion before the neutral circulation envelope reaches its structural limit. This provides a physical ceiling on the compression of charged vortex configurations.

        \paragraph{(iii) Mechanical interpretation: the spring constant of photon--electron interaction.}

        The electron's internal vortex is modelled as a linear spring with stiffness

        \begin{align}
            K_e \;=\; \frac{F_{\rm swirl}^{\max}}{n\,\rc},
            \qquad n = 2,
            \label{eq:spring_K}
        \end{align}

        where $n$ counts the number of half-wavelength segments on the core and is fixed to $n=2$ by the photon--electron swirl matching condition. The natural oscillation frequency of the electron core under this spring is

        \begin{align}
            \omega_{\rm spring}
            \;=\; \sqrt{\frac{K_e}{m_e}}
            \;=\; \sqrt{\frac{F_{\rm swirl}^{\max}}{2\,\rc\,m_e}}
            \;=\; \frac{\omega_c}{\alpha}
            \;=\; \frac{m_e c^2}{\alpha\hbar},
            \label{eq:omega_spring}
        \end{align}

        a frequency intermediate between the atomic Rydberg scale $(\sim\alpha^2\omega_c)$ and the bare Compton scale $(\omega_c)$.

        \textbf{[DERIVED]} The energy stored in this spring at maximum compression $\delta = \rc$ is

        \begin{align}
            E_{\rm spring}
            \;=\; \tfrac{1}{2}\,K_e\,\rc^2
            \;=\; \frac{F_{\rm swirl}^{\max}\,\rc}{4}
            \;=\; \frac{\hbar\omega_c}{8}
            \;=\; \frac{m_e c^2}{8}.
            \label{eq:Espring}
        \end{align}

        One eighth of the electron rest energy is stored as internal vortex-spring potential at maximum compression.

        \textbf{[SPECULATIVE] Physical picture}: during photon absorption, the electron vortex core is compressed for a duration of order one Planck time $t_p$ into a transient mass-like configuration. The restoring force scale of this compression is $F_{\rm swirl}^{\max}$, and the energy scale of the transient is $m_e c^2/8$. The spring releases at the Compton period, generating the recoil that initiates the Coulomb interaction with the proton.

        \textbf{[CRITICAL NOTE]} Equation~\eqref{eq:omega_spring} assumes $n=2$ from empirical matching. A first-principles derivation of the photon wrap number from the delay dynamics of Section~\ref{sec:delay} is an open derivational step.

        \subsubsection{Unified summary}

        The three interpretations are collected in Table~\ref{tab:Fmax_triple}. Together they organize $F_{\rm swirl}^{\max}$ as a \emph{calibrated structural tension scale}. The table is useful for cross-sector bookkeeping, but the equalities are algebraic consequences of the same calibrated electron-scale constants unless a separate medium-dynamical derivation is supplied.

        \begin{table}[h]
        \centering
        \caption{Threefold physical interpretation of $F_{\rm swirl}^{\max}$ within the SST canonical framework.}
        \label{tab:Fmax_triple}
        \begin{tabular}{p{3.2cm}p{5.0cm}p{4.8cm}l}
        \hline
        \textbf{Context} & \textbf{Expression} & \textbf{Physical meaning} & \textbf{Status} \\
        \hline
        Gravitational ceiling &
          $F_{\rm gr}^{\max}/F_{\rm swirl}^{\max} = \alpha/(4\alpha_g)$ &
          Calibrated electron-coupling ratio $\alpha/4\alpha_g$ &
          {[Calibrated identity]} \\[6pt]
        Coulomb ceiling &
          $F_{pp}(\rc) = 4\,F_{\rm swirl}^{\max}$ &
          Quarter of max.\ Coulomb force at circulation scale &
          {[Calibrated identity]} \\[6pt]
        Vortex spring &
          $K_e = F_{\rm swirl}^{\max}/2\rc$ &
          Spring constant of photon--electron interaction &
          {[Speculative]} \\
        \hline
        \end{tabular}
        \end{table}

        \textbf{[SPECULATIVE] Canonical interpretation}: $F_{\rm swirl}^{\max}$ is retained as a candidate tension scale for topological integrity of the structured medium. Its comparison with $F_{\rm gr}^{\max}=c^4/4G$ is a calibrated electron-normalized analogy, not a theorem that gravitational and electromagnetic stiffness have already been derived from the substrate.

% ---------- AXIOMS ----------
\section{Axiomatic Framework}
    \label{sec:axioms}
    
    \subsection{Preliminaries}
        
        We formalize Swirl-String Theory (SST) as an axiomatic continuum theory defined over a three-dimensional spatial domain $\mathcal{M} \subset \mathbb{R}^3$ equipped with a preferred time parameter $t \in \mathbb{R}$.
        
        The theory is constructed from a minimal primitive set $\mathcal{P}$ (Section~\ref{sec:foundations}) and a set of axioms $\mathcal{A}$ governing admissible configurations and dynamics.
    
    \subsection{Axiom 1: Continuum Substrate}
        
        \textbf{Statement.}
        There exists a continuous, incompressible, inviscid medium filling $\mathcal{M}$, characterized by a constant background density $\rhoF$.
        
        \begin{align}
            \nabla \cdot \mathbf{v} = 0
        \end{align}
        
        \textbf{[ORTHODOX]}: This corresponds to the standard incompressible Euler framework.
        
        \textbf{Interpretation.}
        All physical structures arise as configurations within this medium; no discrete particles are assumed at the fundamental level.
    
    \subsection{Axiom 2: Existence of Stable Vortex Filaments}
        
        \textbf{Statement.}
        The medium admits stable, localized, topologically nontrivial vortex filament solutions.
        
        \textbf{[ORTHODOX]} Physics input: vortex filaments are known idealized structures in fluid dynamics.

        \textbf{[SPECULATIVE]} Extension: such filaments are assumed here to be dynamically stable and persist as fundamental carriers.
        
        \textbf{Mathematical form.}
        A filament is represented as a closed embedding:
        
        \begin{align}
            X: S^1 \to \mathcal{M}, \quad \sigma \mapsto X(t,\sigma)
        \end{align}
    
    \subsection{Axiom 3: Circulation Quantization}
        
        \textbf{Statement.}
        Circulation around any stable filament is quantized:
        
        \begin{align}
            \Gamma = \oint \mathbf{v} \cdot d\ell = n \Gamma_0, \quad n \in \mathbb{Z}
        \end{align}
        
        \textbf{[ORTHODOX]}: Quantized circulation appears in superfluid systems.
        
        \textbf{[SPECULATIVE]}: Elevated here to a universal invariant.
    
    \subsection{Axiom 4: Compton Anchoring}
        
        \textbf{Statement.}
        The Compton phase frequency of the canonical horn/circulation carrier satisfies:
        
        \begin{align}
            \frac{\vchar}{\rc} = \omega_c
        \end{align}
        
        \textbf{[ORTHODOX]}: $\omega_c$ is the Compton frequency.
        
        \textbf{[SPECULATIVE]}: Identified as the intrinsic Compton phase frequency of the SST carrier, not as the local vorticity of the resolved tube.
        
        \textbf{Consequence.}
        The horn/circulation radius is not independent but derived (Eq.~\ref{eq:rc_definition}); the local tube radius remains $a_{\rm core}$.
    
    \subsection{Axiom 5: Delay as a Constitutive Principle}
        
        \textbf{Statement.}
        Finite propagation time along closed vortex structures is a fundamental dynamical ingredient.
        
        \begin{align}
            \tau_{\rm circ}(K) = \frac{L_K}{\vchar}
        \end{align}
        
        \textbf{[DERIVED]} Kinematic consequence: the delay follows from finite signal speed along a closed carrier.

        \textbf{Axiomatic role.} The delay is not treated as negligible, but as structurally relevant to the dynamics.
        
        \textbf{Interpretation.}
        Temporal nonlocality is intrinsic to the dynamics.
    
    \subsection{Axiom 6: Energy Localization}
        
        \textbf{Statement.}
        Energy is localized in vortex structures and determines mass via:
        
        \begin{align}
            E = m c^2
        \end{align}
        
        \textbf{[ORTHODOX]}: Relativistic energy relation.
        
        \textbf{[SPECULATIVE]}: Interpreted as rotational energy of the vortex configuration.
    
    \subsection{Axiom 7: Swirl-Clock Relational Time}
        
        \textbf{Statement.}
        Local time is encoded in a scalar field derived from fluid kinematics:
        
        \begin{align}
            \SwirlClock = \sqrt{1 - \frac{\vchar{}^2}{c^2}}
        \end{align}
        
        \textbf{[ORTHODOX]} Functional form: this matches Lorentz time dilation.

        \textbf{[SPECULATIVE]} Identification: the relevant velocity is taken to be the local medium kinematic state.

        \textbf{[DERIVED]} Consequence: within SST this yields a relational clock field tied to local motion.
        
        \textbf{Interpretation.}
        Time is not fundamental but relational to the state of motion in the medium.
    
    \subsection{Axiom 8: Density Saturation and Horn-Envelope Structure}
        
        \textbf{Statement.}
        The canonical horn-envelope exhibits an effective saturated density
        \(\rhohorn\) determined by the mechanical normalization constraints of the
        medium. A resolved local tube density is not specified until
        \(a_{\rm core}\) and a tube profile are supplied.
        
        \textbf{[SPECULATIVE]} Closure hypothesis: the effective saturated
        envelope density is fixed by the closure relation of Eq.~\eqref{eq:core_density}.

        \textbf{[DERIVED]} Consequence: once that closure is adopted, matter
        corresponds to configurations whose neutral circulation envelope reaches
        a structural normalization limit.
    
    \subsection{Logical Structure of the Axioms}
        
        The axioms form a dependency hierarchy:
        
        \begin{align}
            \text{Continuum}
            \rightarrow \text{Vortex Structures}
            \rightarrow \text{Quantization}
            \rightarrow \text{Delay Dynamics}
            \rightarrow \text{Mass/Energy Emergence}
        \end{align}
    
    \subsection{Minimality Principle}
        
        \textbf{Statement.}
        The axioms are minimal in the sense that removing any one of them destroys at least one of the following:
        
        \begin{itemize}
            \item existence of stable structures
            \item emergence of discrete spectra
            \item compatibility with known physics
        \end{itemize}
    
    \subsection{Philosophical Interpretation}
        
        The axiomatic structure implies:
        
        \begin{itemize}
            \item Particles are not fundamental objects but stable dynamical configurations
            \item Discreteness is emergent, not postulated
            \item Time is relational, not absolute
        \end{itemize}
    
    \subsection{Consistency Condition}
        
        A valid SST realization must satisfy:
        
        \begin{itemize}
            \item Compatibility with orthodox physics in appropriate limits
            \item Internal logical closure
            \item Explicit traceability from axioms to predictions
        \end{itemize}
    
    \subsection{Canonical Statement}
        
        \begin{center}
            \textit{
                Swirl-String Theory is defined as the minimal continuum theory in which stable vortex structures, governed by quantized circulation and delay dynamics, give rise to discrete physical phenomena consistent with known physics.
            }
        \end{center}

% ---------- CONSISTENCY LAYER ----------
%    \section{Consistency and Classification Framework}
%        \label{sec:consistency}
        
\section{Consistency and Classification Framework}
    \label{sec:consistency}
    
    \subsection{Classification Scheme}
        
        All statements are labeled as:
        
        \begin{itemize}
            \item \textbf{[ORTHODOX]} — established physics
            \item \textbf{[DERIVED]} — logically deduced within SST
            \item \textbf{[SPECULATIVE]} — novel or unverified
        \end{itemize}

    \subsection{Extended Labeling Convention}

        We distinguish between \emph{epistemic status} and \emph{editorial role}.

        \paragraph{Epistemic status.}
        The primary status labels used in the present Canon are:
        \begin{itemize}
            \item \textbf{[ORTHODOX]} --- established in standard physics or mathematics
            \item \textbf{[DERIVED]} --- obtained internally from the present SST assumptions, definitions, or closures, \emph{without re-using the target quantity as an input}
            \item \textbf{[CALIBRATED]} --- fixed by, or algebraically equivalent to, a measured constant ($\alpha$, $m_e$, $\hbar$, $R_\infty$, \dots) that enters through the primitive set; numerically exact but not an independent prediction
            \item \textbf{[CONDITIONAL]} --- follows only if a stated, not-yet-proven premise holds (e.g.\ an open lemma or selection rule)
            \item \textbf{[SPECULATIVE]} --- conjectural SST extension, interpretation, or unverified identification
        \end{itemize}
        A \textbf{[CRITICAL NOTE]} editorial overlay flags a circularity risk, dimensional issue, or reviewer vulnerability on any of the above.

        \paragraph{Editorial role.}
        Additional descriptors such as \emph{Definition}, \emph{Constraint}, \emph{Interpretation}, \emph{Benchmark}, \emph{Roadmap}, and \emph{Dimensional check} are not independent epistemic classes. They indicate only the role played by the statement.

        \paragraph{Usage rule.}
        When both are needed, the Canon uses the format
        \begin{align}
            \texttt{[STATUS] Role: statement.}
        \end{align}
        For example, one may write \textbf{[ORTHODOX] Constraint}, \textbf{[DERIVED] Consequence}, or \textbf{[SPECULATIVE] Interpretation}.
    
    \subsection{Orthodox Constraints}
        
        The following must remain intact:
        
        \begin{itemize}
            \item Quantum mechanics predictions
            \item Atomic spectroscopy
            \item Relativistic consistency
        \end{itemize}
    
    \subsection{Internal Consistency Rules}
        
        \begin{itemize}
            \item No circular definitions
            \item All primitives explicitly defined
            \item Derived quantities must follow from axioms
        \end{itemize}
    
    \subsection{Conflict Resolution Strategy}
        
        Priority order:
        
        \begin{enumerate}
            \item Logical consistency
            \item Experimental agreement
            \item Canonical simplicity
        \end{enumerate}
    
    \subsection{Epistemic Status Tracking}
        
        Each major result must explicitly state:
        
        \begin{itemize}
            \item Assumptions
            \item Derivation status
            \item Empirical testability
        \end{itemize}

        In particular, a definitional choice must not be presented as a theorem, and an interpretive bridge must not be presented as an orthodox result unless the underlying mathematical structure is itself standard.
    
    \subsection{Canon Integrity Principle}
        
        \begin{center}
            \textit{A valid SST canon must be internally consistent, externally compatible, and explicitly stratified by epistemic status.}
        \end{center}

    \section{Geometric and Topological Sector}
        \label{sec:geometry}

    \subsection{Role of geometry and topology in the canon}
        This sector collects those parts of SST in which the identity of a state depends not only on local field amplitudes but on the global organization of the carrier. Geometry enters through lengths, radii, and closure constraints; topology enters through linking, crossing, chirality, and other invariant features that survive smooth deformations. At canon level, the main role of this sector is selective rather than exhaustive: it provides the organizing variables that feed the topological kernel $\Xi_K$ and the geometric gate $\Pi_K$ in Eq.~\eqref{eq:sst_master_equation}.

        The canon therefore treats geometry and topology as follows:
        \begin{enumerate}
            \item geometry sets admissible scales and shielding ratios,
            \item topology labels persistent identity classes,
            \item only those quantities that survive smooth deformation are allowed to enter the long-lived state classification.
        \end{enumerate}

        This is also why layer or dressing indices are not treated as primary identity labels: they may modify a state, but they do not replace the topology-first classification of the carrier itself.

    \subsection{Euler topological preservation and closure discipline}
        \label{subsec:euler_topological_preservation_closure}

        \textbf{[ORTHODOX].}
        In a smooth incompressible inviscid Euler flow, with no boundaries, forcing,
        viscosity, singular core collapse, or phase-slip defects, vortex lines are
        materially transported by the flow. Consequently, linking and knot type of
        already closed vortex carriers are preserved. A closed unknot cannot become
        a trefoil by ordinary ideal-Euler evolution:
        \begin{align}
            0_1 \not\longrightarrow 3_1
            \qquad
            \text{under smooth ideal Euler flow.}
            \label{eq:euler_no_unknot_to_trefoil}
        \end{align}

        \textbf{[DERIVED SST CONSTRAINT].}
        Trefoil creation must therefore not be described as reconnection of an
        already closed ideal-Euler vortex tube. It is admissible only as a
        boundary-supported closure, phase-slip, nucleation, or other non-Euler
        event occurring before the final carrier has acquired a protected closed
        knot type. After closure, smooth evolution may deform the embedding but
        must preserve the knot class.

        \textbf{[CANON RULE].}
        Topology protects existing closed carriers; it does not by itself provide
        an ideal-Euler mechanism for creating them.

    \subsection{Framed self-linking of the swirl string and the spinorial selection of the lepton ladder}
        \label{subsec:framed_selflinking_spinorial}

        This subsection records the framed-tube self-linking of a closed torus-knot swirl
        string and the boundary class that selects its physical winding. It is the first
        genuinely non-circular topological result in the canon, in deliberate contrast to the
        calibrated $\alpha$-identities of the Formal Foundations (which re-express $\alpha$
        rather than predict it).

        \textbf{[DERIVED] Torus-surface self-linking.}
        Let $K=T(p,q)$ be a closed swirl string on the standard torus and let $U_{\rm tor}$ be
        the torus surface-normal framing. The framed self-linking is
        \begin{align}
            SL_{\rm tor}(T(p,q)) \;=\; pq,
            \label{eq:SL_torus_pq}
        \end{align}
        verified numerically to machine precision for $p\in\{2,3\}$ and $q=3,\dots,11$ with
        \emph{no} target injected. For the lepton family $T(2,q)$, $SL_{\rm tor}=2q$.

        \textbf{[DEFINITIONAL] Physical-framing winding.}
        For any single-valued physical framing $U_{\rm phys}$ of $K$ (the direction of the local core
        swirl-velocity field $\uswirl$), the relative winding
        \begin{align}
            \chi \;:=\; SL(K,U_{\rm phys})-SL(K,U_{\rm tor})
                  \;=\; \frac{1}{2\pi}\oint_K \mathrm{d}\!\left[\angle_{U_{\rm tor}}(U_{\rm phys})\right]
                  \;\in\; \mathbb{Z}
            \label{eq:chi_def}
        \end{align}
        is an integer relative winding number, so $SL_{\rm phys}(T(p,q))=pq+\chi$.

        \textbf{[CRITICAL NOTE]} Any test that fixes the twist by $Tw=\mathrm{target}-Wr$ and
        then confirms the target is circular: it passes for \emph{any} integer target and
        proves only the Calugareanu identity $SL=Wr+Tw$. Only the target-free
        framing~\eqref{eq:SL_torus_pq} and the selection argument below are admissible.

        \textbf{[DERIVED NEGATIVE] Ordinary $U(1)$ closure gives $\chi=0$.}
        With positive longitudinal phase stiffness $A_\parallel>0$, the winding energy
        \begin{align}
            E(\chi)\;=\;E_0+\frac{2\pi^2 A_\parallel}{L}\,\chi^2
            \label{eq:E_chi}
        \end{align}
        is minimized over ordinary (periodic) $U(1)$ windings $\chi\in\mathbb{Z}$ at
        $\chi_{\rm ground}=0$, hence $SL_{\rm phys}=pq$ (even). Pure framed-curve geometry with
        positive stiffness therefore yields a \emph{boson} and cannot produce the lepton
        $pq+1$ ladder.

        \textbf{[CONDITIONAL DERIVED] Spinorial anti-periodicity gives $\chi=\pm1$.}
        If the core phase is a section of a spinor (rather than a $U(1)$) bundle, then under one
        circuit of $K$ the spin frame rotates by $2\pi$ and the phase is anti-periodic, so the
        admissible windings are $\chi\in 2\mathbb{Z}+1$. Minimizing~\eqref{eq:E_chi} over this
        sector gives the degenerate ground winding
        \begin{align}
            \chi_{\rm ground}=\pm1 \quad\Longrightarrow\quad
            SL_{\rm phys}(T(p,q))=pq\pm1,
            \qquad SL_{\rm phys}(T(2,q))=2q\pm1,
            \label{eq:SL_phys_pq1}
        \end{align}
        the two signs being the matter and antimatter sectors. The selection is a
        \emph{superselection} effect: the boson winding $\chi=0$ is always energetically lower
        (by $2\pi^2 A_\parallel/L>0$), so it is the boundary class, not the energetics, that
        renders the carrier fermionic. The result is independent of $(A_\parallel,L)$ for all
        positive values.

        \paragraph{Why the scalar substrate is bosonic and $\theta=\pi$ is the fermionic $\mathbb{Z}_2$ choice.}
        The two boundary classes are not arbitrary continuous tunings.  In an
        $S^2\simeq\mathbb{C}P^1$ order-parameter theory they can be identified with the two values
        of a discrete topological/statistical sector.

        \textbf{[DERIVED NEGATIVE] Scalar $U(1)$ substrate.}
        A single-component order parameter $\psi=|\psi|e^{i\varphi}$ is single-valued only if
        \begin{equation}
            \oint \mathrm{d}\varphi \in 2\pi\mathbb{Z},
        \end{equation}
        so its scalar phase is periodic:
        \begin{equation}
            e^{i\oint \mathrm{d}\varphi}=+1
        \end{equation}
        on every closed loop.  This condition allows integer windings, but it does not force
        the longitudinal framing class to be odd.  With the positive stiffness
        energy~\eqref{eq:E_chi},
        \begin{equation}
            E_\chi = E_0+\frac{2\pi^2 A_\parallel}{L}\,\chi^2,\qquad A_\parallel>0,
        \end{equation}
        minimization over the ordinary scalar class $\chi\in\mathbb{Z}$ selects
        \begin{equation}
            \chi_{\rm ground}=0 .
        \end{equation}
        Hence a purely scalar SST substrate has no intrinsic mechanism for the odd core-phase
        sector required by fermions.  In this operational sense the scalar substrate is bosonic:
        it admits the trivial periodic carrier and does not generate the $pq\pm1$ lepton ladder.

        \textbf{[DERIVED / ALLOWED] $S^2\simeq\mathbb{C}P^1$ admits a $\mathbb{Z}_2$ statistical sector.}
        The matter target is not a scalar $U(1)$ phase but the director/Hopfion target
        \begin{equation}
            M_{\rm SST}=S^2\simeq\mathbb{C}P^1 .
        \end{equation}
        The relevant homotopy data are
        \begin{equation}
            \pi_3(S^2)=\mathbb{Z},\qquad
            \pi_4(S^2)=\mathbb{Z}_2 .
        \end{equation}
        The first class gives the static Hopf charge,
        \begin{equation}
            Q_H = SL_{\rm phys},
        \end{equation}
        while the second allows a $\mathbb{Z}_2$ topological/statistical phase for spacetime
        loops of soliton configurations~\cite{WilczekZee1983}.  More carefully, the fermionic
        sign is not simply a local scalar factor $e^{i\theta Q_H}$ in the static Hopf charge.
        It is the sign of the relevant $\mathbb{Z}_2$ spacetime class,
        \begin{equation}
            \exp(iS_\theta)=(-1)^\nu,\qquad
            \nu\in\mathbb{Z}_2 ,
        \end{equation}
        whose representative may be read, for the relevant $2\pi$ rotation or exchange process,
        as the parity class of the Hopf sector.  Thus the two possible statistical sectors are
        \begin{align}
            \theta=0 &: \quad \exp(iS_\theta)=+1
                \quad\Rightarrow\quad \text{bosonic / periodic sector},\\
            \theta=\pi &: \quad \exp(iS_\theta)=(-1)^\nu
                \quad\Rightarrow\quad \text{fermionic / spinorial sector for odd Hopf parity}.
        \end{align}
        Identifying the spinorial sector with the longitudinal core-phase boundary class gives
        \begin{equation}
            \theta=\pi
            \quad\Longleftrightarrow\quad
            \chi\in2\mathbb{Z}+1 .
        \end{equation}
        Then positive stiffness no longer minimizes over all integers but over the odd class:
        \begin{equation}
            \chi\in2\mathbb{Z}+1
            \quad\Rightarrow\quad
            \chi_{\rm ground}=\pm1 .
        \end{equation}
        Since
        \begin{equation}
            Q_H=SL_{\rm phys}=SL_{\rm tor}+\chi=pq+\chi,
        \end{equation}
        and $pq=2q$ is even for the lepton family $T(2,q)$, the $\theta=\pi$ sector gives
        \begin{equation}
            SL_{\rm phys}(T(2,q))=2q\pm1,
        \end{equation}
        i.e. the odd Hopf/self-linking carrier required for fermionic leptons.

        \textbf{[POSIT, pinned] $\theta=\pi$ is selected by observation, not derived.}
        The framed-curve geometry and the quadratic stiffness energy do not prefer
        $\theta=\pi$ over $\theta=0$.  The value $\theta=\pi$ is therefore a discrete
        superselection input pinned by the empirical existence of fermions.  It is naturally
        identified with the same $\mathbb{Z}_2$ datum appearing in half-quantum vortices of
        multi-component condensates~\cite{Volovik2003,Autti2016} and in the half/full
        circulation assignment of the magnetic-moment sector (Route~B).  With that
        identification, fermion statistics, the $pq\pm1$ ladder, and the Route~B origin of
        $g=2$ are no longer independent posits but three faces of one $\mathbb{Z}_2$ sector.

        \textbf{[CRITICAL NOTE]}
        The statement ``$\pi_4(S^2)=\mathbb{Z}_2$ guarantees a fermionic Hopfion'' is too strong.
        The homotopy group shows that a $\mathbb{Z}_2$ spacetime class exists, but hopfion
        statistics are set by the presence and value of the corresponding topological/statistical
        term.  In modern language this classification can involve bordism/TQFT data rather than
        only bare homotopy groups~\cite{ChenTanizaki2023}.  What is solid here is the reduction:
        \begin{equation}
            \theta=0 \Rightarrow \text{periodic/bosonic substrate},\qquad
            \theta=\pi \Rightarrow \text{spinorial/fermionic substrate}.
        \end{equation}
        The remaining open problem is to derive why the SST substrate realizes the
        $\theta=\pi$ sector rather than the $\theta=0$ sector.

        \textbf{[OPEN] Derivation of the spinorial boundary condition.}
        The spinorial anti-periodicity of the core phase is not yet derived from the swirl-field
        dynamics; it is the single remaining premise of the $pq+1$ ladder.

        \textbf{[CONDITIONAL] Unification with the magnetic moment.}
        The spinor premise above is the \emph{same} posit as the spin-$\tfrac12$ assignment of
        the magnetic-moment sector (Route~B): one spinorial lift of the core phase yields both
        $g=2$ and $\chi=\pm1\Rightarrow pq\pm1$. If established, fermion statistics, the
        gyromagnetic ratio, and the lepton self-linking ladder collapse into a single mechanism.

        \textbf{[DERIVED] Falsifier.}
        An independently computed $T(2,q)$ Hopfion self-linking with $Q_H\neq 2q\pm1$ (e.g.\
        $T(2,5)\neq 11$) falsifies the conditional chain irrespective of the variational
        argument.


\subsection{Spectral Structure of Locked Closed Carriers (Secondary Result)}
    \label{subsec:spectral_locked_carriers}
    
    \paragraph{Scope and status.}
        This subsection records a \emph{secondary consequence} of the kinematics of rationally locked, closed carriers. It does not introduce new primitives. All results are conditional on the existence of a well-defined phase field and a regular mode spectrum. Status: \textbf{derived under spectral assumptions; experimentally unverified}.

    \paragraph{Kinematic reduction.}
        Consider a closed carrier with a single central driving phase $\phi(t)$ and $N_s$ sector variables $\theta_i(t)$ $(i=1,\dots,N_s)$ subject to a rational locking relation
        \begin{equation}
            \theta_i = \sigma_i \kappa_i \phi + \delta_i,
            \qquad
            \sigma_i \in \{+1,-1\},\quad \kappa_i \in \mathbb{Q}.
            \label{eq:locking_relation}
        \end{equation}
        In the rigid-lock limit, the dynamics reduce to a single collective coordinate $\phi(t)$ with effective Lagrangian
        \begin{equation}
            \mathcal{L}_{\mathrm{eff}}
            =
            \frac{I_{\mathrm{eff}}}{2}\dot{\phi}^{\,2}
            -
            V(\phi),
            \label{eq:Leff_canonical}
        \end{equation}
        where $I_{\mathrm{eff}}$ is an effective inertia and $V(\phi)$ is a periodic locking potential.

    \paragraph{Distributed phase field.}
        On a closed carrier of length $L$, allow small phase fluctuations
        \begin{equation}
            \phi(t) \;\longrightarrow\; \phi_0 + \eta(s,t),
            \qquad s\in[0,L],
        \end{equation}
        with periodic boundary condition
        \begin{equation}
            \eta(s+L,t)=\eta(s,t).
        \end{equation}
        A minimal quadratic field model consistent with Eq.~\eqref{eq:Leff_canonical} is
        \begin{equation}
            \mathcal{L}_{\mathrm{field}}
            =
            \int_0^L
            \left[
                \frac{\mu_{\mathrm{eff}}}{2}(\partial_t \eta)^2
            -
                \frac{T_{\mathrm{eff}}}{2}(\partial_s \eta)^2
            -
                \frac{\mu_{\mathrm{eff}}\omega_0^2}{2}\eta^2
            \right] ds,
            \label{eq:field_model}
        \end{equation}
        where $\mu_{\mathrm{eff}}$ is an effective phase inertia density, $T_{\mathrm{eff}}$ an effective stiffness, and $\omega_0$ the small-oscillation frequency set by $V(\phi)$.
        
        Dimensional consistency:
        \[
            [\mu_{\mathrm{eff}}]=\mathrm{kg\,m},\qquad
            [T_{\mathrm{eff}}]=\mathrm{N},\qquad
            [\omega_0]=\mathrm{s^{-1}}.
        \]

    \paragraph{Mode spectrum.}
        Expanding in normal modes
        \begin{equation}
            \eta(s,t) = \sum_{n\in\mathbb{Z}} a_n(t)\, e^{i k_n s},
            \qquad
            k_n = \frac{2\pi n}{L},
        \end{equation}
        yields the dispersion relation
        \begin{equation}
            \boxed{
                \omega_n^2
                =
                \omega_0^2
                +
                c_\phi^2 k_n^2,
                \qquad
                c_\phi^2 = \frac{T_{\mathrm{eff}}}{\mu_{\mathrm{eff}}}.
            }
            \label{eq:dispersion}
        \end{equation}
        For large $n$,
        \begin{equation}
            \omega_n \sim \frac{2\pi c_\phi}{L}\, n.
            \label{eq:asymptotic_linear}
        \end{equation}

    \paragraph{Cubic spectral sums and \texorpdfstring{$\zeta(3)$}{zeta(3)}.}
        Consider a cubic-weighted spectral contribution of the form
        \begin{equation}
            \Delta \mathcal{E}^{(3)}
            =
            A_3 \sum_{n=1}^{\infty} \frac{1}{\omega_n^3},
            \qquad
            [A_3]=\mathrm{J\,s^{-3}}.
            \label{eq:cubic_sum}
        \end{equation}
        Using the asymptotic scaling \eqref{eq:asymptotic_linear},
        \begin{equation}
            \Delta \mathcal{E}^{(3)}
            \;\approx\;
            A_3
            \left(\frac{L}{2\pi c_\phi}\right)^3
            \sum_{n=1}^{\infty} \frac{1}{n^3}.
        \end{equation}
        Hence
        \begin{equation}
            \boxed{
                \Delta \mathcal{E}^{(3)}
                \;\approx\;
                A_3
                \left(\frac{L}{2\pi c_\phi}\right)^3
                \zeta(3).
            }
            \label{eq:zeta3_result}
        \end{equation}

    \paragraph{Interpretation.}
        Equation \eqref{eq:zeta3_result} shows that $\zeta(3)$ arises from the \emph{high-mode spectral tail} of a closed carrier, not from the geometric threefold structure alone. The role of the locking relation \eqref{eq:locking_relation} is indirect: it fixes $I_{\mathrm{eff}}$, $\omega_0$, and admissible boundary conditions, thereby determining the spectral scale $(L,c_\phi)$.

    \paragraph{Universality.}
        Different geometric realizations (e.g., distinct multi-sector carriers) that reduce to the same class of effective phase dynamics and share the asymptotic linear spectrum \eqref{eq:asymptotic_linear} will produce the same $\zeta(3)$ factor in Eq.~\eqref{eq:zeta3_result}, up to system-dependent prefactors.
        
        \begin{equation}
            \boxed{
                \text{$\zeta(3)$ is a spectral invariant of the mode tower, not a direct invariant of the carrier geometry.}
            }
        \end{equation}

    \paragraph{Limitations and falsifiers.}
        The result fails if:
        \begin{enumerate}
            \item the mode spectrum is not asymptotically linear in $n$,
            \item strong damping or nonlocal coupling invalidates the expansion \eqref{eq:field_model},
            \item the system does not admit a well-defined periodic phase field on a closed domain.
        \end{enumerate}

    \paragraph{Minimal experimental test.}
        Measure mode frequencies $\omega_n$ for $n=1,2,3,\dots$ and verify
        \begin{equation}
            \frac{\omega_n}{n} \to \text{constant}
            \quad\text{as } n\to\infty.
        \end{equation}
        Deviation from this scaling excludes Eq.~\eqref{eq:zeta3_result}.

    \paragraph{Status.}
        \textbf{Secondary derived result.} Valid under spectral assumptions; not a defining postulate. The cubic weighting $\sum_n 1/\omega_n^3$ with $\omega_n\sim n$ yields $\zeta(3)$ for \emph{any} asymptotically linear spectrum; this is a generic calculus property $[\text{CALIBRATED}]$, not an SST-specific fluid-dynamical or topological invariant.

\subsection{Symmetry-Sensitive Dark-Knot Classification}
    \label{subsec:dark_knot_classification}

    A compact classification retained from earlier SST topology layers distinguishes visible, quasi-dark, and dark sectors by chirality, helicity, and low-order instability content.

    Let $\chi(K)$ denote a chirality indicator, let
    \begin{align}
        \mathcal{H}[K] = \int \mathbf{v}\cdot(\nabla\times\mathbf{v})\,dV
    \end{align}
    be the helicity functional, and let $N_u^{(1)}(K)$ count low-order unstable modes in a chosen perturbative stability analysis. Then the retained taxonomy is

    \begin{align}
        \text{dark:}&\quad \chi(K)=0,\quad \mathcal{H}[K]\approx 0,\quad N_u^{(1)}(K)=0, \\
        \text{quasi-dark:}&\quad |\chi(K)|\ll 1,\quad \mathcal{H}[K]\approx 0,\quad 0 < N_u^{(1)}(K)\le 2, \\
        \text{visible:}&\quad \chi(K)\neq 0 \quad \text{or}\quad \mathcal{H}[K]\neq 0.
        \label{eq:dark_knot_taxonomy}
    \end{align}

    \textbf{[ORTHODOX]}: helicity is a standard invariant of ideal vortex dynamics \cite{Moffatt1969,Batchelor1967}.

    \textbf{[SPECULATIVE]}: the interpretation of Eq.~\eqref{eq:dark_knot_taxonomy} as a matter-sector taxonomy is an SST model choice.

    \paragraph{Heuristic balance diagnostic from the SST helicity archive.}
    As a supporting diagnostic (not a stand-alone classifier), define
    \begin{align}
        \Delta_H(K)=\abs{\widehat{\mathcal{H}}_{b}(K)+\frac{1}{2}},
    \end{align}
    where $\widehat{\mathcal{H}}_{b}(K)$ is the archived normalized helicity-like base indicator for class $K$.
    Near-balanced values with $\Delta_H(K)\ll 1$ are treated only as candidate evidence for dark or quasi-dark sectors.

    \textbf{[CRITICAL NOTE]}: in canon usage, the diagnostic above is subordinate to Eq.~\eqref{eq:dark_knot_taxonomy}; no single scalar is sufficient for dark-sector membership without the symmetry and instability filters.

    \textbf{[CANON WORKFLOW NOTE]}: if SST-labelled and parallel legacy-labelled helicity exports are numerically identical, only the SST-labelled archive is required for canon workflows.

    \paragraph{Symmetry-first refinement (v0.8.x).}
        In the taxonomy program, candidate sectors are labeled symmetry-first, then filtered by helicity balance and low-order instability content.
        Amphichiral knot types furnish the preferred \emph{candidate} pool for the dark bin in symmetry-first tables; assigning ``dark'' still requires the conjunction of near-zero chirality indicators, helicity balance in the sense of Eq.~\eqref{eq:dark_knot_taxonomy}, and vanishing low-order instability count.
        The implication ``amphichiral $\Rightarrow$ dark'' is a classifier proposal tested against observables and archive diagnostics, not a closed theorem.

    \paragraph{Ambiguity tie-break.}
        If rule conjunctions overlap, precedence is
        \[
            \text{symmetry gate} \;\to\; \text{helicity balance gate} \;\to\; \text{stability gate}.
        \]

\section{Swirl--Electromagnetic Bridge}
    \label{sec:embridge}

    \subsection{Minimal transverse field sector}

        A conservative bridge from the fluid sector to a radiation sector is obtained by introducing a divergence-free transverse field $\mathbf{A}_{\mathrm{eff}}$ with

        \begin{align}
            \nabla\cdot \mathbf{A}_{\mathrm{eff}} = 0,
            \label{eq:div_free_Aeff}
        \end{align}

        and effective Lagrangian density

        \begin{align}
            \mathcal{L}_{\mathrm{rad}}
            =
            \frac{\rhoF}{2}\,\lVert \partial_t \mathbf{A}_{\mathrm{eff}}\rVert^2
            -
            \frac{K_\gamma}{2}\,\lVert \nabla\times\mathbf{A}_{\mathrm{eff}}\rVert^2.
            \label{eq:rad_lagrangian}
        \end{align}

        Varying Eq.~\eqref{eq:rad_lagrangian} and using Eq.~\eqref{eq:div_free_Aeff} gives

        \begin{align}
            \frac{1}{c_\gamma^2}\partial_t^2 \mathbf{A}_{\mathrm{eff}} - \nabla^2 \mathbf{A}_{\mathrm{eff}} = 0,
            \qquad
            c_\gamma^2 = \frac{K_\gamma}{\rhoF}.
            \label{eq:vector_wave_Aeff}
        \end{align}

        Define

        \begin{align}
            \mathbf{E}_{\mathrm{eff}} = -\partial_t \mathbf{A}_{\mathrm{eff}},
            \qquad
            \mathbf{B}_{\mathrm{eff}} = \nabla\times \mathbf{A}_{\mathrm{eff}}.
            \label{eq:effective_em_fields}
        \end{align}

        Then one immediately recovers the source-free identities

        \begin{align}
            \nabla\cdot \mathbf{B}_{\mathrm{eff}} &= 0, \\
            \nabla\times \mathbf{E}_{\mathrm{eff}} &= -\partial_t \mathbf{B}_{\mathrm{eff}}.
            \label{eq:maxwell_like_identities}
        \end{align}

        \textbf{[ORTHODOX FORM]}: Eq.~\eqref{eq:vector_wave_Aeff} and Eq.~\eqref{eq:maxwell_like_identities} are the standard transverse-wave and source-free Faraday identities in vector-potential form \cite{Jackson1999}.

        \textbf{[SPECULATIVE IDENTIFICATION]}: SST identifies $\mathbf{A}_{\mathrm{eff}}$ with a coarse-grained torsional or director displacement field of the medium.

    \subsection{Photon / unknot sector}

        The minimal retained radiation picture is that an R-phase excitation is an unknotted, finite-duration torsional packet propagating in the field $\mathbf{A}_{\mathrm{eff}}$. At the Rosetta level this corresponds to a source-free transverse wave packet with dispersion inherited from Eq.~\eqref{eq:vector_wave_Aeff}.

        \textbf{[SPECULATIVE]}: handedness of the local torsional packet is taken to encode polarization, while additional azimuthal phase winding encodes orbital angular momentum.

        \textbf{[CANON / SEMANTIC RULE]}:
        a photon-sector R-phase torsion packet is not a loose end of a vortex
        filament and does not yet carry a protected closed matter-knot invariant.
        It is a phase/torsion support on a swirl-characteristic. A spacetime
        geodesic is only the orthodox comparison language for the coarse-grained
        propagation path.

    \subsection{R-to-T closure threshold and conservation bookkeeping}
        \label{subsec:rt_closure_threshold_bookkeeping}

        \textbf{[SPECULATIVE SST BRIDGE].}
        If electron-trefoil creation is retained, the conservative bookkeeping
        form is
        \begin{align}
            \gamma_{\rm R}
            +
            \mathcal{B}_{\rm atom}
            \overset{\mathcal{K}_{\rm Kairos}}{\longrightarrow}
            e^-_{3_1}
            +
            \mathcal{R},
            \label{eq:kairos_rt_trefoil_closure}
        \end{align}
        where \(\gamma_{\rm R}\) is an extended R-phase torsion packet,
        \(\mathcal{B}_{\rm atom}\) is a local boundary or trapping condition, and
        \(\mathcal{R}\) carries recoil, momentum balance, and any required
        topological or charge compensation. Equation~\eqref{eq:kairos_rt_trefoil_closure}
        is not a pure Euler reconnection event.

        The Compton anchoring gives the closure energy scale
        \begin{align}
            \omega_c = \frac{\vchar}{\rc},
            \qquad
            \hbar\omega_c = m_e c^2.
            \label{eq:compton_closure_energy_scale}
        \end{align}
        This scale may be used as a seed/closure threshold only in the presence
        of a boundary reservoir or pre-existing attachment support. Free creation
        of a matter carrier from radiation must respect charge, momentum, and
        topological compensation. A lone photon-to-lone-electron channel is
        therefore excluded in the canon-safe reading. The minimal free pair
        bookkeeping is instead
        \begin{align}
            \gamma+\gamma
            \longrightarrow
            3_1+\overline{3_1},
            \qquad
            E_{\rm pair}=2m_e c^2,
            \label{eq:trefoil_antitrefoil_pair_bookkeeping}
        \end{align}
        or a field-assisted version
        \begin{align}
            \gamma+\mathcal{F}
            \longrightarrow
            3_1+\overline{3_1}+\mathcal{F}' .
            \label{eq:field_assisted_pair_bookkeeping}
        \end{align}

    \subsection{Swirl pressure law}
        \label{subsec:swirl_pressure_law}

        For steady axisymmetric swirl, Euler radial balance gives

        \begin{align}
            \frac{1}{\rhoF}\frac{dp_{\mathrm{swirl}}}{dr} = \frac{v_\theta(r)^2}{r}.
            \label{eq:swirl_pressure_law}
        \end{align}

        For a locally rigid rotation profile $v_\theta(r)=\Omega r$, Eq.~\eqref{eq:swirl_pressure_law} integrates to

        \begin{align}
            p_{\mathrm{swirl}}(r) = p_0 + \frac{1}{2}\rhoF \Omega^2 r^2,
            \label{eq:swirl_pressure_rigid}
        \end{align}

        while for an exterior irrotational profile $v_\theta(r)=\Gamma/(2\pi r)$ one finds

        \begin{align}
            p_{\mathrm{swirl}}(r) = p_\infty - \frac{\rhoF\Gamma^2}{8\pi^2 r^2}.
            \label{eq:swirl_pressure_irrotational}
        \end{align}

        In a flat rotation curve with asymptotically constant azimuthal speed $v_\theta(r)\to v_0$ as $r$ grows in a window where the radial balance \eqref{eq:swirl_pressure_law} remains valid, integration yields the logarithmic tail
        \begin{align}
            p_{\mathrm{swirl}}(r)
            =
            p_0 + \rhoF v_0^2 \ln\!\left(\frac{r}{r_0}\right),
            \label{eq:swirl_pressure_flat_tail}
        \end{align}
        for suitable reference radius $r_0$ and offset $p_0$ within that window.

        \textbf{[ORTHODOX] Assumption block:} Eq.~\eqref{eq:swirl_pressure_law} is taken in the usual incompressible, inviscid, stationary, axisymmetric regime with azimuthal-dominant flow, so that radial drift, viscosity, and stratification corrections are subleading in the domain of validity.

        \textbf{[ORTHODOX] Dimensional check:}
        $\bigl[\rhoF^{-1}\,dp/dr\bigr]=\mathrm{m\,s^{-2}}=\bigl[v_\theta^2/r\bigr]$.

        \textbf{[ORTHODOX]}: Eq.~\eqref{eq:swirl_pressure_law} is the ordinary radial Euler balance for steady swirl \cite{Batchelor1967,Saffman1992}.

        \textbf{[DERIVED SST ROLE]}: this pressure law is the canonical bridge between localized circulation and large-scale force analogies. Interpreting $p_{\mathrm{swirl}}$ as a dark-sector classifier or support layer is an SST model step, logically separate from the orthodox fluid derivation.

        \textbf{[DERIVED] Observables / falsifiers:} (i) consistency of a pressure profile reconstructed from archived $v_\theta(r)$ with the radial gradient implied by \eqref{eq:swirl_pressure_law}; (ii) null test: systematic sign or scale mismatch of the reconstructed $dp_{\mathrm{swirl}}/dr$ relative to $v_\theta^2/r$ across validated radial windows excludes the applicability of the balance in those windows.

    \subsection{Lorentz-Type Force Density as Swirl Stress}
    \label{sec:lorentz_type_force_density_as_swirl_stress}

        \paragraph{Status labels.}
        The vector identity used in this section is \textbf{[ORTHODOX]}.  The
        definition of the swirl-force density is \textbf{[DERIVED]} as a hydrodynamic
        force-density structure.  The numerical SST scales are \textbf{[CALIBRATED]}
        with respect to the canonical constants
        \(\rhoF\), \(\rc\), and \(\vchar\).
        No microscopic electromagnetic element-force law is canonized here.

        \paragraph{Hydrodynamic identity.}
        For incompressible inviscid flow,
        \[
        \nabla\cdot\mathbf v=0,
        \qquad
        \boldsymbol{\omega}=\nabla\times\mathbf v .
        \]
        The convective acceleration satisfies the exact identity
        \[
        (\mathbf v\cdot\nabla)\mathbf v
        =
        \nabla\!\left(\frac12|\mathbf v|^2\right)
        -
        \mathbf v\times\boldsymbol{\omega}.
        \]
        Equivalently,
        \[
        \boxed{
        \mathbf v\times\boldsymbol{\omega}
        =
        \nabla\!\left(\frac12|\mathbf v|^2\right)
        -
        (\mathbf v\cdot\nabla)\mathbf v .
        }
        \]
        This provides the hydrodynamic structure behind a Lorentz-type transverse
        force-density term.

        \paragraph{Swirl-force density.}
        SST defines the local swirl-force density
        \[
        \boxed{
        \mathbf f_{\circlearrowleft}
        =
        \rhoF\mathbf v\times\boldsymbol{\omega}.
        }
        \]
        Its dimensional consistency is
        \[
        [\rhoF\mathbf v\times\boldsymbol{\omega}]
        =
        \mathrm{kg\,m^{-3}}\,
        \mathrm{m\,s^{-1}}\,
        \mathrm{s^{-1}}
        =
        \mathrm{N\,m^{-3}},
        \]
        as required for a force density.

        \paragraph{Bernoulli relation and sign convention.}
        The scalar quantity
        \[
        p_{\circlearrowleft}
        =
        \frac12\rhoF|\mathbf v|^2
        \]
        is the dynamic swirl-pressure scale.  It should not be confused with the
        static pressure, which decreases in high-swirl regions in the usual
        Bernoulli sense.  For stationary Euler flow with constant
        \(\rhoF\) and no external body force,
        \[
        \rhoF(\mathbf v\cdot\nabla)\mathbf v
        =
        -\nabla p_{\mathrm{stat}} .
        \]
        Using the identity above gives
        \[
        \boxed{
        \rhoF\mathbf v\times\boldsymbol{\omega}
        =
        \nabla\!\left(
        p_{\mathrm{stat}}
        +
        \frac12\rhoF|\mathbf v|^2
        \right).
        }
        \]
        Thus the swirl-force density is canonically tied to the gradient of total
        Euler pressure, not to an independently postulated electromagnetic force law.

        \paragraph{Closed-loop observables.}
        SST assigns canonical status to integrated and topological observables rather
        than to a unique microscopic segment-to-segment force kernel.  The canonical
        loop observables include circulation,
        \[
        \Gamma
        =
        \oint_C \mathbf v\cdot d\boldsymbol{\ell},
        \]
        helicity,
        \[
        H
        =
        \int_V \mathbf v\cdot\boldsymbol{\omega}\,d^3x,
        \]
        and the Gauss linking number,
        \[
        Lk(C_1,C_2)
        =
        \frac{1}{4\pi}
        \oint_{C_1}\oint_{C_2}
        \frac{
        (d\mathbf x_1\times d\mathbf x_2)\cdot(\mathbf x_1-\mathbf x_2)
        }{
        |\mathbf x_1-\mathbf x_2|^3
        }.
        \]
        Distinct local force kernels may yield identical closed-loop observables.
        Therefore SST canonizes circulation, helicity, linking number, pressure
        differences, and net forces, but not a unique microscopic electromagnetic
        element-force law.

        \paragraph{Canonical scales.}
        Using the canonical SST constants
        \[
        \vchar
        =
        1.09384563\times10^6\,\mathrm{m\,s^{-1}},
        \qquad
        \rc
        =
        1.40897017\times10^{-15}\,\mathrm m,
        \qquad
        \rhoF
        =
        7.0\times10^{-7}\,\mathrm{kg\,m^{-3}},
        \]
        the natural circulation scale is
        \[
        \Gamma_0
        =
        2\pi \rc\vchar
        =
        9.68361920\times10^{-9}\,\mathrm{m^2\,s^{-1}}.
        \]
        The associated dynamic swirl-pressure scale is
        \[
        p_{\circlearrowleft,0}
        =
        \frac12\rhoF\vchar{}^2
        =
        4.1877439\times10^{5}\,\mathrm{Pa}.
        \]
        This value is the canonical SST stress scale for Lorentz-type
        force-density analogies and coincides with the calibrated source-pressure
        scale used in the swirl-pressure sector.

        \paragraph{Canonical conclusion.}
        The canonized result is not a replacement of Maxwell--Lorentz electrodynamics.
        The canonized result is the hydrodynamic statement that a Lorentz-type
        transverse force-density structure has the exact SST form
        \[
        \boxed{
        \mathbf f_{\circlearrowleft}
        =
        \rhoF\mathbf v\times\boldsymbol{\omega}
        =
        \rhoF
        \nabla\!\left(\frac12|\mathbf v|^2\right)
        -
        \rhoF(\mathbf v\cdot\nabla)\mathbf v .
        }
        \]
        Accordingly, magnetic-type force densities in SST are interpreted as
        effective projections of pressure, vorticity, and closed-loop topology in the
        substrate.

    \subsection{Minimal effective Lagrangian bridge}

        A compact field-theoretic summary retaining the clock, radiation, and matter sectors is

        \begin{align}
            \mathcal{L}_{\mathrm{eff}}
            =
            \mathcal{L}_{T}
            +
            \mathcal{L}_{\mathrm{rad}}
            +
            \sum_K \bar\Psi_K\bigl(i\gamma^\mu D_\mu - m_K^{(\mathrm{sol})}\bigr)\Psi_K
            + \cdots,
            \label{eq:minimal_effective_lagrangian_bridge}
        \end{align}

        where $\mathcal{L}_{T}$ is the clock/foliation sector introduced in Section~\ref{sec:time}, $\mathcal{L}_{\mathrm{rad}}$ is given by Eq.~\eqref{eq:rad_lagrangian}, and the ellipsis denotes higher-order self-interaction, topological, or medium-response operators.

        \textbf{[SPECULATIVE BRIDGE]}: Eq.~\eqref{eq:minimal_effective_lagrangian_bridge} is a bookkeeping layer, not yet a completed first-principles derivation.

% ---------- DELAY THEORY ----------
%    \section{Delay and Mode Selection Theory}
%        \label{sec:delay}
        
\section{Delay and Mode Selection Theory}
    \label{sec:delay}
    
    \subsection{Circulation Delay as a Physical Primitive}
        
        We introduce the circulation delay as a constitutive element of the theory.
        For a closed vortex filament of effective length $L$ and characteristic propagation speed $\vchar$, the circulation time is defined as:
        
        \begin{align}
            \tau_{\rm circ}(K) = \frac{L_K}{\vchar}
            \label{eq:delay_time}
        \end{align}
        
        In the special case of a minimal closed loop associated with the neutral circulation scale, this reduces to:
        
        \begin{align}
            \tau_{\rm circ,c} = \frac{2\pi \rc}{\vchar} = \frac{2\pi}{\omega_c}
            \label{eq:compton_delay}
        \end{align}
        
        This establishes a direct identification between the circulation delay and the inverse Compton frequency.
        
        \textbf{[DERIVED]}: The intrinsic delay of the neutral circulation loop is identical to the Compton period.
    
    \subsection{Delayed Self-Interaction Dynamics}
        
        We model the phase $\phi(t)$ of a circulating excitation as a delayed self-interacting oscillator:
        
        \begin{align}
            \dot{\phi}(t) = \omega_0 + \kappa \sin\big(\phi(t - \tau_{\rm circ}) - \phi(t)\big)
            \label{eq:delay_phase}
        \end{align}
        
        This equation represents a minimal effective description of a circulating field with finite propagation time.
        
        \textbf{[ORTHODOX]}: Delay-differential equations of this type are well-established in feedback systems \cite{SST_delay_basic}.
    
    \subsection{Emergence of Discrete Stable Modes}
        
        We seek steady-state solutions of the form:
        
        \begin{align}
            \phi(t) = \Omega t
        \end{align}
        
        Substituting into Eq.~\ref{eq:delay_phase} yields $\phi(t-\tau_{\rm circ})-\phi(t)=-\Omega\tau_{\rm circ}$, hence:
        
        \begin{align}
            \Omega = \omega_0 - \kappa \sin(\Omega \tau_{\rm circ})
            \label{eq:mode_condition}
        \end{align}
        
        This transcendental equation admits a discrete set of stable solutions $\Omega_n$ determined by the delay parameter $\tau_{\rm circ}$. The sign convention is not physically unique---one may absorb it into $\kappa$ or reverse the delay argument---but the displayed form is the direct consequence of Eq.~\eqref{eq:delay_phase} as written.
        
        \textbf{[CONDITIONAL DERIVED]}: Given the effective DDE ansatz of Eq.~\eqref{eq:delay_phase}, the delayed feedback equation admits discrete locked branches. The claim that the SST medium obeys this specific DDE remains a bridge assumption, not a first-principles Euler theorem.
    
    \subsection{Delay-Induced Quantization Mechanism}
        
        Unlike conventional quantum mechanics, where discreteness is postulated, here it emerges dynamically:
        
        \begin{itemize}
            \item No boundary condition required
            \item No operator quantization assumed
            \item Discreteness arises from stability selection
        \end{itemize}
        
        \textbf{[SPECULATIVE]}: Particle spectra may be interpreted as delay-selected stable circulation modes.
    
    \subsection{Extension to Nonlinear Schrödinger Dynamics}
        
        To connect with field-theoretic descriptions, we extend the delay mechanism to a nonlinear Schrödinger framework:
        
        \begin{align}
            i \hbar \partial_t \psi = -\frac{\hbar^2}{2m} \nabla^2 \psi + F_{\text{delay}}[\psi(t - \tau_{\rm circ})]
        \end{align}
        
        where the delayed functional introduces temporal nonlocality into the field evolution.
        
        \textbf{[SPECULATIVE]}: This provides a bridge between classical delay systems and quantum wave dynamics.
    
    \subsection{Physical Interpretation}
        
        The delay $\tau_{\rm circ}$ acts as:
        
        \begin{itemize}
            \item a memory kernel
            \item a nonlocal temporal coupling
            \item a selection mechanism for stable modes
        \end{itemize}
        
        Thus:
        
        \begin{center}
            \textit{Delay $\Rightarrow$ Stability $\Rightarrow$ Discreteness}
        \end{center}
        
        This mechanism forms one of the central pillars of the SST canon.

%    \section{Atomic Bridge Model}
%        \label{sec:atomic}
        
\section{Atomic Bridge Model}
    \label{sec:atomic}

    \subsection{Introduction and epistemic framing}
        This section records the minimal SST atomic bridge rather than a full derivation of atomic structure. Its canon role is deliberately narrow: recover the orthodox Coulomb limit when the SST-specific sector is switched off, define a branch-resolved circulation coupling fixed by canonical medium scales, and make explicit which many-electron and shell-structure claims remain open. It is therefore a benchmark layer, not yet a complete atomic theory.
    
    \subsection{Baseline Hamiltonian}
        
        We adopt the standard Coulomb Hamiltonian as the orthodox baseline:
        
        \begin{align}
            H_0 = -\frac{\hbar^2}{2\mu} \nabla^2 - \frac{Ze^2}{4\pi \epsilon_0 r}
            \label{eq:coulomb}
        \end{align}
        
        \textbf{[ORTHODOX]}: This reproduces the hydrogenic spectrum.

    \subsection{Legacy Soft-Core Regularization}

        Early canon layers introduced a finite-core regularization of the hydrogenic potential. In the present document this is retained only as a controlled bridge ansatz, not as a replacement for Eq.~\eqref{eq:coulomb}.

        Define the effective coupling scale
        \begin{align}
            \Lambda_{\mathrm{SST}} = 4\pi \, \rhocore \, \vchar{}^{2} \, \rc^{4}
            \label{eq:lambda_sst}
        \end{align}
        with dimensions
        \begin{align}
            [\Lambda_{\mathrm{SST}}] = \mathrm{J\,m} = \mathrm{N\,m^2}.
        \end{align}

        A finite-core bridge potential is then
        \begin{align}
            V_{\mathrm{soft}}(r) = -\frac{\Lambda_{\mathrm{SST}}}{\sqrt{r^2+r_c^2}}
            \label{eq:softcore_potential}
        \end{align}
        satisfying the large-radius limit
        \begin{align}
            V_{\mathrm{soft}}(r) \xrightarrow{r\gg r_c} -\frac{\Lambda_{\mathrm{SST}}}{r}.
        \end{align}

        \textbf{[ORTHODOX]} Effective-model input: finite-core regularizations of singular central potentials are standard tools.\newline
        \textbf{[DERIVED]} Canonical consequence: the SST scale in Eq.~\eqref{eq:lambda_sst} is fixed by canonical constants rather than introduced phenomenologically.\newline
        \textbf{[ORTHODOX] Constraint}: this sector is admissible only if the large-radius limit reproduces the standard Coulomb problem to spectroscopic accuracy.

        In particular, Eq.~\eqref{eq:softcore_potential} is retained as a bridge ansatz only. It is not promoted here to the primary atomic potential.

    \subsection{SST Circulation Correction}
        
        We introduce a circulation-dependent perturbation:
        
        \begin{align}
            \delta H^{(\sigma)} = -2\sigma E_R \eta_0 \gamma + \frac{E_R \eta_0 \sigma}{2} \left\{ \frac{\hat{L}_z}{\hbar}, \gamma \right\}
            \label{eq:sst_correction}
        \end{align}
        
        \textbf{[BRIDGE ANSATZ]}: This is an orthodox-operator perturbation template whose coefficient depends on the SST circulation amplitude $\gamma(r,t)$. Because \(\hat L_z\) is imported from the standard quantum Hamiltonian, Eq.~\eqref{eq:sst_correction} is not a first-principles derivation of quantum angular momentum from the substrate.

        \textbf{[CRITICAL NOTE]}: The fractional shift scales as $\delta E_n/E_n \sim \eta_0\gamma=(\alpha/2)\gamma$. With $\gamma=\mathcal{O}(1)$ this is $\mathcal{O}(\alpha)$, far larger than the orthodox fine-structure splitting $\mathcal{O}(\alpha^2)$ (fractional). Compatibility with the atomic spectroscopy bounds (Sec.~7) therefore requires $\gamma$ to be $\alpha$-suppressed; this suppression is not yet derived.

    \subsection{Circulation source chain}
        The atomic perturbation is not introduced as an arbitrary extra scalar, but is intended to descend from a circulation-bearing source chain,
        \begin{align}
            \xi \rightarrow \omega_z \rightarrow u_\theta \rightarrow \Gamma_{\mathrm{loc}} \rightarrow \gamma,
        \end{align}
        where an axial displacement field $\xi$ generates local vorticity, the vorticity sources azimuthal swirl, and the loop circulation is normalized by the canonical quantum $\Gamma_0$. This makes the circulation amplitude $\gamma$ the natural canon-level coupling variable rather than a free phenomenological insertion.
    
    \subsection{Branch Structure}
        
        Two chiral sectors exist:
        
        \begin{align}
            \sigma = \pm 1
        \end{align}
        
        These correspond to:
        
        \begin{itemize}
            \item $\swirlarrow$ (left-handed)
            \item $\swirlarrowcw$ (right-handed)
        \end{itemize}
    
    \subsection{Electron trefoil identity under ordinary atomic binding}
        \label{subsec:electron_trefoil_atomic_binding}

        \textbf{[DERIVED SST CONSTRAINT].}
        Ordinary propagation through vacuum and ordinary atomic binding do not
        change the electron carrier's knot class. In the trefoil assignment, the
        transition is
        \begin{align}
            3_1^{\rm free}
            \longrightarrow
            3_1^{\rm bound},
            \label{eq:free_to_bound_trefoil}
        \end{align}
        not a breaking or unknotting of the carrier. Binding changes the embedding,
        phase envelope, branch coupling, and dressing of the carrier:
        \begin{align}
            e^-_{\rm atom}
            =
            3_1^{\rm core}
            \otimes
            \Psi_{n\ell m\sigma}^{\rm envelope}.
            \label{eq:atomic_electron_trefoil_tensor_envelope}
        \end{align}

        \textbf{[CANON RULE].}
        Ordinary electron capture by an ion is a binding and phase-locking process.
        Destruction or conversion of the separate electron trefoil requires a
        non-Euler process such as annihilation with an anti-trefoil or weak-sector
        capture, with the necessary compensation carried by the final state.

    \subsection{Consistency Requirement}
        
        To remain compatible with atomic physics:
        
        \begin{align}
            |\delta E_n| \ll |E_n|
        \end{align}
        
        \textbf{[ORTHODOX] Constraint}: SST corrections must not violate known spectroscopy.
    
    \subsection{Interpretation}
        
        \textbf{[STATUS]}: The atomic bridge is a controlled benchmark sector. The orthodox Coulomb problem is imported as the baseline; the circulation-dependent term is the SST addition; shell counting, exclusion-like behavior, and heavy-atom extensions remain programmatic rather than derived here. The purpose of the section is therefore structural compatibility, not yet a complete replacement of atomic quantum mechanics.

%    \section{Spectroscopic Constraints}
%        \label{sec:spectroscopy}
        
\section{Spectroscopic Constraints}
    \label{sec:spectroscopy}
    
    \subsection{Motivation}
        
        Any extension of atomic structure must remain consistent with high-precision spectroscopic measurements.
        
        \textbf{[ORTHODOX]}:
        Atomic energy levels are experimentally verified to extremely high precision.
        
        \textbf{Requirement:}
        \begin{align}
            |\delta E_n| \ll |E_n|
            \label{eq:spectro_constraint}
        \end{align}
        
        where $\delta E_n$ represents any SST-induced correction.
    
    \subsection{Internal Mode Hypothesis}
        
        We consider the possibility that atomic bound states possess an additional internal excitation sector associated with vortex filament dynamics.
        
        \textbf{[SPECULATIVE]}:
        Each bound state carries an internal mode spectrum $\{E_{m,n}\}$.
        
        The effective energy becomes:
        
        \begin{align}
            E_n^{\text{eff}} = E_n^{(0)} + \delta E_n
            \label{eq:effective_energy}
        \end{align}
    
    \subsection{Generic Coupling Structure}
        
        We parameterize the coupling between internal modes and observable energy levels:
        
        \begin{align}
            |\delta E_n| \leq \lambda_n \sum_m E_{m,n} \, \nu_{m,n}
            \label{eq:coupling_bound}
        \end{align}
        
        where:
        \begin{itemize}
            \item $\lambda_n$ is the coupling efficiency
            \item $\nu_{m,n}$ is the occupation factor
        \end{itemize}
        
        \textbf{[DERIVED]}: This form captures a general upper bound independent of microscopic details.
    
    \subsection{Ungapped Spectrum: Instability}
        
        If the internal spectrum is ungapped:
        
        \begin{itemize}
            \item arbitrarily low-energy excitations exist
            \item no universal suppression mechanism applies
        \end{itemize}
        
        \textbf{Consequence:}
        \begin{align}
            \delta E_n \not\to 0 \quad \text{generically}
        \end{align}
        
        \textbf{[CRITICAL]}:
        An ungapped internal sector is generically incompatible with spectroscopy unless:
        
        \begin{itemize}
            \item $\lambda_n \ll 1$ (extremely weak coupling), or
            \item the density of states is highly suppressed
        \end{itemize}
    
    \subsection{Gapped Spectrum: Decoupling Mechanism}
        
        We introduce an excitation gap $\Delta_K$:
        
        \begin{align}
            E_{m,n} \geq \Delta_K
        \end{align}
        
        Then occupation factors follow Boltzmann suppression:
        
        \begin{align}
            \nu_{m,n} \sim \exp\left(-\frac{\Delta_K}{k_B T_{\text{eff}}}\right)
        \end{align}
        
        Substituting into Eq.~\ref{eq:coupling_bound}:
        
        \begin{align}
            \delta E_n \propto \exp\left(-\frac{\Delta_K}{k_B T_{\text{eff}}}\right)
            \label{eq:boltzmann_suppression}
        \end{align}
        
        \textbf{[DERIVED]}:
        Exponential suppression ensures compatibility with spectroscopy.
    
    \subsection{Physical Origin of the Gap}
        
        \textbf{[SPECULATIVE]}:
        The gap $\Delta_K$ may arise from:
        
        \begin{itemize}
            \item geometric constraints on filament curvature
            \item torsional rigidity
            \item topological restrictions
            \item core structure of the vortex
        \end{itemize}
    
    \subsection{Order-of-Magnitude Estimate}
        
        Using filament parameters:
        
        \begin{align}
            \Delta_K \sim \frac{\Gamma^2}{4\pi \xi L}
        \end{align}
        
        where:
        \begin{itemize}
            \item $\xi$ is the core scale
            \item $L$ is the loop length
        \end{itemize}
        
        For typical SST scales:
        
        \begin{align}
            \Delta_K = \mathcal{O}(10^2 - 10^3 \, \text{eV})
        \end{align}
        
        \textbf{[DERIVED]}:
        This places the internal sector well above atomic energy scales ($\sim$ eV).
    
    \subsection{Consistency Condition}
        
        To preserve atomic physics:
        
        \begin{align}
            \frac{\Delta_K}{k_B T_{\text{eff}}} \gg 1
        \end{align}
        
        \textbf{Interpretation:}
        The internal sector is effectively frozen at atomic energies.
    
    \subsection{Observable Consequences}
        
        \textbf{[SPECULATIVE]}:
        Residual effects may include:
        
        \begin{itemize}
            \item tiny level shifts
            \item suppressed line broadening
            \item weak temperature dependence
        \end{itemize}
    
    \subsection{Falsifiability}
        
        The theory is falsifiable via:
        
        \begin{itemize}
            \item detection of unexplained spectral shifts
            \item absence of expected suppression
            \item deviation from Boltzmann scaling
        \end{itemize}
    
    \subsection{Integration with SST Framework}
        
        The spectroscopic constraint enforces:
        
        \begin{itemize}
            \item compatibility with the Atomic Bridge (Section~\ref{sec:atomic})
            \item constraints on delay-induced excitations (Section~\ref{sec:delay})
            \item bounds on allowable vortex configurations
        \end{itemize}
    
    \subsection{Canonical Statement}
        
        \begin{center}
            \textit{
                Any viable SST extension must include a mechanism—such as an excitation gap—that ensures exponential suppression of internal modes at atomic energy scales.
            }
        \end{center}

    \subsection{Precision electroweak bridge: lessons from the CMS \texorpdfstring{$W$}{W}-mass measurement}
        \label{subsec:cms-w-mass-bridge}
    
        \paragraph{Status.}
            \textbf{[ORTHODOX]} External input: the CMS measurement is a precision constraint on any SST effective weak-sector realization.
            \textbf{[DERIVED]} Methodological consequence: the result supports a forward-modelling strategy for SST phenomenology.
            \textbf{[CRITICAL NOTE]} The paper does \emph{not} provide direct evidence for a structured-substrate ontology of the $W$ boson.
    
        \paragraph{Empirical anchor.}
            Let the experimentally inferred resonance mass be denoted by
            \begin{equation}
                m_W^{\rm exp} = 80.3602 \pm 0.0099~{\rm GeV}.
            \end{equation}
            In the CMS analysis this quantity is not read off from a single observable, but extracted from a profiled likelihood over finely binned charged-lepton kinematics. This motivates the following SST principle:
            
            \begin{quote}
                \emph{Collider observables constrain SST through forward maps from latent swirl-string states to detector-level distributions, not through direct identification of a single internal parameter with a measured mass.}
            \end{quote}
    
        \paragraph{Latent-to-observable map.}
            Introduce a latent SST parameter set for an effective spin-1 weak excitation,
            \begin{equation}
                \Theta_{W}^{\rm SST}
                =
                \Bigl(
                K_W,\,
                \chi_W,\,
                \mathcal{L}_W,\,
                \mathcal{H}_W,\,
                \SwirlClock_W,\,
                \rhoF^{\rm eff},\,
                \uswirl^{\rm eff}
                \Bigr),
            \end{equation}
            where:
            \begin{itemize}
                \item $K_W$ is the topology class or excitation family,
                \item $\chi_W$ is chirality / handedness label,
                \item $\mathcal{L}_W$ is an effective geometric length scale,
                \item $\mathcal{H}_W$ is a helicity/linkage content,
                \item $\SwirlClock_W$ is the local clock factor for the mode,
                \item $\rhoF^{\rm eff}$ and $\uswirl^{\rm eff}$ are the coarse-grained medium density and effective local swirl-velocity field seen by the excitation.
            \end{itemize}
            
            The experimentally accessible distribution is then not a bare function of $m_W$ alone, but a projected density
            \begin{equation}
                \mathcal{P}_{\rm meas}(x)
                =
                \int dz\; R(x|z)\,\mathcal{P}_{\rm prod}(z\mid \Theta_W^{\rm SST}),
                \label{eq:sst-forward-map}
            \end{equation}
            where $z$ denotes truth-level production/decay variables and $R(x|z)$ is the detector response kernel.
    
        \paragraph{Mass as an effective resonance parameter.}
            In SST one should distinguish:
            \begin{equation}
                m_W^{\rm internal}
                \neq
                m_W^{\rm eff}
            \end{equation}
            in general, where:
            \begin{itemize}
                \item $m_W^{\rm internal}$ is the internal excitation-energy scale of the structured mode,
                \item $m_W^{\rm eff}$ is the resonance parameter inferred from scattering distributions.
            \end{itemize}
            The experimentally constrained quantity is $m_W^{\rm eff}$.
            
            A minimal canon-safe matching condition is therefore
            \begin{equation}
                m_W^{\rm eff}(\Theta_W^{\rm SST}) = m_W^{\rm exp} + \delta_W,
                \qquad
                |\delta_W| \lesssim 10~{\rm MeV},
                \label{eq:mw-match}
            \end{equation}
            with the current precision scale set by experiment.
    
        \paragraph{Distribution-level fit target.}
            The correct SST comparison object is a likelihood built from binned observables:
            \begin{equation}
                \mathcal{L}\bigl(\Theta_W^{\rm SST},\nu\bigr)
                =
                \prod_{b=1}^{N_{\rm bins}}
                    {\rm Poisson}\!\left(
                                       n_b \,\middle|\, \mu_b(\Theta_W^{\rm SST},\nu)
                \right)
                \,
                \Pi(\nu),
                \label{eq:sst-likelihood}
            \end{equation}
            where:
            \begin{itemize}
                \item $n_b$ is the observed count in bin $b$,
                \item $\mu_b$ is the SST prediction after production, decay, and detector folding,
                \item $\nu$ denotes nuisance parameters,
                \item $\Pi(\nu)$ is the nuisance prior / constraint factor.
            \end{itemize}
            The profile likelihood for the physically interesting SST parameters is then
            \begin{equation}
                \mathcal{L}_{\rm prof}(\Theta_W^{\rm SST})
                =
                \max_{\nu}\,
                \mathcal{L}(\Theta_W^{\rm SST},\nu).
                \label{eq:profile-like}
            \end{equation}
    
        \paragraph{Helicity decomposition as a falsifier.}
            A particularly sharp SST test is not the central value of $m_W$ alone, but whether the effective spin-1 mode reproduces the standard angular basis. Write
            \begin{equation}
                \frac{d\sigma}{d\Omega}
                =
                \sum_{i=0}^{8} A_i\, f_i(\theta,\phi)
                +
                \epsilon_{\rm SST}\,\Delta(\theta,\phi),
                \label{eq:helicity-decomp-sst}
            \end{equation}
            where the $A_i f_i$ are the standard helicity structures and $\epsilon_{\rm SST}\Delta$ is the first SST correction. Then:
            \begin{itemize}
                \item if $\epsilon_{\rm SST}=0$ within experimental precision, SST survives this channel,
                \item if $\epsilon_{\rm SST}\neq 0$, the deviation is directly falsifiable through angular distributions.
            \end{itemize}
    
        \paragraph{Canon interpretation.}
            This leads to the following conservative reading:
            \begin{enumerate}
                \item The weak sector must be \emph{SM-like at the effective distribution level}.
                \item Large $\mathcal{O}(10^{-4})$--$\mathcal{O}(10^{-3})$ relative distortions in the effective $W$ mass are disfavoured.
                \item The best SST discovery channels are therefore likely to be
                \begin{equation}
                    \text{small angular deviations},\quad
                    \text{polarization asymmetries},\quad
                    \text{off-shell tails},\quad
                    \text{rare differential distortions},
                \end{equation}
                rather than a large shift in the central value of $m_W$.
            \end{enumerate}
    
        \paragraph{Minimal bridge claim.}
            The CMS result does \emph{not} force the $W$ boson to be fundamental in the ontological sense. It does force any SST realization of the weak sector to satisfy the stronger requirement:
            \begin{equation}
                \mathcal{P}_{\rm prod}^{\rm SST}(z\mid \Theta_W^{\rm SST})
                \approx
                \mathcal{P}_{\rm prod}^{\rm SM}(z\mid m_W,\text{SM nuisances})
            \end{equation}
            to current experimental precision after detector folding and nuisance profiling.
    
        \paragraph{Practical SST program.}
            A precision-compatible SST weak-sector program should therefore proceed in three layers:
            \begin{align}
                \text{(i) effective resonance matching:}&\quad
                m_W^{\rm eff} \to m_W^{\rm exp}, \\
                \text{(ii) helicity matching:}&\quad
                A_i^{\rm SST} \to A_i^{\rm SM}, \\
                \text{(iii) residual search:}&\quad
                \Delta \neq 0 \text{ only in subleading differential structure.}
            \end{align}
    
        \paragraph{Conclusion.}
            The CMS $W$-mass result is best understood in SST as a \emph{precision filter}:
            it does not supply ontology, but it strongly constrains phenomenology. Any viable SST weak-sector realization must be distributionally SM-like to present precision, while leaving room only for small, structured residuals in higher-order observables.

% ---------- TIME ----------
    \section{Relational Time Framework}
        \label{sec:time}

    \subsection{Canonical Time Ontology}
        \label{subsec:canonical_time_ontology}

        \paragraph{Status.}
        \textbf{[CANON / SEMANTIC ORGANIZATION].}
        This subsection fixes the time-language of SST.  It does not introduce a
        new dynamical postulate.  Its purpose is to prevent the absolute foliation
        parameter, laboratory time, local proper time, carrier loop time, the
        Swirl-Clock factor, and Kairos transition events from being conflated.

        \paragraph{No bare-\(\tau\) convention.}
        Because SST uses both relativistic local clock accumulation and finite
        circulation delay, the unqualified symbol \(\tau\) is avoided in mixed
        contexts.  Circulation delay is denoted by
        \(\tau_{\rm circ}\), while local proper time carried by a localized
        process \(A\) is denoted by \(\tau_A\) or \(\tau_{\rm loc}\).

        \paragraph{Absolute foliation time.}
        The symbol \(T\) denotes the global foliation parameter of the
        incompressible SST substrate.  Each value of \(T\) labels a constraint
        slice
        \begin{align}
            \Sigma_T .
        \end{align}
        The pressure field and incompressibility constraint are reconstructed on
        \(\Sigma_T\).  The parameter \(T\) is bookkeeping time for the medium
        foliation; it is not by itself a local clock reading.  This foliation
        \(T\) is distinct from the coarse-grained clock field \(T(x)\) introduced
        later in relational observables from conserved currents.

        \paragraph{External comparison time.}
        The symbol \(t_{\infty}\) denotes asymptotic laboratory time or external
        comparison time.  It is the time coordinate used when matching SST clock
        predictions to standard relativistic observables.  In weakly perturbed
        laboratory settings one may identify \(t_{\infty}\) with \(T\), but the
        distinction should be retained whenever local clock rates are being
        discussed.

        \paragraph{Swirl-Clock factor.}
        The local Swirl-Clock factor is the dimensionless rate-conversion field
        already defined in Eq.~\eqref{eq:swirl_clock}:
        \begin{align}
            \SwirlClock(x,T)
            =
            \sqrt{
                1
                -
                \frac{
                    \left\|
                    \uswirl(x,T)
                    \right\|^2
                }{c^2}
            } .
            \label{eq:canonical_time_ontology_swirl_clock}
        \end{align}
        It is not a duration.  It converts the foliation rate into the local
        proper-clock rate of a resolved carrier or process.

        \paragraph{Local proper time.}
        The symbol \(\tau_A\) (or \(\tau_{\rm loc}\)) denotes local proper clock time.  Along a localized
        carrier \(A\),
        \begin{align}
            d\tau_A
            =
            \SwirlClock[A;T]\,dT .
            \label{eq:canonical_time_ontology_local_tau}
        \end{align}
        Thus two localized carriers \(A\) and \(B\) may share the same foliation
        slice \(\Sigma_T\) while accumulating different proper times:
        \begin{align}
            \frac{d\tau_A}{dT}=\SwirlClock[A;T],
            \qquad
            \frac{d\tau_B}{dT}=\SwirlClock[B;T],
            \qquad
            \frac{d\tau_C}{dT}=\SwirlClock[C;T].
            \label{eq:absolute_T_local_tau_layers}
        \end{align}

        \paragraph{Proper loop or turnover time.}
        The symbol \(\tau_{\rm circ}\) denotes a carrier-internal circulation or
        delay time.  For a loop of resolved length \(L_K\),
        \begin{align}
            \tau_{\rm circ}(K)
            =
            \frac{L_K}{\vchar}.
            \label{eq:canonical_time_ontology_loop_time}
        \end{align}
        At the canonical core scale this reduces to
        \begin{align}
            \tau_{\rm circ,c}
            =
            \frac{2\pi \rc}{\vchar}
            =
            \frac{2\pi}{\omega_c}.
        \end{align}
        The turnover time \(\tau_{\rm circ}\) counts internal circulation cycles; it is
        therefore distinct from the dimensionless Swirl-Clock factor \(\SwirlClock\).

        \paragraph{Kairos transition event.}
        A Kairos event is not a continuous clock variable.  It is a discrete
        topological transition marker,
        \begin{align}
            \mathcal{K}_{\rm Kairos}:
            K_-
            \longrightarrow
            K_+ ,
            \label{eq:canonical_time_ontology_kairos}
        \end{align}
        used when a process changes topological sector, chirality class, closure
        state, or reconnection class.  It may occur at some foliation value \(T\),
        but it is not itself a time coordinate.

        \paragraph{Summary schema.}
        \begin{center}
        \begin{tabular}{lll}
        \hline
        Layer & Symbol & Role \\
        \hline
        Foliation slice & \(\Sigma_T\) & global incompressible constraint slice \\
        Absolute foliation time & \(T\) & ordering parameter of the substrate \\
        External comparison time & \(t_{\infty}\) & laboratory/asymptotic comparison time \\
        Swirl-Clock factor & \(\SwirlClock\) & dimensionless local rate conversion \\
        Local proper time & \(\tau_A\), \(\tau_{\rm loc}\) & accumulated local clock time \\
        Loop/circulation delay & \(\tau_{\rm circ}\) & internal circulation period \\
        Kairos event & \(\mathcal{K}_{\rm Kairos}\) & discrete topological transition \\
        \hline
        \end{tabular}
        \end{center}

        \paragraph{Units guard.}
        \(\SwirlClock\) is dimensionless; \(T\), \(t_{\infty}\), \(\tau_A\), and
        \(\tau_{\rm circ}\) have units of time; \(\mathcal{K}_{\rm Kairos}\) is an
        event, not a coordinate.

    \subsection{Swirl-Clock and local proper time}
        Section~\ref{subsec:canonical_time_ontology} fixes the time-language of SST.
        The canonical time variable is not introduced as an independent primitive. It is defined relationally from the local kinematic state of the medium through the Swirl-Clock already introduced in Eq.~\eqref{eq:swirl_clock}. In its coarse-grained form,
        \begin{align}
            d\tau_{\rm loc} = \SwirlClock(x)\,dt_{\infty},
        \end{align}
        consistent with Eq.~\eqref{eq:canonical_time_ontology_local_tau}.  Here $t_{\infty}$ denotes the asymptotic laboratory time and $\tau_{\rm loc}$ the effective proper time associated with a localized process in the medium. In SST, this factor is interpreted as a fluid-kinematic clock field rather than as a purely geometric spacetime postulate.

        \paragraph{No-signalling research-track note.}
        The pressure field \(p_{\swirlarrow}(\mathbf{x},T)\) is a constraint
        field on \(\Sigma_T\).  Its elliptic reconstruction may be instantaneous
        in the \(T\)-bookkeeping sense, but this does not by itself define an
        operational superluminal signal channel.  Observable clock and photon
        phases are read through the local Swirl--Clock projection
        \(d\tau_A/dT=\SwirlClock[A;T]\).  The no-signalling status of the pressure
        constraint is a research-track theorem target, not an assumed device
        capability.

    \subsection{Carrier-local clocking and the turnover scale}
        For localized carriers the natural microscopic rate is the turnover frequency
        \begin{align}
            \Omega_0 = \frac{\vchar}{\rc} = \omega_c,
        \end{align}
        which coincides with the Compton anchoring scale and satisfies
        \(\tau_{\rm circ,c}=2\pi/\omega_c\) (Eq.~\eqref{eq:compton_delay}).
        A complete SST clock description therefore involves two linked but distinct levels:
        \begin{itemize}
            \item the \emph{microscopic turnover scale} $\Omega_0$, which counts circulation cycles of the carrier,
            \item the \emph{coarse-grained Swirl-Clock} $S_{(t)}$, which relates local rates to asymptotic time.
        \end{itemize}
        This resolves a recurring ambiguity in earlier canon versions: local cycle counting and coarse-grained time dilation are related, but they are not identical observables.

    \subsection{Relational observables from conserved currents}
        A purely operator-based notion of time is avoided in SST by defining time observables through conserved currents. Let $J^\mu$ denote a matter current and let $j_{\mathrm{ev}}^\mu$ denote a coarse-grained event current associated with structured activity of the medium. The relational clock field $T(x)$ is then recovered only after coarse graining,
        \begin{align}
            \partial_\mu T(x) \approx \frac{1}{\nu_0}\,\langle j^{\mathrm{ev}}_\mu \rangle_{\ell},
        \end{align}
        where $\ell$ is the coarse-graining scale and $\nu_0$ a reference event frequency. The canon interpretation is that time is read from correlations between currents, not from an abstract external operator.

    \subsection{Time-of-arrival as a relational field observable}
        Let $\mathcal{W}$ be a detector world-tube with boundary $\Sigma = \partial \mathcal{W}$. Then the arrival-time distribution associated with a detector label $\Theta$ is written as
        \begin{align}
            p(\Theta) = \frac{1}{\mathcal{N}} \int_{\Sigma} d\sigma_\mu \, J^\mu(x)\,\delta\!\big(T(x)-\Theta\big).
        \end{align}
        This formulation is canonically important because it bypasses the need for a self-adjoint time operator conjugate to a bounded Hamiltonian. The measured time is instead the value of a physical clock field at the point where matter flux crosses the detector boundary.

    \subsection{Clock broadening and continuum limits}
        If the event current has a finite variance, then the observed time-of-arrival distribution is broadened relative to the ideal classical prediction. Writing $P_{\mathrm{cl}}$ for the classical arrival distribution and $\sigma_\tau^2$ for the effective clock variance, one expects a convolution structure of the form
        \begin{align}
            P_{\mathrm{obs}}(\Theta)
            =
            \int dt\,P_{\mathrm{cl}}(t)
            \frac{1}{\sqrt{2\pi\sigma_\tau^2}}
            \exp\!\left[-\frac{(\Theta-t)^2}{2\sigma_\tau^2}\right].
        \end{align}
        In the limit $\ell \gg r_c$, the event structure is smoothed and one recovers a continuum clock field. In the opposite limit $\ell \to r_c$, the underlying discreteness of carrier events becomes operationally relevant.

    \subsection{Connection to gravity and geometry}
        In the canon, the relational-time sector and the geometric-response sector are not independent. Variations in the Swirl-Clock track variations in local swirl energy density and therefore provide the time component of the effective geometric response of the medium. At this level, the SST reading is:
        \begin{center}
            \textit{gravity and time are two coarse-grained aspects of the same clock-bearing medium.}
        \end{center}
        This statement should be read as a canon-level organizing principle, not yet as a completed field-theoretic proof.

    \subsection{Clock Field and Foliation Variable}

        \textbf{[DERIVED]} EFT bridge: the scalar Swirl-Clock field is promoted from the kinematic relation
        of Eq.~\eqref{eq:swirl_clock} to a long-wavelength foliation variable
        \begin{align}
            \chi(x),
        \end{align}
        whose normalized gradient defines a preferred local clock frame:
        \begin{align}
            u_\mu
            =
            \frac{\nabla_\mu \chi}
            {\sqrt{g^{\alpha\beta}\nabla_\alpha\chi\,\nabla_\beta\chi}}.
            \label{eq:clock_foliation_vector}
        \end{align}

        \textbf{[ORTHODOX]} Equation~\eqref{eq:clock_foliation_vector} matches the standard hypersurface-orthogonal
        foliation construction used in khronometric / Einstein--\AE ther effective field theory
        \cite{Jacobson2008,BlasPujolasSibiryakov2011}.

        \textbf{[DERIVED]} Within SST, the interpretation is different: \(u_\mu\) is not taken as a fundamental
        spacetime structure, but as the coarse-grained clock-flow direction associated with the medium.

    \subsection{Clock-Sector Effective Action}

        \textbf{[ORTHODOX]} EFT template: the most general second-order covariant clock-sector action compatible
        with hypersurface orthogonality is
        \begin{align}
        S_\chi
        =
        \int d^4x \sqrt{-g}\,
        \Big[
        & c_1 (\nabla_\mu u_\nu)(\nabla^\mu u^\nu)
        + c_2 (\nabla_\mu u^\mu)^2 \nonumber\\
        & + c_3 (\nabla_\mu u_\nu)(\nabla^\nu u^\mu)
        + c_4 u^\mu u^\nu (\nabla_\mu u_\alpha)(\nabla_\nu u^\alpha)
        \Big].
        \label{eq:clock_sector_action}
        \end{align}

        \textbf{[SPECULATIVE]} SST interpretation: Equation~\eqref{eq:clock_sector_action} is not itself derived
        from first-principles filament dynamics in the present Canon. It is recorded as the minimal orthodox EFT
        bridge into a relativistic clock-field description.

    \subsection{Poisson Limit and Effective Gravity Bridge}

        \textbf{[DERIVED]} Bridge statement: in the weak-field, slowly varying limit, the scalar clock field is taken to obey
        a Poisson-type closure
        \begin{align}
            \nabla^2 \chi(x)
            =
            4\pi G_{\mathrm{eff}} \,\rho_{\mathrm{matter}}(x),
            \label{eq:poisson_clock}
        \end{align}
        so that outside localized sources
        \begin{align}
            \chi(r) \propto \frac{1}{r},
            \qquad
            \mathbf{g}_{\mathrm{eff}} = - \nabla \chi.
            \label{eq:effective_gravity_clock}
        \end{align}

        \textbf{[ORTHODOX]} The \(1/r\) scalar potential and \(1/r^2\) force law are standard consequences of
        Eq.~\eqref{eq:poisson_clock}.

        \textbf{[SPECULATIVE]} The identification of \(\chi\) with the Swirl-Clock and the interpretation of
        Eq.~\eqref{eq:poisson_clock} as an emergent gravity law are SST-specific.

    \subsection{Observational Constraint from Tensor-Mode Speed}

        \textbf{[ORTHODOX] Constraint}: If the clock-sector EFT is used, multimessenger gravitational-wave
        observations require \cite{Monitor2017}
        \begin{align}
            c_{13} \equiv c_1 + c_3 = 0,
            \label{eq:c13_constraint}
        \end{align}
        to enforce luminal tensor-mode propagation at observational precision.

        \textbf{[ORTHODOX] Constraint}: any SST realization employing the EFT bridge of Eq.~\eqref{eq:clock_sector_action}
        must preserve Eq.~\eqref{eq:c13_constraint}.

    \subsection{Internal Tensor-Speed Naturalness}
    \label{subsec:tensor_speed_naturalness}

        \textbf{[DERIVED, conditional]} In the clock-field / Einstein--\AE ther
        comparison layer, the spin-2 tensor-mode speed is conventionally written as
        \begin{align}
            c_T^2
            =
            \frac{c^2}{1-c_{13}},
            \qquad
            c_{13}:=c_1+c_3 .
            \label{eq:eaether_tensor_speed}
        \end{align}
        Multimessenger observations require \(c_T=c\) to high precision and therefore
        impose \(c_{13}\simeq 0\) observationally \cite{Monitor2017,Abbott2017GW170817}. In SST there is a stronger conditional
        interpretation: if the metric tensor mode of the coarse-grained clock-field
        sector is identified with the transverse excitation mode of the underlying
        medium, and if \(c\) is the canonical transverse propagation speed of that
        medium, then
        \begin{align}
            c_T=c
            \quad\Longrightarrow\quad
            c_{13}=0 .
            \label{eq:c13_internal_identity}
        \end{align}

        \textbf{Status.} Equation~\eqref{eq:c13_internal_identity} is not an
        independent derivation of the Einstein--Hilbert dynamics. It is a conditional
        mode-identification result inside the orthodox Einstein--\AE ther comparison
        layer \cite{JacobsonMattingly2001,FosterJacobson2006}. Its canon role is to upgrade the luminal tensor-speed condition from a
        purely external observational constraint to an internal naturalness condition,
        provided the SST tensor-mode identification is valid.

        \textbf{[CRITICAL NOTE]} The proposition applies only to the spin-2/tensor
        sector. It does not by itself constrain the scalar and vector \AE{}-modes,
        nor does it prove monometricity for all matter, gauge, and knot species. The
        full monometricity target and the Einstein-dynamics routes are developed in the
        research track (Relativity Emergence Ladder).

    \subsection{Interpretive Role in the Canon}

        This section has a deliberately limited role:
        \begin{itemize}
            \item it upgrades the Swirl-Clock from a purely local kinematic factor to a candidate field variable;
            \item it records a minimal orthodox EFT bridge for relativistic comparison;
            \item it supplies the internal reference point needed by Section~\ref{sec:unification} for the time--gravity link.
        \end{itemize}

        \textbf{Status summary.}
        \begin{itemize}
            \item Eq.~\eqref{eq:clock_foliation_vector}: \textbf{[ORTHODOX]} mathematical form; \textbf{[DERIVED]} role as an SST bridge variable
            \item Eq.~\eqref{eq:clock_sector_action}: \textbf{[ORTHODOX]} EFT template
            \item Eq.~\eqref{eq:poisson_clock}: \textbf{[DERIVED]} bridge closure
            \item gravity interpretation: \textbf{[SPECULATIVE]}
        \end{itemize}

    \subsection{Relation to the Chronos--Kelvin law}

        Section~\ref{subsec:chronos_kelvin} provides the local transport identity of the Swirl-Clock on a material loop. The present section supplies the long-wavelength field-theoretic and analogue-metric reading of the same clock sector. The two layers are therefore complementary:

        \begin{itemize}
            \item Eq.~\eqref{eq:clock_radius_transport} is the loop-level transport law;
            \item Eq.~\eqref{eq:clock_foliation_vector}--Eq.~\eqref{eq:poisson_clock} provide the coarse-grained field description.
        \end{itemize}


    \subsection{Canonical Electromagnetic--Gravity Closure}
        \label{subsec:canonical_em_gravity_closure}

        \textbf{[CANONICAL BRIDGE LAW].} The electromagnetic and gravitational sectors of SST are canonically linked by a common local swirl state, but they remain distinct projections. Electromagnetic response is carried by transverse phase/flux variables and topological piercing changes. Gravitational response is carried by pressure gradients, swirl stress, and the local clock factor. Thus SST does not identify the ordinary electromagnetic field with gravity. It identifies both as different observable projections of the same incompressible, vorticity-bearing medium.

        The canonical source state for the coupled EM--gravity bridge is
        \begin{align}
            \boxed{
            \mathcal{X}^{\mathrm{EMG}}_{\swirlarrow}
            =
            \left(
            \mathbf{u}_{\swirlarrow},
            p_{\swirlarrow},
            \boldsymbol{\omega}_{\swirlarrow},
            \varrho_{\swirlarrow},
            \boldsymbol{\sigma}^{\mathrm{swirl}},
            \rhoE,
            \rhoM
            \right),
            \qquad
            \rhoM=\frac{\rhoE}{c^2}.
            }
            \label{eq:emg_source_state}
        \end{align}
        The topological piercing density \(\varrho_{\swirlarrow}\) belongs to the electromagnetic impulse channel, while \(p_{\swirlarrow}\), \(\rhoM\), and \(\boldsymbol{\sigma}^{\mathrm{swirl}}\) belong to the force and gravity channel.

        \textbf{[CANON / PREFACTOR GUARD]} The SI-normalized electromagnetic bridge, when invoked, uses the electron mass-to-charge factor
        \begin{align}
            \mathbf{A}_{\rm SI}
            =
            \frac{m_e}{e}\,\uswirl(\mathbf{x},t),
            \qquad
            \mathbf{B}_{\rm SST}
            =
            \nabla\times\mathbf{A}_{\rm SI}
            =
            \frac{m_e}{e}\,\boldsymbol{\omega}.
            \label{eq:canonical_em_prefactor_guard}
        \end{align}
        The inverse factor \(e/m_e\) is rejected in this bridge because it has the wrong SI dimensions for mapping vorticity to magnetic flux density. At the Compton angular frequency this calibrated bridge gives
        \begin{align}
            B_{\rm core}^{\rm SST}
            =
            \frac{m_e}{e}\omega_c
            =
            \frac{m_e^2c^2}{e\hbar}
            =
            B_{\rm QED},
            \label{eq:canonical_qed_field_guard}
        \end{align}
        while preserving the epistemic status of the result as a normalization bridge, not a derivation of QED from Euler dynamics.

        \paragraph{Electromagnetic projection.}
        The canonical topological electromagnetic impulse is
        \begin{align}
            \boxed{
            \Delta\Phi_{\swirlarrow}
            =
            \Phi_0\,\Delta N,
            \qquad
            \Phi_0=\frac{h}{2e},
            \qquad
            \Delta N\in\mathbb{Z}.
            }
            \label{eq:canonical_em_flux_impulse_closure}
        \end{align}
        In local differential notation, the canonical part of the effective Faraday projection is
        \begin{align}
            \boxed{
            \nabla\times\mathbf{E}_{\mathrm{eff}}
            =
            -\frac{\partial\mathbf{B}_{\mathrm{eff}}}{\partial t}
            -
            \Phi_0\frac{\partial\varrho_{\swirlarrow}}{\partial t}\,\hat{\mathbf n}
            -
            \mathbf{b}^{(\sigma)}_{\swirlarrow}.
            }
            \label{eq:canonical_effective_faraday_emg}
        \end{align}
        Here \(\mathbf{b}^{(\sigma)}_{\swirlarrow}\) is a calibrated stress-projection correction. It is not a primitive canonical constant. If no stress-projection closure has been derived for a given apparatus, set
        \begin{align}
            \mathbf{b}^{(\sigma)}_{\swirlarrow}=0.
            \label{eq:stress_projection_zero_default}
        \end{align}
        Equation~\eqref{eq:canonical_effective_faraday_emg} is an effective-medium projection. In the R-phase null limit,
        \begin{align}
            \partial_t\varrho_{\swirlarrow}\to0,
            \qquad
            \mathbf{b}^{(\sigma)}_{\swirlarrow}\to0,
        \end{align}
        and the ordinary source-free Maxwell--Faraday identity is recovered.

        \paragraph{Gravity projection.}
        The canonical SST gravitational force density is
        \begin{align}
            \boxed{
            \mathbf{f}_{\mathrm{grav},\swirlarrow}
            =
            \rhoM\mathbf{g}_{\swirlarrow}
            +
            \nabla\cdot\boldsymbol{\sigma}^{\mathrm{swirl}}.
            }
            \label{eq:canonical_sst_grav_force_density}
        \end{align}
        The bulk acceleration field is
        \begin{align}
            \boxed{
            \mathbf{g}_{\swirlarrow}
            =
            -\nabla\Phi_{\swirlarrow}
            -
            \frac{\partial\mathbf{A}_{\swirlarrow}}{\partial t},
            \qquad
            \mathbf{B}_{\swirlarrow}
            =
            \nabla\times\mathbf{A}_{\swirlarrow}.
            }
            \label{eq:canonical_gravito_swirl_fields}
        \end{align}
        In the weak, stationary, coarse-grained limit the scalar sector obeys the Poisson-type closure
        \begin{align}
            \boxed{
            \nabla\cdot\mathbf{g}_{\swirlarrow}
            =
            -4\pi G_{\mathrm{swirl}}\rhoM,
            \qquad
            \nabla^2\Phi_{\swirlarrow}
            =
            4\pi G_{\mathrm{swirl}}\rhoM.
            }
            \label{eq:canonical_weak_gravity_closure}
        \end{align}
        This closure is valid only after coarse graining over swirl-string microstructure and only in the weak-field, slow-flow regime where pressure reconstruction is consistent with Euler balance.

        The canonical pressure anchor remains the incompressible Euler radial balance
        \begin{align}
            \boxed{
            \frac{1}{\rhoF}\frac{dp_{\mathrm{swirl}}}{dr}
            =
            \frac{v_\theta(r)^2}{r}.
            }
            \label{eq:canonical_emg_euler_anchor}
        \end{align}
        Therefore the Newtonian-looking gravitational acceleration is not fundamental curvature in SST; it is the coarse-grained acceleration produced by pressure/stress gradients in the swirl medium.

        \paragraph{Clock and optical closure.}
        The clock sector supplies the shared time-rate projection:
        \begin{align}
            \boxed{
            \frac{d\tau_{\rm loc}}{dt_\infty}
            =
            \SwirlClock
            =
            \sqrt{1-\frac{\lVert\uswirl\rVert^2}{c^2}},
            \qquad
            n_\gamma=\SwirlClock^{-1}.
            }
            \label{eq:canonical_emg_clock_optical_closure}
        \end{align}
        Photon phase therefore probes \(n_\gamma\), while bulk matter probes Eq.~\eqref{eq:canonical_sst_grav_force_density}. These are coupled through \(\uswirl\), but they need not have identical spatial gradients.

        \paragraph{Electron-normalized coupling calibration.}
        The weak gravitational coupling is written in electron-normalized SST variables as
        \begin{align}
            \boxed{
            G_{\mathrm{swirl}}
            =
            \frac{\vchar c^3 t_p^2}{\rc M_e}
            =
            6.6743013\times10^{-11}\,\mathrm{m^3\,kg^{-1}\,s^{-2}}.
            }
            \label{eq:canonical_gswirl_emg_closure}
        \end{align}
        Thus the EM--gravity bridge is not an unlabeled free-parameter assertion. However, because \(t_p\) already contains \(G\) in orthodox units, Eq.~\eqref{eq:canonical_gswirl_emg_closure} is a \textbf{[CALIBRATED]} electron-normalized consistency identity unless a future SST sector derives \(t_p\) or \(G_{\mathrm{swirl}}\) independently from medium dynamics.

        \paragraph{Compton--horn reduction guard.}
        Substituting the calibrated horn closure \(\rc=\vchar/\omega_c\) into Eq.~\eqref{eq:canonical_gswirl_emg_closure} gives
        \begin{align}
            G_{\mathrm{swirl}}
            =
            \frac{\omega_c c^3 t_p^2}{M_e}.
            \label{eq:canonical_gswirl_compton_reduction}
        \end{align}
        With the orthodox definitions
        \begin{align}
            \omega_c=\frac{M_ec^2}{\hbar},
            \qquad
            t_p^2=\frac{\hbar G}{c^5},
        \end{align}
        Eq.~\eqref{eq:canonical_gswirl_compton_reduction} reduces identically to \(G\). This confirms that the electron-normalized bridge is algebraically closed under the v0.8.18 calibration chain, but it also fixes the epistemic status: the identity must not be cited as an independent derivation of Newton's constant unless \(t_p\), \(G_{\mathrm{swirl}}\), or the Compton--horn closure is derived without importing \(G\) through orthodox Planck units.

        \textbf{[CANON / STATUS GUARD]} A dimensionless balance angle \(\phi_G\) may be introduced only as a research-track diagnostic of Eq.~\eqref{eq:canonical_gswirl_emg_closure}. Since the calibrated chain gives \(\tan^2\phi_G=1\), the associated \(45^\circ\) angle is a visualization of equal radial-confinement and tangential-swirl weighting, not a physical direction of gravitational force and not an additional gravitational law. See Sec.~\ref{sec:rt_compton_horn_balance_angle}.

        \textbf{[ORTHODOX GR REPARAMETERIZATION].}
        In orthodox units the electron Schwarzschild radius may be written as
        \begin{align}
            r_s(e)
            =
            \frac{2Gm_e}{c^2}
            =
            \frac{4\pi \ell_P^2}{\lambda_c},
            \qquad
            \ell_P=\sqrt{\frac{\hbar G}{c^3}}.
            \label{eq:orthodox_rs_planck_reparam}
        \end{align}
        This is a standard GR identity and notational bridge. It does not derive \(G\) from SST medium dynamics.

        \paragraph{Dimensional checks.}
        The electromagnetic source term has units
        \begin{align}
            \left[\Phi_0\partial_t\varrho_{\swirlarrow}\right]
            =
            (\mathrm{V\,s})(\mathrm{m^{-2}\,s^{-1}})
            =
            \mathrm{V\,m^{-2}},
        \end{align}
        matching \([\nabla\times\mathbf{E}_{\mathrm{eff}}]\). The gravitational force density has units
        \begin{align}
            [\rhoM\mathbf{g}_{\swirlarrow}]
            =
            \mathrm{kg\,m^{-3}}\,\mathrm{m\,s^{-2}}
            =
            \mathrm{N\,m^{-3}},
            \qquad
            [\nabla\cdot\boldsymbol{\sigma}^{\mathrm{swirl}}]
            =
            \mathrm{N\,m^{-3}}.
        \end{align}
        Finally,
        \begin{align}
            \left[
            \frac{\lVert\vswirl\rVert c^3 t_p^2}{\rc M_e}
            \right]
            =
            \frac{(\mathrm{m\,s^{-1}})(\mathrm{m^3\,s^{-3}})(\mathrm{s^2})}
            {\mathrm{m}\,\mathrm{kg}}
            =
            \mathrm{m^3\,kg^{-1}\,s^{-2}}.
        \end{align}

        \paragraph{Numerical hierarchy.}
        At the canonical speed \(\vchar=1.09384563\times10^6\,\mathrm{m\,s^{-1}}\),
        \begin{align}
            \frac{1}{2}\rhoF\vchar{}^2
            &=
            4.1877439\times10^5\,\mathrm{Pa},
            \\
            \SwirlClock
            &=
            0.9999933436,
            \\
            n_\gamma-1
            &=
            6.6564858\times10^{-6}.
            \label{eq:canonical_emg_numerical_hierarchy}
        \end{align}
        Hence pressure/stress channels can be macroscopically large at canonical scale, while single-pass optical-index shifts are small unless a resonant optical path enhances them. This hierarchy is now part of the canonical interpretation of the EM--gravity bridge.

        \paragraph{Canonical exclusion rules.}
        The following statements are excluded from the canonical EM--gravity bridge:
        \begin{enumerate}
            \item Ordinary electromagnetic fields are not identified with gravity: \(\mathbf{B}_{\mathrm{eff}}\neq\mathbf{B}_{\swirlarrow}\).
            \item Uncalibrated vorticity terms may not be inserted directly into Maxwell source equations.
            \item Device-specific gains such as \(Q_\gamma\), \(\mathcal{G}^{(\sigma)}_{\swirlarrow}\), and impedance gains are research-track unless derived from a resonator or stress-closure model.
            \item Apparent laboratory forces must first be decomposed into electromagnetic, thermal, acoustic, electrohydrodynamic, mechanical, and ordinary vorticity-impulse channels before any residual is interpreted as SST-specific.
        \end{enumerate}

        \paragraph{Canonical summary.}
        The promoted law is therefore
        \begin{align}
            \boxed{
            \text{EM response and gravity response are distinct sectoral projections of }\mathcal{X}^{\mathrm{EMG}}_{\swirlarrow}.
            }
            \label{eq:canonical_emg_summary_law}
        \end{align}
        The electromagnetic projection is controlled by topological flux impulse \(\Delta\Phi_{\swirlarrow}=\Phi_0\Delta N\). The gravitational projection is controlled by \(\rhoM\mathbf{g}_{\swirlarrow}+\nabla\cdot\boldsymbol{\sigma}^{\mathrm{swirl}}\) with weak closure \(\nabla\cdot\mathbf{g}_{\swirlarrow}=-4\pi G_{\mathrm{swirl}}\rhoM\). Both recover their orthodox limits when the SST projection channels vanish.


    \subsection{Triadic Gravity-Response Corollary}
        \label{subsec:triadic_gravity_response_corollary}

        \textbf{[DERIVED] Canonical corollary.}
        Because the electromagnetic, gravitational, optical, and boundary-stress
        responses are distinct projections of the same SST state
        \(\mathcal{X}^{\mathrm{EMG}}_{\swirlarrow}\), a laboratory observation must not be
        classified as ``gravity'' from force response alone. The admissible coarse
        response separates into three limiting channels:
        \begin{align}
            \mathcal{R}_{\mathrm{plume}}
            &:=
            \left(
            \Gamma_{\swirlarrow}\neq 0,
            \quad
            \nabla n_\gamma \approx 0,
            \quad
            \mathbf{g}_{\swirlarrow}\approx 0
            \right),
            \label{eq:triadic_plume_mode}\\
            \mathcal{R}_{\gamma}
            &:=
            \left(
            \nabla_\perp n_\gamma \neq 0,
            \quad
            \Gamma_{\swirlarrow}\approx 0,
            \quad
            \mathbf{f}_{\mathrm{matter}}\approx 0
            \right),
            \label{eq:triadic_photon_mode}\\
            \mathcal{R}_{\Pi}
            &:=
            \left(
            \nabla\cdot\boldsymbol{\sigma}^{\mathrm{swirl}}\neq0,
            \quad
            \nabla n_\gamma\approx0
            \right).
            \label{eq:triadic_boundary_mode}
        \end{align}
        Here \(\mathcal{R}_{\mathrm{plume}}\) is an impedance or damping response of a
        plume or flame-like medium, \(\mathcal{R}_{\gamma}\) is an optical phase/caustic
        response, and \(\mathcal{R}_{\Pi}\) is a boundary-stress or elastic-shell response.
        The three modes may share a common source state but are not mutually identical.

        \paragraph{Local Euler locking of pressure and optical projections.}
        Under the additional assumptions of a passive, stationary, inviscid Euler
        region with constant \(\rhoF\), the pressure and optical projections are not
        independent.  Bernoulli reconstruction gives
        \begin{align}
            \delta p_{\mathrm{swirl}}
            =
            -\frac{1}{2}\rhoF\lVert\uswirl\rVert^2,
        \end{align}
        while the clock/optical closure gives
        \begin{align}
            n_\gamma
            =
            \left(
                1-\frac{\lVert\uswirl\rVert^2}{c^2}
            \right)^{-1/2}
            =
            1+\frac{\lVert\uswirl\rVert^2}{2c^2}
            +O\!\left(\frac{\lVert\uswirl\rVert^4}{c^4}\right).
        \end{align}
        Eliminating \(\lVert\uswirl\rVert^2\) yields the parameter-free leading-order
        locking relation
        \begin{align}
            \boxed{
            \delta n_\gamma
            =
            -\frac{\delta p_{\mathrm{swirl}}}{\rhoF c^2}
            }
            \label{eq:canonical_pressure_optical_locking}
        \end{align}
        with validity restricted to the passive Euler regime.  Forcing, damping,
        viscosity, thermal gradients, acoustic drive, plasma response, or explicit
        Mode-I terms such as \(\Gamma_{\swirlarrow}\mathbf{u}_{\mathrm{plume}}\)
        break the assumptions and move the prediction to the research-track transfer
        model.  Thus the refined diagnostic statement is not three fully independent
        channels, but two locally locked projections plus one nonlocal Poisson/bulk
        response whenever Eq.~\eqref{eq:canonical_weak_gravity_closure} is invoked.

        With \(\rhoF=7.0\times10^{-7}\,\mathrm{kg\,m^{-3}}\),
        \begin{align}
            \frac{1}{\rhoF c^2}
            &=
            1.58950008\times10^{-11}\,\mathrm{Pa^{-1}},
            \\
            \frac{1}{2}\rhoF\vchar{}^2
            &=
            4.1877439\times10^5\,\mathrm{Pa},
            \\
            \delta n_\gamma
            &=
            6.6565\times10^{-6}
        \end{align}
        at canonical swirl speed, consistent with Eq.~\eqref{eq:canonical_emg_numerical_hierarchy}.

        Therefore
        \begin{equation}
        \boxed{
        \text{matter stabilization}
        \neq
        \text{photon caustic}
        \neq
        \text{boundary shell force}.
        }
        \end{equation}

        \textbf{[CRITICAL NOTE].}
        This corollary canonizes only the diagnostic separation of response channels.
        It does not canonize any specific witness claim, device claim, or anti-gravity
        interpretation. Device-specific excitation factors, optical gains, and shell
        stiffness parameters remain research-track unless derived from an explicit
        resonator, stress-closure, or calibrated transfer model.



%    \section{Unified Interpretation}
%        \label{sec:unification}
        
\section{Unified Interpretation}
    \label{sec:unification}
    
    \subsection{Overview}
        
        Swirl-String Theory (SST) proposes a unified interpretation in which fundamental physical phenomena emerge from the dynamics of a continuous medium supporting stable vortex structures.
        
        This section synthesizes the results of the previous sections into a coherent conceptual framework.
    
    \subsection{Ontology of the Theory}
        
        \textbf{[SPECULATIVE] Ontological claim}:
        The fundamental ontology consists of:
        
        \begin{itemize}
            \item A continuous incompressible medium (Section~\ref{sec:axioms})
            \item Stable vortex filament configurations (Section~\ref{sec:geometry})
            \item Circulation as a conserved topological invariant
        \end{itemize}
        
        \textbf{Interpretation:}
        Particles are not elementary objects, but persistent dynamical structures.
    
    \subsection{Mass as Trapped Circulation Energy}
        
        From Section~\ref{sec:foundations}, and in particular from Eq.~\eqref{eq:sst_master_equation}, the mass sector is organized as the product of a medium scale, a geometric gate, a topological kernel, and a local clock-impedance factor. The canon-level reading is therefore not merely that ``energy is stored in rotation,'' but that persistent inertial scale is assigned only after geometry, topology, and local clocking have all been specified.
        
        \begin{align}
            M_K(x)=\rhoM(x)\,\rc^3\,\Pi_K\,\Xi_K\,\SwirlClock(x)^{-2}.
        \end{align}
        
        \textbf{[SPECULATIVE] Unified interpretation:} mass corresponds to localized, self-sustained circulation energy in the medium, modulated by state geometry, topology, and local clock impedance.
    
    \subsection{Time as a Relational Quantity}
        
        From Section~\ref{sec:time} and Section~\ref{sec:foundations}, time is not introduced as an independent primitive but as a relational field read from the local state of the medium. The same kinematic quantity that enters the Swirl-Clock also appears in the organizing equation of the mass sector, so the canon treats proper-time modulation and inertial scaling as two projections of one shared substrate state.
        
        \begin{align}
            \SwirlClock(x,t)=\sqrt{1-\frac{\norm{\uswirl(x,t)}^2}{c^2}}.
        \end{align}
        
        \textbf{[SPECULATIVE] Unified interpretation:} time is not fundamental, but emerges from the kinematic state of the medium and is read operationally through clock-bearing processes.
    
    \subsection{Discreteness without Quantization Postulates}
        
        From Section~\ref{sec:delay}:
        
        \begin{itemize}
            \item Delay introduces temporal nonlocality
            \item Stability selects discrete modes
        \end{itemize}
        
        \textbf{[DERIVED] Unified consequence}:
        \begin{center}
            Discrete physical spectra arise dynamically from delay-induced mode selection.
        \end{center}
        
        This replaces operator-based quantization with a dynamical mechanism.
    
    \subsection{Charge and Circulation}
        
        \textbf{[SPECULATIVE] Provisional bridge only.} The canon retains circulation as a structurally important invariant, but it does \emph{not} yet claim a completed derivation of observed gauge charges from circulation data alone. At most, the present document records a working hypothesis that orientation and magnitude of circulation may enter a later charge-assignment program.
    
    \subsection{Unification of Interactions}
        
        The canon records a \emph{unification program}, not yet a completed unification theorem. In the present document:
        \begin{itemize}
            \item electromagnetic effects are represented by the minimal transverse-field bridge of Section~\ref{sec:embridge},
            \item gravitational effects are represented by the clock-field and Poisson-limit bridge of Section~\ref{sec:time},
            \item inertial mass is organized by the medium/topology/geometry/clock structure of Eq.~\eqref{eq:sst_master_equation}.
        \end{itemize}
        
        \textbf{[SPECULATIVE] Programmatic claim:} these sectors may ultimately be read as different effective descriptions of one substrate, but their full common derivation remains open.
    
    \subsection{Atomic Structure as an Emergent Layer}
        
        From Section~\ref{sec:atomic}:
        
        \begin{itemize}
            \item Standard quantum mechanics appears as an effective description
            \item SST introduces circulation-dependent corrections
        \end{itemize}
        
        \textbf{[SPECULATIVE] Unified interpretation}:
        \begin{center}
            Quantum mechanics is an emergent effective theory of underlying vortex dynamics.
        \end{center}
    
    \subsection{Consistency with Spectroscopy}
        
        From Section~\ref{sec:spectroscopy}:
        
        \begin{itemize}
            \item Internal modes are suppressed by an excitation gap
            \item Low-energy physics remains unchanged
        \end{itemize}
        
        \textbf{[ORTHODOX] Constraint}:
        The unified picture does not violate known experimental constraints.
    
    \subsection{Hierarchy of Description}
        
        The theory admits multiple levels:
        
        \begin{enumerate}
            \item \textbf{Microscopic level:} vortex filament dynamics
            \item \textbf{Mesoscopic level:} delay-induced mode structure
            \item \textbf{Macroscopic level:} effective quantum description
        \end{enumerate}
    
    \subsection{Conceptual Summary}
        
        The SST framework can be summarized as:
        
        \begin{center}
            \begin{tabular}{ll}
                \textbf{Mass} & $\rightarrow$ trapped circulation energy \\
                \textbf{Time} & $\rightarrow$ relational swirl-clock field \\
                \textbf{Charge} & $\rightarrow$ provisional circulation bridge \\
                \textbf{Quantization} & $\rightarrow$ delay-induced stability \\
            \end{tabular}
        \end{center}
    
    \subsection{Canonical Statement}
        
        \begin{center}
            \textit{
                All observed physical structure emerges from the dynamics of a continuous medium in which stable vortex configurations, governed by quantized circulation and delay dynamics, produce the phenomena conventionally attributed to particles, fields, and quantum mechanics.
            }
        \end{center}
    
    \subsection{Scope and Limits}
        
        \textbf{[IMPORTANT]}:
        
        \begin{itemize}
            \item The framework is not yet a complete fundamental theory
            \item Several mechanisms remain to be derived rigorously
            \item Many results are conditional or phenomenological
        \end{itemize}
    
    \subsection{Final Consistency Principle}
        
        \begin{center}
            \textit{
                A valid unified interpretation must preserve all experimentally verified results while providing a deeper structural explanation of their origin.
            }
        \end{center}

    \section{Integration of Prior Canon Versions}
        \label{sec:integration}

    \subsection{Integration policy (v0.5.12 reintegration)}

        This block records only material from earlier SST canon layers that can be retained inside the v0.8.1 architecture without weakening the explicit classification discipline of Section~\ref{sec:consistency}.

        Older canon documents are therefore \emph{not} cited as authorities. Transferable content is rewritten here as current-canon statements, each labeled according to its present status.

    \subsection{Retained v0.5.12 kernel sectors}

        The present draft now retains the following v0.5.12 kernel sectors in rewritten form:
        \begin{itemize}
            \item the Chronos--Kelvin invariant and clock--radius transport law (Section~\ref{subsec:chronos_kelvin});
            \item the reparametrized zero-parameter corollary and coarse-graining rule (Section~\ref{subsec:zero_parameter_corollary});
            \item a minimal Swirl--EM bridge and photon / unknot sector (Section~\ref{sec:embridge});
            \item the swirl pressure law as the Euler-level force bridge (Section~\ref{subsec:swirl_pressure_law});
            \item the analogue-metric / clock-field bridge (Section~\ref{sec:time}).
        \end{itemize}

    \subsection{Wave--particle phase rule and measurement bridge}
        \label{subsec:measurement_bridge}

        A retained qualitative rule from earlier canon layers is that SST distinguishes two limiting phases:
        \begin{align}
            \text{R-phase} &:\quad \text{extended, unknotted, radiative}, \\
            \text{T-phase} &:\quad \text{localized, knotted, tangible}.
        \end{align}

        A minimal dynamical bookkeeping law for phase conversion is

        \begin{align}
            \partial_t P_T
            =
            \Gamma_{R\to T}[\mathcal{K}]\,P_R
            -
            \Gamma_{T\to R}[\mathcal{K}]\,P_T,
            \qquad
            P_R + P_T = 1,
            \label{eq:rt_phase_rate_equation}
        \end{align}

        where $\mathcal{K}$ denotes an environment- and coherence-dependent kernel functional. The intended interpretation is that ``measurement'' corresponds not to an axiomatically imposed collapse, but to a medium-induced transfer of support from the extended R-phase to the localized T-phase.

        \textbf{[ORTHODOX LIMIT]}: in weak coupling and after coarse graining, Eq.~\eqref{eq:rt_phase_rate_equation} should reduce to ordinary environment-induced decoherence models \cite{Zurek2003}.

        \textbf{[SPECULATIVE]}: the identification of decoherence with an actual R$\leftrightarrow$T phase transition is SST-specific.

    \subsection{Gauge and weak-mixing bridge}
        \label{subsec:gauge_weakmix_bridge}

        The retained gauge-level statement is that the effective low-energy algebra should reduce to

        \begin{align}
            \mathfrak{g}_{\mathrm{eff}} = \mathfrak{su}(3)\oplus\mathfrak{su}(2)\oplus\mathfrak{u}(1),
            \label{eq:effective_gauge_algebra}
        \end{align}

        with gauge potentials interpreted as holonomy data of coarse-grained multi-director transport.

        A minimal bridge ansatz for the weak mixing angle is

        \begin{align}
            \sin^2\theta_W^{\mathrm{eff}} = \frac{K_Y}{K_Y + K_W},
            \label{eq:weak_mixing_bridge}
        \end{align}

        where $K_Y$ and $K_W$ are effective Abelian and weak-sector stiffness scales. Eq.~\eqref{eq:weak_mixing_bridge} is not a completed derivation, but it preserves the v0.5.12 idea that weak mixing should descend from a ratio of medium-response coefficients rather than appear as an unstructured free number.

        \textbf{[ORTHODOX]}: $\theta_W$ is the standard electroweak mixing parameter of the SM \cite{Weinberg1967,PDG2024}.

        \textbf{[SPECULATIVE]}: Eq.~\eqref{eq:weak_mixing_bridge} is an SST bridge ansatz only.

    \subsection{What the v0.5.12 reintegration does not import}

        The present reintegration does \emph{not} import v0.5.12 cosmogony, engineering subsystems, cosmological-term closures, or highly application-specific protocols into the main canon body. Those remain suitable for appendices or standalone papers until they are more tightly tied to the minimal axiom system.

    \subsection{Purpose of this integration layer}
        Version~0.8.1 is not a fresh theory written from zero, but an explicit consolidation of earlier canon lines. This section records what has been retained, what has been renamed or reorganized, and what has been demoted from canon status to research-track status.

    \subsection{Retained from the early canon layers}
        From the v0.1.x--v0.3.x line, the present canon retains the primitive closure logic:
        \begin{itemize}
            \item a continuum substrate with effective density $\rhoF$,
            \item a core scale $r_c$ fixed by Compton anchoring,
            \item circulation quantization through $\Gamma_0 = 2\pi \rc \, \vchar$,
            \item the interpretation of matter as stable organized structures rather than point primitives.
        \end{itemize}
        These elements remain non-negotiable because later sectors depend on them directly.

    \subsection{Retained from the middle canon layers}
        From the v0.4.x--v0.6.x line, the canon retains three durable structures:
        \begin{enumerate}
            \item the Rosetta strategy for mapping SST statements into orthodox language,
            \item the photon/torsion sector as the minimal radiation bridge,
            \item delay-induced mode selection as the main route to discreteness without operator postulates.
        \end{enumerate}
        These are now distributed across Sections~\ref{sec:delay}, \ref{sec:time}, and \ref{sec:unification} rather than treated as isolated notes.

    \subsection{Retained from the late canon layers}
        From the v0.7.x line, v0.8.1 retains the strongest structurally reusable sectors:
        \begin{itemize}
            \item the Swirl-Clock as the canonical relational-time field,
            \item the atomic bridge as a benchmark consistency layer rather than a complete atomic derivation,
            \item the spectroscopic gap logic as a hard consistency filter,
            \item the topology-first mass program, with topology separated from thermodynamic or radiative dressing,
            \item the EFT-style foliation and gauge-emergence program as a long-range development track.
        \end{itemize}

    \subsection{Notation and density cleanup}
        Earlier canon versions sometimes conflated background density, mass-equivalent density, and saturated envelope-equivalent density. The present version adopts the cleaned hierarchy
        \begin{align}
            \rhoF
            &\quad \text{effective fluid density},
            \nonumber\\
            \rhoE
            &\quad \text{swirl energy density},
            \nonumber\\
            \rhoM = \rhoE/c^2
            &\quad \text{mass-equivalent density},
            \nonumber\\
            \rhohorn
            &\quad \text{canonical saturated envelope-equivalent density},
            \nonumber\\
            \rhocore
            &\quad \text{legacy alias, deprecated unless explicitly defined}.
            \label{eq:density_symbol_roles}
        \end{align}
        Any older passage that identifies a single symbol with more than one of these roles is superseded by the present separation.

    \subsection{Topology-first reinterpretation of layers}
        A major consolidation from the recent lepton and knot-mass program is the separation between topology and layering. Knot type is retained as the primary identity label, while Golden-layer or thermal ladder indices are treated as dressing or coherence variables rather than species-defining labels. This prevents the canon from over-reading numerical layer fits as ontological identities.

    \subsection{Status of optional sectors}
        The following sectors are retained as important but non-primitive research layers:
        \begin{itemize}
            \item full gauge emergence and charge assignments,
            \item dual-vacuum phenomenology and Unruh-echo interpretation,
            \item complete many-electron atomic closure,
            \item full weak-sector and collider phenomenology,
            \item explicit topology-to-particle dictionary beyond the best-supported benchmark assignments.
        \end{itemize}
        They remain in the canon as organized programs, not as closed theorem-level results.

    \subsection{Integration principle}
        The purpose of v0.8.1 is therefore not to maximize novelty per page, but to minimize duplication and to expose the dependency graph of the theory. Whenever two earlier canon fragments make the same structural claim in different language, the present version keeps one canonical form and demotes the others to historical context.

    \subsection{Research-track companion}
        The sectors that are useful for continuity but too long or too speculative for a minimal main-canon reading are also collected in the companion file \texttt{canon-0.8.1-research-track.tex}. That excerpt includes the thermodynamic-radiation interface, the minimal-coupling QED analogy, the ribbon-invariant chirality refinement, the particle-candidate mapping layer, the hydrodynamic exchange / Pauli-barrier estimate, the numerical-benchmark policy block, the gauge-sector roadmap, the four topology/stability appendices on non-Abelian protection, reconnection decay, fractional torus-phase winding, and topological-fluid mechanics, and the \emph{dark-sector and galactic canonization execution package} (frozen topology toolkit, benchmark protocol template, Euler-pressure anchor details, galactic branch gate, and probe--equation map). The integration subsections below retain the full in-context treatments with labels and cross-references; the companion file provides a portable standalone copy. Editorial separation does not reject these sectors; it prevents the core narrative from over-claiming closure.

    \subsection{Research-track knot-energy normal form}
        \label{sec:research_track_knot_energy_normal_form}

        \paragraph{Status.}
        \textbf{[RESEARCH TRACK / NON-DERIVED SCREENING FUNCTIONAL].}
        The main canon does not assume a closed universal knot-energy functional beyond the calibrated constant chain and the explicitly derived hydrodynamic energy scales.  For classification and numerical screening, the research track uses the dimensionless normal form
        \begin{equation}
            \frac{E_{\rm eff}[K]}{E_0}
            =
            \alpha_C C(K)
            +\alpha_L L(K)
            +\alpha_H\widehat{\mathcal H}(K)
            +\alpha_G I_G(K)
            +\cdots,
            \label{eq:research_track_energy_normal_form}
        \end{equation}
        where $C(K)$ denotes a crossing/contact-complexity proxy, $L(K)$ a ropelength or relaxed-length invariant, $\widehat{\mathcal H}(K)$ a nondimensional helicity measure, and $I_G(K)$ a scalar index extracted from a possible non-Abelian topological charge sector $Q_G(K)$.  The map
        \begin{equation}
            I_G:Q_G(K)\longrightarrow \mathbb R_{\ge 0}
        \end{equation}
        is required because $Q_G(K)$ is generally a sector label, not a real-valued energy contribution.  Equation~\eqref{eq:research_track_energy_normal_form} is therefore a research-track screening functional, not a derived SST Hamiltonian or mass theorem.

    \subsection{Dark-sector and galactic canonization status (v0.8.x)}
        \label{sec:dark-canonicalization-status}

        This block records promotion labels for the dark-sector / galactic program consistent with the internal canonization plan. Detailed freezes, the mandatory benchmark-protocol checklist, the galactic profile branch gate (Branch~A: legacy-profile refine with benchmark traceability), and the probe falsification table live in \texttt{canon-0.8.1-research-track.tex} (Section titled \texttt{Dark-sector and galactic canonization execution package}).

        \begin{table}[h]
        \centering
        \caption{Master tracker: dark-sector and galactic program (v0.8.x).}
        \label{tab:dark-sector-master-tracker}
        \begin{tabular}{p{4.2cm}p{2.3cm}p{2.5cm}p{5.5cm}}
        \hline
        Item & Prior state & Target & Status after Stage 0--4 freeze \\
        \hline
        Dark-sector Euler pressure anchor & BRIDGE & CANON & CANON-CANDIDATE (derivation, assumptions, limits, falsifiers frozen; see Section~\ref{subsec:swirl_pressure_law}) \\
        Dark taxonomy & HOLD & CANON & BRIDGE (symmetry-first classifier + filters; Section~\ref{subsec:dark_knot_classification}) \\
        Galactic swirl law & BRIDGE & CANON & BRIDGE (Branch~A closure; companion execution package) \\
        Experimental probes & HOLD & CANON-BRIDGE & CANON-BRIDGE (equation / null-result map; companion) \\
        Benchmark / reproducibility & BRIDGE & CANON & CANON-CANDIDATE (protocol template frozen; repository-wide reuse pending) \\
        Topological reference identities & CANON-CANDIDATE & CANON & CANON-ready scope/notation (Appendix~\ref{sec:appendix_topology_toolkit}) \\
        \hline
        \end{tabular}
        \end{table}

        \textbf{[DERIVED] Stage 5 review:} cross-item dependencies are respected (pressure anchor and benchmark layer precede full galactic closure; taxonomy strengthens other sectors). No sector is promoted on fit quality alone; each carries explicit falsifiability or reproducibility hooks as in Table~\ref{tab:dark-sector-master-tracker} and the companion execution package.

    \subsection{Integration Policy}

        This section records only material from prior internal SST development that can be retained inside the
        v0.8.1 architecture \emph{without} weakening the epistemic classification framework of
        Section~\ref{sec:consistency}.

        \textbf{Rule.}
        Older canon documents are \emph{not} cited as external authorities. Instead, transferable content is
        rewritten here as current-canon statements, each labeled as \textbf{[ORTHODOX]}, \textbf{[DERIVED]},
        or \textbf{[SPECULATIVE]}.

    \subsection{Selective Recovery from Pre-v0.7 Canon Layers}

        The present draft absorbs only those legacy ingredients that can be rewritten in the current epistemic format without appealing to earlier canon versions as authorities.

        \paragraph{Analogue line-element bridge.}
        Legacy SST drafts recorded a schematic analogue line element and a metric--clock identification of the form
        \begin{align}
            ds^2 &= -c^2\,\SwirlClock(x)^2\,dt^2 + \delta_{ij}\,dx^i\,dx^j,
            \label{eq:legacy_metric}\\
            g_{00}(x) &= -\SwirlClock(x)^2 .
            \label{eq:metric_swirlclock_bridge}
        \end{align}

        Recovered legacy sectors include:
        \begin{itemize}
            \item the finite-core atomic bridge, Eqs.~\eqref{eq:lambda_sst}--\eqref{eq:softcore_potential};
            \item the analogue line-element bridge, Eqs.~\eqref{eq:legacy_metric}--\eqref{eq:metric_swirlclock_bridge};
            \item compact research-track interfaces to thermodynamic radiation and minimal-coupling QED analogies;
            \item topological chirality refinements based on standard ribbon invariants.
        \end{itemize}

        Legacy material not imported includes version history, cosmogonic narrative, and phenomenological galaxy-scale laws that are not yet sufficiently integrated with the present minimal axiom system.

    \subsection{Thermodynamic Radiation Interface}

        A recurring legacy idea is that a local swirl-energy density may define an effective temperature field and thereby an emergent radiation sector.

        Let
        \begin{align}
            T_{\mathrm{eff}} = T_{\mathrm{eff}}\!\left(\rho_E,\,\omega,\,S_{(t)}\right)
            \label{eq:teff_placeholder}
        \end{align}
        denote an unspecified constitutive relation. A minimal legacy interface then asks whether Planck-type spectra can emerge from a medium whose local excitation scale is set by $T_{\mathrm{eff}}$.

        \textbf{[ORTHODOX]}: blackbody spectra are governed by Planck's law once a valid temperature variable is defined.\newline
        \textbf{[SPECULATIVE]}: SST does not yet derive Eq.~\eqref{eq:teff_placeholder}, so this sector remains a research program rather than canonically closed physics.

    \subsection{Minimal Coupling Interface to Quantum Electrodynamics}

        Early research notes repeatedly explored whether an effective vector potential could be built from circulation variables. The most conservative placeholder takes the form
        \begin{align}
            \nabla \rightarrow \nabla - i\,\frac{m}{\hbar}\,\mathbf{A}_{\mathrm{eff}},
            \qquad
            \mathbf{A}_{\mathrm{eff}} = \chi\,\mathbf{v},
            \label{eq:legacy_minimal_coupling}
        \end{align}
        where $\chi$ is a scalar weight and $\mathbf{v}$ is the local carrier velocity field.

        \textbf{[ORTHODOX]}: minimal coupling is the standard gateway from free to gauge-coupled wave dynamics.\newline
        \textbf{[SPECULATIVE]}: Eq.~\eqref{eq:legacy_minimal_coupling} is only an analogy template unless a full gauge-covariant derivation is supplied.

    \subsection{Ribbon-Invariant Chirality Refinement}

        The older canon layers also suggested using the C\u{a}lug\u{a}reanu--White--Fuller relation
        \begin{align}
            Lk = Tw + Wr
            \label{eq:calugareanu}
        \end{align}
        as a refinement of the knot-based matter/antimatter dictionary.

        \textbf{[ORTHODOX]}: Eq.~\eqref{eq:calugareanu} is a standard geometric identity for framed curves and ribbons.\newline
        \textbf{[DERIVED BOOKKEEPING]}: once a framed curve is specified, the identity separates writhe and twist contributions.\newline
        \textbf{[SPECULATIVE]}: the claim that writhe-dominated and twist-dominated realizations are dynamically stable physical isomers, or map to particle/antiparticle sectors, remains a model hypothesis.


    \subsection{Framed-Tube Ontology and the Attachment Gate}
        \label{subsec:framed_tube_attachment_gate}

        The particle carrier in SST is not a bare centerline knot alone. The canonical object is a framed, phase-bearing tube,
        \begin{align}
            \mathcal{K}_{\mathrm{SST}}
            =
            \big(K,\,\mathcal{F},\,\Gamma,\,\theta_{\mathrm{core}}\big),
        \end{align}
        where $K$ is the centerline knot or link, $\mathcal{F}$ is a ribbon/framing, $\Gamma$ is circulation, and $\theta_{\mathrm{core}}$ is a material or order-parameter phase attached to the core.

        The current Attachment Lemma is the statement that a local Compton/core cycle can be promoted to a globally closed framed-tube holonomy,
        \begin{align}
            \frac{1}{2\pi}\oint_K d\theta_{\mathrm{core}}
            = \chi,
            \qquad
            \chi\in\mathbb{Z}.
            \label{eq:attachment_gate_chi}
        \end{align}
        A compensated form allows local swirl-clock redistribution while preserving the exact global integer,
        \begin{align}
            \chi_{\mathrm{total}}
            =
            \chi_{\mathrm{local}}+\chi_{\mathrm{attachment}}
            = \chi,
            \qquad
            \delta Q=0.
            \label{eq:compensated_attachment_gate}
        \end{align}

        \textbf{[CANON / OPEN GATE]}: Eq.~\eqref{eq:attachment_gate_chi} is not yet a theorem of the Euler--SST dynamics. It is a load-bearing gate. Integer energy selection alone is insufficient; exact holonomy is required for canon status.

        \textbf{[CANON / NOTATION]}: the framing phase follows the Compton angular cycle $\omega_c$, while the associated local vorticity scale is $2\omega_c$. The factor of two must not be absorbed into $\chi$.


    \subsection{Particle-Candidate Mapping Layer}
        \label{subsec:integration_knotmap}

        \textbf{[SPECULATIVE]} Working dictionary: a minimal knot/composition mapping retained from earlier SST work is summarized in Table~\ref{tab:particle_candidate_mapping}. This table is organizational rather than definitive; it records current topological assignments and mechanical analogues used by the mass and taxonomy program.

        \begin{table}[h]
        \centering
        \caption{Minimal knot-class and analogue summary retained in v0.8.3, with current semantic corrections.}
        \label{tab:particle_candidate_mapping}
        \begin{tabularx}{\linewidth}{p{0.28\linewidth}p{0.22\linewidth}X}
        \hline
        Carrier / composition & Candidate role & Status / interpretation \\
        \hline
        $3_1$ trefoil & electron / positron sector & baseline chiral lepton calibration; not used for fractional quark charge \\
        $T(2,5)$ or equivalent higher torus class & muon & lepton hierarchy candidate \\
        $T(2,7)$ or equivalent higher torus class & tau & lepton hierarchy candidate \\
        twist/clasp knots, quasi-unknot loops & quark-like defects & research-track fractional boundary winding candidates \\
        $5_2+5_2+6_1$ and nearby composites & baryon candidates & legacy mass-functional dictionary, still under audit \\
        triple-gear locked analogue & proton-like charged baryon analogue & mechanical model for nonzero exterior collective rotation \\
        Borromean-type analogue & neutron-like neutral baryon analogue & mechanical/topological model for zero exterior charge with internal triple linking \\
        \hline
        \end{tabularx}
        \end{table}

        \textbf{[CRITICAL NOTE]}: Table~\ref{tab:particle_candidate_mapping} is a working dictionary only. Electric charge is not identified with raw Hopf helicity or pairwise linking number. The safer current placeholder is a boundary-winding charge variable $q_{\partial}$, while knot/link data supply mass, chirality, confinement, and internal topological structure.

        \textbf{[CANON / DECISION]}: The trefoil assignment is retained for the electron/positron sector only. Fractional values such as $2/3$ belong, if anywhere, to the twist-knot/quark-like boundary-winding research track rather than to the trefoil-electron calibration.

    \subsection{Hydrodynamic Exchange and the Pauli Barrier}
        \label{subsec:integration_pauli}

        \textbf{[ORTHODOX]} Fluid-mechanical input: for two thin filamentary loops \(C_1\) and \(C_2\) in an
        incompressible medium, the Biot--Savart interaction energy is \cite{Batchelor1967}
        \begin{align}
            E_{\mathrm{int}}
            =
            \frac{\rho_{\!f}\Gamma_1\Gamma_2}{4\pi}
            \oint_{C_1}\oint_{C_2}
            \frac{d\mathbf{l}_1\cdot d\mathbf{l}_2}
            {\lVert \mathbf{r}_1 - \mathbf{r}_2 \rVert}.
            \label{eq:biot_savart_interaction}
        \end{align}

        \textbf{[ORTHODOX]} Regularized filament-energy template: if two identical electron-candidate loops are forced into the same spatial state, a short-distance cutoff \(a_{\rm cut}\) gives the overlap penalty
        \begin{align}
            V_{\mathrm{Pauli}}(\mathbf{d};a_{\rm cut})
            \approx
            \frac{\rho_{\!f}\Gamma_0^2}{4\pi}
            \oint\!\!\oint
            \frac{d\mathbf{l}_1\cdot d\mathbf{l}_2}
            {\sqrt{\lVert \mathbf{r}_1(\theta_1)-\mathbf{r}_2(\theta_2)+\mathbf{d}\rVert^2+a_{\rm cut}^2}},
            \label{eq:pauli_barrier_regularized}
        \end{align}
        with maximal-overlap scale
        \begin{align}
            V_{\mathrm{Pauli}}^{\max}(a_{\rm cut})
            \approx
            \frac{\rho_{\!f}\Gamma_0^2}{4\pi}\left(\frac{L}{a_{\rm cut}}\right)\mathcal{O}(1).
            \label{eq:pauli_barrier_scale}
        \end{align}

        Using the canonical values
        \begin{align}
            \rho_{\!f} &= 7.0 \times 10^{-7}\ \mathrm{kg\,m^{-3}}, \nonumber\\
            r_c &= 1.40897017 \times 10^{-15}\ \mathrm{m}, \nonumber\\
            \Gamma_0 &= 2\pi r_c \vchar
            \approx 9.6836\times 10^{-9}\ \mathrm{m^2\,s^{-1}},
        \end{align}
        and \(L \sim 2\pi a_0\) at the hydrogenic scale, the legacy benchmark corresponds to the cutoff calibration \(a_{\rm cut}=r_c\):
        \begin{align}
            V_{\mathrm{Pauli}}^{\max}(a_{\rm cut}=r_c)
            \sim 1.23\times 10^{-18}\ \mathrm{J}
            \sim 7.6\ \mathrm{eV}.
            \label{eq:pauli_barrier_numeric}
        \end{align}
        Away from this cutoff calibration the scaling is
        \begin{align}
            V_{\mathrm{Pauli}}^{\max}(a_{\rm cut})
            \sim
            V_{\mathrm{Pauli}}^{\max}(r_c)
            \left(\frac{r_c}{a_{\rm cut}}\right).
            \label{eq:pauli_barrier_acut_scaling}
        \end{align}

        \textbf{[CRITICAL NOTE]} The cutoff \(a_{\rm cut}\) is not the resolved physical tube radius \(a_{\rm core}\).  The benchmark choice \(a_{\rm cut}=r_c\) is retained only as a legacy ultraviolet regularization scale.  It must not be read as \(a_{\rm core}=r_c\), and it is not a theorem-level derivation of Pauli exclusion from ideal Euler dynamics.

        As written, Eq.~\eqref{eq:pauli_barrier_numeric} is an order-of-magnitude, cutoff-calibrated benchmark attached to the canonical constants. A physical resolved-core calculation must provide \(a_{\rm core}\), the tube profile, and the relation between that profile and \(a_{\rm cut}\) independently.

    \subsection{Numerical Benchmark Layer}
        \label{subsec:integration_benchmarks}

        \textbf{[DERIVED]} Computational-program requirement: earlier SST work contained explicit benchmark tables for particle and
        atomic masses. In v0.8.1, the benchmark layer is retained as a reproducibility requirement rather than as a
        standalone authority claim.

        \begin{table}[h]
        \centering
        \caption{Benchmark summary structure to be populated from canonical scripts and archived CSV output.}
        \label{tab:benchmark_summary_structure}
        \begin{tabular}{lrrrrc}
        \hline
        Particle & $m_{\mathrm{exp}}$ (MeV) & $m_{\mathrm{SST}}$ (MeV) & rel.\ error & mode & class \\
        \hline
        $e$   & 0.51099895 & 0.51099895 & 0 & predictive/closure & $3_1$ \\
        $\mu$ & 105.6583745 & \dots & \dots & predictive & candidate torus class \\
        $\tau$& 1776.86 & \dots & \dots & predictive & candidate torus class \\
        $p$   & 938.272088 & \dots & \dots & predictive or closure & composite class \\
        $n$   & 939.565420 & \dots & \dots & predictive or closure & composite class \\
        \hline
        \end{tabular}
        \end{table}

        \textbf{[CRITICAL NOTE]} A future v0.8.x benchmark appendix should include:
        \begin{itemize}
            \item the exact script or package path used for generation,
            \item archived CSV outputs,
            \item mode definitions (predictive vs.\ exact-closure),
            \item uncertainty propagation on canonical constants and geometric factors,
            \item theorem-style pass/fail metadata (including tail-quality checks, postchecks, and extrapolation summaries),
            \item explicit separation between practical best-estimate closure and theorem-level closure,
            \item normalization of legacy constant names into SST notation (\(\rhoF,\rc,\Gamma_0\)) while preserving numerical values.
        \end{itemize}

        This requirement is included to keep the benchmark layer falsifiable and reproducible, rather than merely illustrative.

    \subsection{Gauge-Sector Roadmap}
        \label{subsec:integration_gauge}

        \textbf{[SPECULATIVE]} Structural roadmap: a retained minimal statement of the gauge program is:
        \begin{align}
            \mathfrak{g}_{\mathrm{eff}}
            =
            \mathfrak{su}(3)\oplus\mathfrak{su}(2)\oplus\mathfrak{u}(1),
        \end{align}
        with effective gauge phases organized as holonomy data of multi-director transport.

        A minimal phase-based representation is
        \begin{align}
            A_\mu^{\,a}(x) = \partial_\mu \theta^a(x),
            \label{eq:gauge_phase_potential}
        \end{align}
        and clock coupling acts as a common phase reference,
        \begin{align}
            \theta^a(x) \mapsto \theta^a(x) + q^a \chi(x)
            \quad\Rightarrow\quad
            A_\mu^{\,a} \mapsto A_\mu^{\,a} + q^a \partial_\mu \chi.
            \label{eq:gauge_clock_shift}
        \end{align}

        \textbf{[CRITICAL NOTE]} The gauge layer is retained in v0.8.1 only as a roadmap statement. Running couplings,
        anomaly cancellation, and realistic mixing matrices remain open work.

        \paragraph{Research-track contact-channel guard.}
        Any QCD-like factor-nine or confinement-sheet tension claim must be routed through the resolved colored contact-channel rank test in Sec.~\ref{sec:rt_colored_contact_rank_nine}.  In the core canon this remains \textbf{[RESEARCH TRACK / MATCHING DIAGNOSTIC]}: the number is not inferred from \(\dim SU(3)=8\), nor from a fitted hadronic scale, but only from a future positive-rank condition on the finite-thickness contact-stress matrix.

    \section{Discussion and Limits}
        \label{sec:discussion}

    \subsection{Status of the recovered v0.5.12 sectors}

        The material reintroduced here does not have uniform epistemic status.

        \begin{itemize}
            \item The Chronos--Kelvin and pressure-law sectors are orthodox fluid-mechanical structures with SST reinterpretations.
            \item The analogue-metric and clock-field EFT layers are orthodox comparison tools, but only partially derived within SST.
            \item The R/T measurement bridge and weak-mixing stiffness ansatz remain explicitly speculative.
        \end{itemize}

    \subsection{What v0.8.1 now recovers from v0.5.12}

        After the present patch, v0.8.1 no longer omits the main v0.5.12 kernel sectors that are structurally compatible with the current architecture:
        \begin{itemize}
            \item Chronos--Kelvin / clock transport,
            \item zero-parameter reparameterization discipline,
            \item Swirl--EM and photon / unknot bridge,
            \item swirl pressure law,
            \item analogue metric and clock-field gravity bridge,
            \item R/T measurement bridge,
            \item weak-mixing bridge ansatz.
        \end{itemize}

    \subsection{Open items that remain outside the present closure}

        The following remain open and are \emph{not} upgraded to full canonical closure in this draft:
        \begin{itemize}
            \item a first-principles derivation of the full SM gauge sector;
            \item a completed EWSB derivation tied to the canon constants;
            \item a rigorous derivation of the R/T measurement kernel from continuum dynamics;
            \item full CANON promotion of the dark-sector taxonomy and galactic closure beyond the current BRIDGE labels (Section~\ref{sec:dark-canonicalization-status}).
        \end{itemize}

    \subsection{What is established at canon level}
        The present document supports the following as canon-level structures:
        \begin{itemize}
            \item the primitive constant chain $(\rhoF, v_{\swirlarrow}, \omega_c) \to R_{\mathrm{horn}}\equiv r_c \to \Gamma_0$,
            \item the Swirl-Clock as the coarse-grained relational time factor,
            \item delay-induced mode selection as the preferred discreteness mechanism,
            \item the atomic bridge and spectroscopic sectors as compatibility filters,
            \item the topology-first program for mass organization.
        \end{itemize}
        These claims are not equally proven, but they are the sectors on which the rest of the canon is structurally built.

    \subsection{What remains closure-level or conjectural}
        Several central SST ambitions remain open:
        \begin{enumerate}
            \item a rigorous derivation of full particle spectra from topology alone,
            \item a first-principles derivation of charge assignments and gauge couplings,
            \item a complete many-body derivation of shell filling and exclusion,
            \item a full demonstration that delay plus topology reproduces the observed quantum rule set without importing it.
        \end{enumerate}
        The canon records these as active programs rather than as settled conclusions.

    \subsection{Compatibility limits}
        Any viable SST realization must remain compatible with:
        \begin{itemize}
            \item atomic spectroscopy,
            \item relativistic effective kinematics,
            \item precision electroweak data where an SST bridge is claimed,
            \item the absence of large low-energy deviations from standard quantum mechanics.
        \end{itemize}
        This is why the canon repeatedly imposes suppression, weak-coupling, and low-energy recovery conditions.

    \subsection{Minimal falsifiers}
        The present canon would be materially weakened by any of the following:
        \begin{itemize}
            \item failure of the constant chain to remain dimensionally and numerically self-consistent,
            \item absence of any workable suppression mechanism for internal modes at atomic energies,
            \item inability to formulate a non-circular route from delay dynamics to discrete stable spectra,
            \item systematic contradiction between any SST atomic bridge and precision spectroscopy.
        \end{itemize}
        Conversely, strong support would come from a successful theorem ladder showing how discrete spectra, topological stabilization, and low-energy effective quantum structure arise from one shared substrate model.

    \subsection{Methodological limit}
        The canon is deliberately hybrid in style. Some sectors are written as exact closures, some as effective models, and some as programmatic roadmaps. This is a feature rather than a flaw, provided the epistemic status of each claim remains explicit. The document should therefore not be judged as though every section were intended to be a stand-alone journal theorem.

    \subsection{Near-term priorities}
        The next canon-strengthening tasks are clear:
        \begin{enumerate}
            \item complete the theorem ladder for the mass and clock sectors,
            \item expand numerical validation blocks using the canonical constants,
            \item refine the atomic bridge into sharper benchmark observables,
            \item decide which optional sectors belong in the core canon and which should migrate to separate technical notes or the research-track companion.
        \end{enumerate}

    \subsection{Final limit statement}
        SST is presently best understood as a structured unification program with several strong internal bridges and several open derivational gaps. The canon is therefore mature enough to organize the theory coherently, but not yet so mature that every phenomenological sector should be presented as theorem-complete.

    \subsection{Status of Legacy Recoveries}

        The additional material imported here from older canon layers does not have uniform status.

        \begin{itemize}
            \item The analogue-metric bridge and ribbon identity are mathematically orthodox tools.
            \item Their SST interpretation is partially derived but not fully closed dynamically.
            \item The thermal-radiation and QED-interface sectors remain explicitly research-track.
        \end{itemize}

        This section therefore serves as a controlled retention layer, not as a blanket validation of all legacy material.

        Accordingly, these recoveries should be read as structured extensions of the minimal core, not as fully promoted final doctrine.

    \subsection{What v0.8.1 Now Retains}

        After integration, the present Canon retains the following prior-development sectors in current-canon form:
        \begin{itemize}
            \item a particle-candidate topological dictionary (Section~\ref{subsec:integration_knotmap});
            \item a hydrodynamic exchange-energy estimate for fermionic exclusion
            (Section~\ref{subsec:integration_pauli});
            \item a benchmark / reproducibility layer (Section~\ref{subsec:integration_benchmarks});
            \item a minimal gauge-sector roadmap (Section~\ref{subsec:integration_gauge});
            \item a clock-field EFT bridge and Poisson gravity limit (Section~\ref{sec:time}).
        \end{itemize}

    \subsection{What is Deliberately Not Claimed}

        The present document still does \emph{not} claim:
        \begin{itemize}
            \item a completed first-principles derivation of the Standard Model;
            \item a unique finalized knot-to-particle dictionary;
            \item a closed proof of the exclusion barrier beyond the current scaling benchmark;
            \item a fully renormalized relativistic field theory for the gauge and clock sectors.
        \end{itemize}

        These omissions are deliberate. v0.8.1 is intended to remain logically strict even when phenomenological
        sectors are reintroduced.

    \subsection{Canon Rule Going Forward}

        Any further reintegrated material must satisfy all of the following:
        \begin{enumerate}
            \item it can be written without citing older internal canon versions as authorities;
            \item it can be assigned an explicit epistemic label;
            \item it includes either an internal derivation path or a reproducible benchmark route;
            \item it preserves compatibility with the orthodox limits recorded earlier in the document.
        \end{enumerate}

% ================================================================================= %
% APPENDICES
% ================================================================================= %
\newpage
\section*{Appendices}
    \addtocontents{toc}{\protect\setcounter{tocdepth}{1}}
    \setcounter{tocdepth}{1}
    \appendix

    \section{Topological reference toolkit}
        \label{sec:appendix_topology_toolkit}

        The following orthodox items form the mandatory internal mathematical toolkit cited by topology-first SST sectors. Statements are standard; SST-specific mappings appear in the main text with explicit status labels.

        \begin{enumerate}
            \item Gauss linking integral and linking number for disjoint closed curves.
            \item Helicity $\mathcal{H}=\int \mathbf{v}\cdot(\nabla\times\mathbf{v})\,dV$ in ideal (inviscid) fluid flow and its conservation in the appropriate limit.
            \item Thin-tube reduction of helicity to flux--linkage form in the standard filament limit.
            \item Hopf invariant for localized field configurations (degree--type integral) where used as a topological density.
            \item Basic torus-knot invariants as needed in the mass program (e.g.\ minimal crossing number, braid index, genus) with fixed naming consistent with Rolfsen notation in the main text.
            \item C\u{a}lug\u{a}reanu--White--Fuller relation $Lk=Tw+Wr$ for framed ribbons (already recalled as Eq.~\eqref{eq:calugareanu} in the integration layer).
        \end{enumerate}

        \textbf{[ORTHODOX]}: each item is classical mathematics \cite{Moffatt1969,White1969,Batchelor1967}. \textbf{[DERIVED]}: scope and notation for v0.8.x are frozen to this list plus house symbols $(\rhoF,\Gamma,\rc,\mathcal{H})$ without parallel aliases.

    \section{Full Derivations}
        \label{sec:appendix_derivations}

        \subsection{Canonical horn-radius constant chain}
            The present canon uses the closure chain
            \begin{align}
                r_e &= \frac{\alpha \hbar}{m_e c}, \\
                r_c &= \frac{r_e}{2}, \\
                \omega_c &= \frac{m_e c^2}{\hbar}, \\
                v_{\swirlarrow} &= \omega_c r_c, \\
                \Gamma_0 &= 2\pi r_c v_{\swirlarrow}.
            \end{align}
            Substituting $r_c = \alpha\hbar/(2m_e c)$ into $v_{\swirlarrow}=\omega_c r_c$ gives
            \begin{align}
                v_{\swirlarrow}
                =
                \left(\frac{m_e c^2}{\hbar}\right)
                \left(\frac{\alpha \hbar}{2m_e c}\right)
                = \frac{\alpha}{2}c,
            \end{align}
            hence the canonical ratio
            \begin{align}
                \eta_0 = \frac{v_{\swirlarrow}}{c} = \frac{\alpha}{2}.
            \end{align}

        \subsection{Horn-Envelope Density Normalization}
            The electron rest-energy normalization may be written as an effective horn-envelope density relation,
            \begin{align}
                \rho_{\mathrm{horn}}^{\mathrm{eff}}
                \equiv
                \frac{m_e c^2}{2\pi \vchar{}^2 R_{\mathrm{horn}}^3}.
                \label{eq:rho_horn_eff}
            \end{align}
            This relation is dimensionally exact,
            \begin{align}
                \left[\frac{m_e c^2}{v^2 r^3}\right]
                =
                \mathrm{kg\,m^{-3}},
            \end{align}
            but its interpretation is now restricted.

            \textbf{[CANON / SEMANTIC RULE]}: Eq.~\eqref{eq:rho_horn_eff} is an effective horn-envelope or normalization density. It is not a measured local material density of the resolved vortex tube. Any resolved-tube model must use $a_{\mathrm{core}}$, a profile function, or ropelength/finite-thickness data rather than silently identifying $a_{\mathrm{core}}=\rc$.

            \textbf{[CANON / MASS-FUNCTIONAL NOTE]}: Existing mass-functionals that use $\rho_{\mathrm{core}}\rc^3$ should be read as horn-envelope normalizations unless a separate local-core profile is explicitly introduced.

        \subsection{Geometric gate identity}
            Using the Compton wavelength $\lambda_c = h/(m_e c)$ and the horn-radius closure $r_c=R_{\mathrm{horn}} = \alpha\hbar/(2m_e c)$, one obtains
            \begin{align}
                \frac{\lambda_c}{\pi r_c}
                =
                \frac{h/(m_e c)}{\pi \alpha \hbar /(2m_e c)}
                =
                \frac{2h}{\pi \alpha \hbar}
                =
                \frac{4}{\alpha}.
            \end{align}
            This is the canonical geometric shielding factor used throughout the late mass-functional papers.

        \subsection{Useful numerical anchors}
            With the canonical values
            \begin{align}
                v_{\swirlarrow} &\approx 1.09384563\times 10^6\ \mathrm{m\,s^{-1}}, \\
                r_c &\approx 1.40897017\times 10^{-15}\ \mathrm{m}, \\
                \Gamma_0 &\approx 9.6836\times 10^{-9}\ \mathrm{m^2\,s^{-1}},
            \end{align}
            one has the Compton turnover time
            \begin{align}
                \tau_{\rm circ,c} = \frac{2\pi \rc}{\vchar} = \frac{2\pi}{\omega_c},
            \end{align}
            and the geometric gate
            \begin{align}
                \frac{\lambda_c}{\pi \rc} \approx 5.48\times 10^2.
            \end{align}

    \section{Delay Mathematics}
        \label{sec:appendix_delay}

        \subsection{Fixed-point condition for locked modes}
            Starting from the delayed phase equation
            \begin{align}
                \dot{\phi}(t) = \omega_0 + \kappa \sin\big(\phi(t-\tau_{\rm circ})-\phi(t)\big),
            \end{align}
            seek a rigidly rotating solution
            \begin{align}
                \phi(t)=\Omega t + \phi_\star.
            \end{align}
            Substitution gives the algebraic mode condition
            \begin{align}
                \Omega = \omega_0 - \kappa \sin(\Omega\tau_{\rm circ}),
            \end{align}
            where the minus sign follows from $\sin[-\Omega\tau_{\rm circ}]=-\sin(\Omega\tau_{\rm circ})$. This admits multiple branches labelled by the effective phase winding across the delay interval.

        \subsection{Linear stability of a locked branch}
            Let
            \begin{align}
                \phi(t) = \Omega t + \epsilon(t),
                \qquad |\epsilon|\ll 1.
            \end{align}
            Linearizing yields
            \begin{align}
                \dot{\epsilon}(t)
                =
                \kappa \cos(\Omega\tau_{\rm circ})
                \big(\epsilon(t-\tau_{\rm circ})-\epsilon(t)\big).
            \end{align}
            Using the exponential ansatz $\epsilon(t)\propto e^{\lambda t}$ gives the characteristic equation
            \begin{align}
                \lambda
                =
                \kappa \cos(\Omega\tau_{\rm circ})
                \left(e^{-\lambda\tau_{\rm circ}}-1\right).
            \end{align}
            Stability requires $\Re(\lambda)<0$ for all nontrivial roots. Thus discrete spectra in SST are not arbitrary: only those branches satisfying both the fixed-point equation and the stability inequality survive dynamically.

        \subsection{Small-delay expansion}
            For $|\lambda\tau_{\rm circ}|\ll 1$,
            \begin{align}
                e^{-\lambda\tau_{\rm circ}} = 1 - \lambda\tau_{\rm circ} + \mathcal{O}((\lambda\tau_{\rm circ})^2),
            \end{align}
            so the characteristic equation becomes
            \begin{align}
                \lambda \approx -\kappa\tau_{\rm circ}\cos(\Omega\tau_{\rm circ})\,\lambda.
            \end{align}
            Nontrivial damping therefore requires the effective factor $\kappa\tau_{\rm circ}\cos(\Omega\tau_{\rm circ})$ to have the correct sign. This provides a simple branch-selection criterion at weak delay.

        \subsection{Interpretive consequence}
            The important canon point is that discreteness is produced in two stages:
            \begin{align}
                \text{delay} \;\longrightarrow\; \text{multiple algebraic branches} \;\longrightarrow\; \text{stability filter}.
            \end{align}
            Quantization is therefore not inserted by hand; it is the residue of a dynamical selection problem on a delayed closed carrier.

    \section{Conversation-Derived Insights}
        \label{sec:appendix_conversations}

        \subsection{Role of this appendix}
            This appendix records recurrent insights that emerged during the project-level development of the canon and are useful for guiding future work, even when they do not yet qualify as theorem-level results.

        \subsection{Topology first, layers second}
            A central refinement of the recent canon is the separation
            \begin{align}
                \text{particle identity} \;\longleftrightarrow\; \text{topology},
                \qquad
                \text{excitation or dressing state} \;\longleftrightarrow\; \text{layer index}.
            \end{align}
            This prevents Golden-layer fits from being mistaken for a species ontology and keeps the topology program structurally cleaner.

        \subsection{Spectral origin of \texorpdfstring{$\zeta(3)$}{zeta(3)}}
            The project discussions repeatedly converged on the conclusion already encoded in Section~\ref{subsec:spectral_locked_carriers}: the appearance of $\zeta(3)$ is best interpreted as a spectral invariant of a cubic-weighted closed-carrier mode tower, not as a direct invariant of threefold geometry by itself.

        \subsection{Atomic bridge status discipline}
            Another recurring project insight is that the atomic bridge is strongest when treated as a benchmark layer with explicit status tags. Hydrogenic consistency, branch-sensitive couplings, and the spectroscopic gap form a robust bridge; shell filling, exclusion, and full heavy-atom structure remain open programs rather than completed derivations.

        \subsection{Master-equation viewpoint}
            A major organizational gain of the recent conversations is the recognition that many SST sectors can be read as different projections of one shared structure:
            \begin{align}
                \text{medium scale} \times \text{topology} \times \text{geometry} \times \text{clocking}.
            \end{align}
            Even when this is not written as a single canonical equation inside the present draft, it remains a useful guide for reducing duplication across the mass, time, and thermodynamic sectors.

        \subsection{State-transition taxonomy from the trefoil-closure discussion}
            The present conversation also fixes a useful state-transition taxonomy:
            \begin{align}
                \text{smooth deformation/binding} &: K=3_1\ \text{preserved},\\
                \text{R-to-T closure/nucleation} &: K\ \text{created at a boundary event},\\
                \text{annihilation or weak capture} &: K\ \text{converted with compensation}.
            \end{align}
            This taxonomy prevents ordinary atomic binding, non-Euler closure,
            and true particle conversion from being conflated.

        \subsection{Dual-vacuum and photon sector status}
            The project-level discussions also clarified a canon policy: the photon/torsion bridge and the dual-vacuum phenomenology are valuable and fertile, but they should not silently override the simpler core canon. They are best treated as phenomenological extensions built on top of the primitive constant chain and Swirl-Clock structure.

        \subsection{Canon edition notes}

        \subsubsection{v0.8.20}
            \textbf{v0.8.20} adds Chronos--Kelvin topological hitting conditions with a first-hitting kink-versus-reconnection split and Kairos boundary ledger (\(\ref{subsubsec:chronos_kelvin_hitting_conditions}\), \(\ref{eq:sst_first_hitting_time}\)), a contact-pressure saturation and Swirl-Clock kinematic shielding block (\(\ref{subsec:rt_contact_pressure_saturation_swirl_clock}\), \(\ref{eq:rt_contact_pressure_split}\)), a main-text research-track guard and colored rank-nine contact-channel diagnostic (\(\ref{sec:rt_colored_contact_rank_nine}\), \(\ref{eq:rt_rank9_factorization}\)), a Separatrix Surface-Density Lift monopole DtN research-track route (\(\ref{subsec:ssdl_monopole_dtn}\)) with audit v0.2 numeric support, and a canon-safe Research-Track particle/clock ansatz with preregistered mass falsifiers (\(\ref{app:research-track}\), \(\ref{app:minimal-numerical-programme}\)) on top of v0.8.19.

        \subsubsection{v0.8.19}
            \textbf{v0.8.19} adds an orthodox radius-convention ropelength/thickness foundation, a main-text guard that knot type alone does not determine a unique physical representative, a scalarized research-track screening functional \(E_{\rm screen}\) with \(Q_G\) retained only as a sector label unless mapped to \(I_G\), a derived leading thin-filament length/log energy anchor, a density and dimensional audit for geometric mass branches, a geometric impedance bridge demoted to research-track pointer form with an explicit \(T_{\rm core}\) free-symbol guard, stricter reintegration rules for derived Hamiltonian terms versus screening surrogates, a consolidated audit bundle (unified \(\SwirlClock\) discipline, F\_max hygiene, Poisson--clock bridge, EMG guards, trefoil precision, pure-geometric mass collapse), a Foucault-pendulum rotation-to-clock inference protocol in the research track, projected swirl-vorticity frequency diagnostics, symmetry-metadata guards for dark-knot taxonomy, benchmark CSV schema hardening, operator-valued Swirl-Clock visibility diagnostics, a cosmological-impedance route for \(\rho_f\) (\S\ref{sec:rt_rhof_cosmological_impedance_route}, Eq.~\ref{eq:rt_rhof_lambda_chi}) with horn-effective density label hygiene, and Taylor-column / rotating-fluid Swirl-Clock parity benchmarks (\S\ref{sec:taylor_column_finite_thickness_swirl_strings}, \S\ref{sec:rotating_fluid_swirl_clock_parity}) on top of v0.8.18.

        \subsubsection{v0.8.18}
            \textbf{v0.8.18} adds calibration-chain guardrails (rc/$\alpha$ gates, kernel normalization, EM prefactor, $\rhohorn$ baseline fix), research-track observational falsifier bounds (GW170817, pulsar preferred-frame limits), orthodox resolved-tube reach/thickness definitions, the four topology/stability research appendices (non-Abelian protection, reconnection decay, fractional torus-phase winding, and topological-fluid mechanics), the research-track knot-energy normal form with $I_G$ scalarization, Euler/Magnus probe-transport polish, canonical time ontology with no-bare-$\tau$ policy, golden-layer hyperbolic rapidity guards, Planck-suppressed $F_{\rm swirl}^{\max}$ reparameterization, Compton--horn $G_{\mathrm{swirl}}$ reduction guard and balance-angle diagnostic, hybrid density-source Swirl-Clock benchmark, and the contact-stress geometry research appendix on top of v0.8.17.

        \subsubsection{v0.8.17}
            \textbf{v0.8.17} adds the T-foliation versus local Swirl-Clock remark (Eq.~\ref{eq:absolute_T_local_tau_layers}), the ropelength/trefoil-$\alpha$ convention guard, Route-I SST-63/23/56 integration in the research track, and aligned updates in SST-73, SST-63, SST-64 v2, and SST-34 on top of v0.8.16. RC2 release-clean polish: p.119 summary quote, `\AE{}`-modes fix, `tabularx` particle-candidate tables, and targeted version-phrase cleanup.

        \subsubsection{v0.8.16}
            \textbf{v0.8.16} adds the pressure--optical locking relation (Eq.~\ref{eq:canonical_pressure_optical_locking}) in the Triadic Gravity-Response Corollary, the research-track no-monopole transport-stress audit (\S\ref{sec:rt_no_monopole_transport_stress_audit}), and the relativity falsifier ladder on top of v0.8.15. SST-73 receives matching critical notes on Candidate C and $[\alpha_{\mathrm{grav}}]=\mathrm{s^2}$.

        \subsubsection{v0.8.15}
            \textbf{v0.8.15} adds the canon hydrodynamic Lorentz-type swirl-force density (\S\ref{sec:lorentz_type_force_density_as_swirl_stress}: Bernoulli tie, loop observables, calibrated scales) and the research-track EM-to-swirl correspondence with $\lambda_{\mathrm{EM}\to\circlearrowleft}$, SST-44 stress closure, and pending $\mathbf f_{\mathrm{link}}$ (\S\ref{sec:research_em_to_swirl_force_density_correspondence}) on top of v0.8.14.

        \subsubsection{v0.8.14}
            \textbf{v0.8.14} adds the two-speed clock discipline (\S\ref{subsec:two_speed_clock_discipline}: $\vchar$ vs.\ $c$, NLSE core sound speed $c_s=\vchar/\sqrt{2}$, research-track gate), the $\alpha$-gate guard on $\lambda_c/(\pi\rc)=4/\alpha$, and the research-track core--torsion impedance-matching lemma (\S\ref{sec:rt_core_torsion_impedance_matching}) on top of v0.8.13.

        \subsubsection{v0.8.13}
            \textbf{v0.8.13} adds the relativity-emergence audit: the internal tensor-speed naturalness proposition (\S\ref{subsec:tensor_speed_naturalness}, conditional $c_{13}=0$) in the main canon, the research-track Relativity Emergence Ladder (\S\ref{sec:rt_relativity_emergence_ladder}) separating derivable SR/GR kinematics from the open Einstein-dynamics program, and relativity/induced-gravity/cosmology bibliography on top of v0.8.12.

        \subsubsection{v0.8.12}
            \textbf{v0.8.12} applies the Gemini round-3 epistemic audit: $\mathcal{P}_{\mathrm{cal}}$-aligned identity relabels, coarse-grain density caveat, Pauli $a_{\rm cut}$ disambiguation, galaxy-scale $\rhoF$ note, and research-track EM-bridge limits on top of v0.8.11.

        \subsubsection{v0.8.11}
            \textbf{v0.8.11} completes the final hygiene pass: consistent $\mathcal{P}_{\mathrm{cal}}$, $\vchar$ vs.\ $\uswirl$ usage in EMG, framed self-linking, W-boson sector, and research-track numerical bridges on top of v0.8.10.

        \subsubsection{v0.8.10}
            \textbf{v0.8.10} applies the Gemini round-2 audit: canonical speed $\vchar$ vs.\ local $\uswirl(x,t)$ discipline, $\eta_0$ algebraic-within-calibration status, delay-equation sign convention, Rydberg/finite-cell relabeling, and research-track $v_K$ disambiguation on top of v0.8.9.

        \subsubsection{v0.8.9}
            \textbf{v0.8.9} adds the triadic gravity-response corollary (\S\ref{subsec:triadic_gravity_response_corollary}) and research-track flame/caustic/shell diagnostics plus atomic condensate closure notes on top of v0.8.8.

        \subsubsection{v0.8.8}
            \textbf{v0.8.8} applies the Gemini audit: primitive calibration notation ($\mathcal{P}_{\mathrm{cal}}$, $\vchar$, $\uswirl$), conditional Kelvin/DDE labels, Pauli $a_{\rm cut}$ disambiguation, and research-track epistemic relabeling on top of v0.8.7.

        \subsubsection{v0.8.7}
            \textbf{v0.8.7} adds the scalar-vs-$\mathbb{C}P^1$ substrate argument, the $\theta=\pi$ fermionic $\mathbb{Z}_2$ sector, and supporting bibliography on top of v0.8.6.

        \subsubsection{v0.8.6}
            \textbf{v0.8.6} adds framed-tube self-linking ($SL=pq$), spinorial $\chi=\pm1$ selection, and the lepton ladder (\S\ref{subsec:framed_selflinking_spinorial}) on top of v0.8.5.

        \subsubsection{v0.8.5}
            \textbf{v0.8.5} adds the highres conversation audit (Euler discipline, R-to-T bookkeeping, twist-ladder research track, horn/envelope wording) and the CALIBRATED circularity-honesty relabeling on top of v0.8.4.

        \subsubsection{v0.8.4}
            \textbf{v0.8.4} integrates the canonical EM--gravity bridge, finite-cell $\alpha$ obstruction summary, and extended research-track taxonomy blocks.

        \subsubsection{v0.8.3}
            \textbf{v0.8.3} integrates framed-tube ontology, attachment gate, trefoil closure discipline, updated particle dictionary, and Pauli $a_{\rm core}$ regularization per conversation, highres, and trefoil diffs.

        \subsubsection{v0.8.2}
            \textbf{v0.8.2} integrates radius/density semantic correction: $\rc \equiv R_{\mathrm{horn}}$, separate $a_{\mathrm{core}}$, and horn-envelope density normalization per conversation and highres diffs.


        \subsection{Trefoil closure and persistence discipline}
            The project-level conclusion is that electron-trefoil creation, if
            retained, must be formulated as a boundary-supported R-to-T closure
            event rather than as ordinary reconnection of an already closed
            ideal-Euler vortex tube. The canon distinction is therefore
            \begin{align}
                \text{pre-closure R-phase support}
                &\xrightarrow{\mathcal{K}_{\rm Kairos}}
                \text{closed }3_1\text{ carrier},\\
                \text{closed }3_1\text{ carrier}
                &\xrightarrow{\rm smooth\ Euler}
                \text{same knot class }3_1 .
            \end{align}
            This keeps the orthodox frozen-in topology constraint intact while
            preserving the SST closure program as a clearly labeled bridge.

            The radius notation is correspondingly split into the local tube
            radius \(a_{\rm core}\), the neutral circulation radius \(\rc\), and
            the classical electron envelope radius \(r_e=2\rc\):
            \begin{align}
                a_{\rm core} \neq \rc,
                \qquad
                r_e=2\rc .
            \end{align}
            Ordinary atomic binding changes the phase envelope of the carrier,
            not its knot class:
            \begin{align}
                e^-_{\rm atom}
                =
                3_1^{\rm core}\otimes\Psi_{n\ell m\sigma}^{\rm envelope} .
            \end{align}


        \paragraph{Ropelength representative and thickness-constrained ground state.}
        For a tame knot or link type \(K\), orthodox ropelength theory proves the
        existence of a ropelength-minimizing representative
        \cite{CantarellaKusnerSullivan2002,DenneDiaoSullivan2006}.
        With the radius convention for thickness,
        \begin{align}
            \operatorname{Rop}_{\rm rad}(\gamma)
            =
            \frac{\operatorname{Len}(\gamma)}
                 {\operatorname{Thi}_{\rm rad}(\gamma)} ,
        \end{align}
        and the minimizer is \(C^{1,1}\), with bounded curvature, but need not be
        \(C^2\) \cite{CantarellaKusnerSullivan2002}.

        \textbf{[ORTHODOX GEOMETRY / RESOLVED-TUBE DEFINITION].}
        For a \(C^1\) embedded centerline, orthodox ropelength theory identifies
        thickness with both reach and normal injectivity radius
        \cite{CantarellaKusnerSullivan2002}.  Therefore the SST resolved-tube
        radius is, when an ideal-tube representative is being used,
        \begin{align}
            a_{\rm core}(K)
            &\equiv
            \operatorname{Thi}_{\rm rad}(\gamma_K)
            =
            \operatorname{reach}(\gamma_K)
            =
            \operatorname{NIR}(\gamma_K),
            \label{eq:a_core_reach_thickness}\\
            \operatorname{Thi}_{\rm rad}(\gamma_K)
            &=
            \min\!\left(\inf_s\frac{1}{\kappa(s)},\,
            \frac12 d_{\rm dcsd}(\gamma_K)\right),
            \label{eq:thickness_curvature_dcsd_guard}
        \end{align}
        where \(d_{\rm dcsd}\) is the doubly-critical self-distance.  Equations
        \eqref{eq:a_core_reach_thickness}--\eqref{eq:thickness_curvature_dcsd_guard}
        are geometric definitions of the resolved tube; they do not license the
        identification \(a_{\rm core}=\rc\).
        In an SST sector where the resolved vortex-string energy is
        tension dominated, with an orthodox knotted flux-tube analogue in the
        glueball literature \cite{BuniyKephart2003},
        \begin{align}
            E_{\rm SST}[\gamma]
            =
            T_{\rm eff}\operatorname{Len}(\gamma)
            +
            E_{\rm sub}[\gamma],
            \qquad
            \operatorname{Thi}_{\rm rad}(\gamma)\ge r_c ,
        \end{align}
        the leading-order ground-state problem reduces to ropelength minimization:
        \begin{align}
            E_{\rm SST}^{(0)}(K)
            \approx
            T_{\rm eff} r_c\,\operatorname{Rop}_{\rm rad}(K).
        \end{align}
        This is a derived-conditional reduction, not an independent proof that all
        shape-dependent observables are ropelength functions.

        \paragraph{Critical note: degeneracy and subleading splitting.}
        Ropelength minimizers need not be unique.  Therefore, while the scalar
        mass proxy \(M_K\propto\operatorname{Rop}(K)\) may remain well-defined,
        shape-sensitive observables such as dipole structure, magnetic moment,
        moment of inertia, writhe distribution, or local swirl-stress maps may
        require subleading splitting terms:
        \begin{align}
            E_{\rm sub}
            =
            E_{\rm bend}
            +
            E_{\rm twist}
            +
            E_{\rm Biot\text{-}Savart}
            +
            E_{\rm core}
            +
            \cdots .
        \end{align}

        \paragraph{Convention guard for trefoil alpha comparisons.}
        The mathematical literature uses both radius- and diameter-based thickness
        conventions.  If
        \begin{align}
            \operatorname{Thi}_{\rm diam}=2\operatorname{Thi}_{\rm rad},
        \end{align}
        then
        \begin{align}
            \operatorname{Rop}_{\rm diam}
            =
            \frac12\operatorname{Rop}_{\rm rad}.
        \end{align}
        For the trefoil, using certified ideal-knot data
        \cite{AshtonCantarellaPiatekRawdon2011},
        \begin{align}
            \operatorname{Rop}_{\rm rad}(3_1)\approx32.7436,
            \qquad
            \operatorname{Rop}_{\rm diam}(3_1)\approx16.3718 .
        \end{align}
        The leading coincidence
        \begin{align}
            \alpha^{-1}_{\rm lead}
            =
            \frac{8\pi}{3}\operatorname{Rop}_{\rm diam}(3_1)
            \approx
            137.1561
        \end{align}
        is therefore a robust sub-promille numerical coincidence, not an exact
        or ppm-certified derivation of the fine-structure constant.

        \subsection{Finite-cell $\alpha$ coincidence and obstruction summary}
            \label{subsec:finite_cell_alpha_obstruction}

            \textbf{[BRIDGE ANSATZ / SPECULATIVE FIT].}
            A parameter-light finite-cell variational construction for an ideal trefoil filament,
            containing partially open modelling choices, yields a leading dimensionless scale
            \begin{align}
                \alpha^{-1}_{\mathrm{lead}}
                =
                \frac{8\pi}{3}\,\mathcal{L}_{3_1}
                \approx 137.15471,
                \label{eq:alpha_lead_finite_cell}
            \end{align}
            using the numerically computed ideal trefoil ropelength
            $\mathcal{L}_{3_1}=16.371637$ \cite{AshtonCantarellaPiatekRawdon2011}
            as a \emph{benchmark} topological-geometric input.
            This lies within $866$~ppm of CODATA $\alpha^{-1}=137.035999$ and is an
            $\alpha$-free sub-per-mille coincidence, not a part-per-million derivation.

            The prefactor $8\pi/3$ contains a mode-count factor $N_p=4$ derived from the
            soft sector of the spherical finite-cell surface spectrum
            $k_\ell=\ell(\ell+1)-2$ (zero and unstable modes only). The sector volume
            $4\pi/3$ and overall normalization remain partially open modelling choices.

            \textbf{[NEGATIVE / OBSTRUCTION RESULT].}
            Closing the residual to ppm accuracy is underdetermined: at fixed cell-radius
            gate the required transverse shell weight is $w_\perp=1.0675$, not the round
            value $1$ (which reaches only $10.5$~ppm). An independent Gross--Pitaevskii
            core-profile proxy returns $w_\perp\approx 2.076$, overshooting to $-157$~ppm.
            The $10.5$~ppm residual can be absorbed independently by any of five
            sub-percent coefficients (shell weight, cell-radius gate, sector volume,
            ropelength value, or a third-order term), so ppm closure carries no
            discriminating power without pinning all inputs to $\sim10^{-6}$ relative
            precision---beyond present ideal-knot ropelength certification
            \cite{AshtonCantarellaPiatekRawdon2011,CantarellaKusnerSullivan2002}.

            The gate hierarchy is therefore
            \begin{align}
                G_0\ (\text{input provenance})
                \longrightarrow
                G_1\ (\text{identifiability})
                \longrightarrow
                G_2\ (\text{operator Hessians})
                \longrightarrow
                G_3\ (\text{ppm closure}),
                \label{eq:finite_cell_gate_hierarchy}
            \end{align}
            with $G_0$ and $G_1$ blocking closure before Hessian-level work becomes relevant.

            The ring-constant proximity $\alpha_{\mathrm{ring}}\approx\varphi$ is treated in
            Subsection~\ref{subsec:g7_ring_constant_vs_trefoil_closure} of the
            research track. It is not a canon-level trefoil curvature closure.

            \textbf{[CANON / NON-CLAIM].}
            This block does not assert a first-principles derivation of the fine-structure
            constant. The defensible content is the $\alpha$-free sub-per-mille leading
            coincidence with partially derived $N_p=4$, plus a localized obstruction to
            ppm closure. Full scripts, appendices, and audit details are in
            \texttt{finite\_cell\_obstruction.tex} and \texttt{Derive\_Constants/}.
        \subsection{Open theorem ladder}
            The most important unfinished task identified across the project is not adding new sectors, but tightening the derivational ladder between existing ones. In practice this means:
            \begin{enumerate}
                \item turning heuristic closures into explicit lemmas,
                \item attaching numerical validation to each nontrivial scaling law,
                \item separating benchmark results from full ontological claims.
            \end{enumerate}
            This appendix therefore serves as a reminder that canon quality increases more by sharpening structure than by increasing conceptual breadth.

% ===========================
% SUPPLEMENTAL APPENDIX
% ===========================
 
        \newpage
    
    
    \section{Explicit Topological Mass Closures}
    \label{sec:appendix_explicit_mass}
    
    \subsection{Purpose and status}
        This appendix supplies an explicit closure layer for the mass program without replacing the canonical master equation of Eq.~\eqref{eq:sst_master_equation}. Its role is to make concrete those invariant choices that were historically used in pre-v0.8 layers whenever actual particle-mass calculations were attempted.
        
        \textbf{[DERIVED]} Canon rule: the abstract kernel \(\Xi_K\) of the main text may be specialized to more explicit invariant combinations provided the resulting branch remains dimensionless, benchmarkable, and subordinate to the topology-first hierarchy.
        
        \textbf{[CRITICAL NOTE]} Nothing in this appendix upgrades the explicit kernel to a theorem of first-principles filament dynamics. The present role is closure discipline and reproducible benchmarking.

    \subsection{Pure geometric baseline branch}
        \textbf{[DERIVED-CONDITIONAL / NORMALIZATION BRANCH].}
        A pure geometric baseline mass scale for topology class $T$ is
        \begin{align}
            M_0(T)
            =
            2\pi^3\,\rhohorn\,\frac{\rc^5}{\lambda_c^2}\,\mathcal{L}_{\mathrm{tot}}(T).
            \label{eq:pure_geometric_mass_baseline}
        \end{align}
        This branch is subordinate to Eq.~\eqref{eq:sst_master_equation}. It records a normalization-level geometric closure used in historical mass-program benchmarks; it does not replace the canonical master functional or the explicit kernel branch below. The density in Eq.~\eqref{eq:pure_geometric_mass_baseline} is the horn-envelope normalization density of Eq.~\eqref{eq:core_density}; using the background mass-equivalent density $\rhoM$ here would move the benchmark down by the factor $\rhoM/\rhohorn\simeq1.20\times10^{-30}$ and is therefore a symbol-collision error, not a physical branch.

    \subsection{Explicit invariant branch of the canonical kernel}
        A conservative explicit specialization is obtained by resolving the topological kernel into ropelength, braid complexity, and Seifert-genus data:
        \begin{align}
            \Xi_K^{(\mathrm{exp})}
            =
            \left[
                \alpha_{\mathcal{L}}
                \frac{\mathcal{L}_{\mathrm{tot}}(K)}{\mathcal{L}_{\mathrm{tot}}(3_1)}
                +
                \alpha_b \bigl(b(K)-1\bigr)
            +
                \alpha_g\, g(K)
            \right]
            \varphi^{-2k(K)}.
            \label{eq:explicit_kernel}
        \end{align}
        Here:
        \begin{itemize}
            \item \(\mathcal{L}_{\mathrm{tot}}(K)\) is the total ropelength-like invariant used for the chosen embedding or ideal representative,
            \item \(b(K)\) is the braid index,
            \item \(g(K)\) is the Seifert genus,
            \item \(k(K)\) is the canonical layer / dressing index already used in the main text.
        \end{itemize}
        
        \textbf{[ORTHODOX]} The invariants \(\mathcal{L}_{\mathrm{tot}}(K)\), \(b(K)\), and \(g(K)\) are standard topological or geometric knot descriptors.
        
        \textbf{[DERIVED / MATCHING-ANSATZ GUARD]} The dimensional admissibility of Eq.~\eqref{eq:explicit_kernel} is derived. The displayed three-coefficient form and the $\varphi^{-2k(K)}$ dressing are a benchmark matching ansatz; the coefficients $\alpha_{\mathcal{L}}$, $\alpha_b$, and $\alpha_g$ remain calibrated unless a separate variational or statistical-mechanical calculation derives them.
        
        Substituting Eq.~\eqref{eq:explicit_kernel} into Eq.~\eqref{eq:sst_master_equation} yields the explicit benchmark mass branch
        \begin{align}
            M_K^{(\mathrm{exp})}(x)
            =
            \rhoM(x)\,\rc^3\,\Pi_K\,\Xi_K^{(\mathrm{exp})}\,\SwirlClock(x)^{-2}.
            \label{eq:explicit_mass_branch}
        \end{align}
    
    \subsection{Baryon benchmark closure}
        For composite baryons built from constituent topology classes \(K_u\) and \(K_d\), a conservative additive benchmark branch is
        \begin{align}
            M_B^{(\mathrm{add})}
            =
            \Pi_B\,M_0
            \left(
                n_u\,\Xi_{K_u}^{(\mathrm{exp})}
                +
                n_d\,\Xi_{K_d}^{(\mathrm{exp})}
            \right),
            \label{eq:baryon_additive_branch}
        \end{align}
        where \(n_u\) and \(n_d\) are constituent counts and \(M_0\) is the common medium scale extracted from the main mass functional.
        
        A retained legacy golden-closure benchmark is
        \begin{align}
            M_p^{(\varphi)}
            &=
            \varphi^{-2} 3^{-1/\varphi}\,\bigl(2M_u + M_d\bigr),
            \label{eq:proton_phi_benchmark}\\
            M_n^{(\varphi)}
            &=
            \varphi^{-2} 3^{-1/\varphi}\,\bigl(M_u + 2M_d\bigr).
            \label{eq:neutron_phi_benchmark}
        \end{align}
        
        \textbf{[SPECULATIVE]} Equations~\eqref{eq:proton_phi_benchmark}--\eqref{eq:neutron_phi_benchmark} are retained as baryon-sector closure benchmarks only. They are not promoted to universal mass laws.
    
    \subsection{Canonical use policy}
        The explicit closures above are intended for:
        \begin{enumerate}
            \item theorem-ladder benchmarking,
            \item archived CSV reproduction,
            \item side-by-side comparison of topology branches,
            \item explicit baryon fits.
        \end{enumerate}
        They are \emph{not} intended to replace the abstract organizing form of the main text.
    
    \section{Effective Quantum Bridge}
    \label{sec:appendix_quantum_bridge}
    
    \subsection{Purpose and status}
        The main canon states that quantum mechanics should appear as an effective low-energy description of the medium. The present appendix supplies the minimal bridge formulas needed to make that statement mathematically legible without claiming a completed derivation of the full quantum rule set.
        
        \textbf{[ORTHODOX]} The Madelung decomposition and Schr\"odinger-form continuum limit are standard.
        
        \textbf{[DERIVED]} SST may use these formulas as a coarse-grained field representation once a local amplitude density and phase have been specified.
        
        \textbf{[SPECULATIVE]} The identification of the resulting fields with branch-resolved swirl carriers remains an SST interpretation.

        \textbf{[CRITICAL NOTE]}: $\rho_\psi$ is a coarse-grained mode-support density of the emergent carrier field, \emph{not} a compression of the incompressible substrate $\rhoF$ (Axiom~1, $\nabla\cdot\mathbf v=0$). The two densities are distinct objects, so the Madelung amplitude does not violate substrate incompressibility.
    
    \subsection{Amplitude--phase decomposition}
        Introduce an effective complex mode field
        \begin{align}
            \psi(x,t) = \sqrt{\rho_\psi(x,t)}\,e^{i\theta(x,t)},
            \label{eq:madelung_ansatz_appendix}
        \end{align}
        with \(\rho_\psi \ge 0\) and real phase \(\theta\). The associated effective transport velocity is
        \begin{align}
            \mathbf{u}_\psi
            =
            \frac{\hbar}{m_{\mathrm{eff}}}\nabla\theta.
            \label{eq:effective_phase_velocity}
        \end{align}
        
        Assuming a Schr\"odinger-form bridge equation
        \begin{align}
            i\hbar \partial_t \psi
            =
            \left(
                -\frac{\hbar^2}{2m_{\mathrm{eff}}}\nabla^2 + V_{\mathrm{eff}}
            \right)\psi,
            \label{eq:schrodinger_bridge_appendix}
        \end{align}
        one obtains the equivalent real system
        \begin{align}
            \partial_t \rho_\psi + \nabla\cdot\bigl(\rho_\psi \mathbf{u}_\psi\bigr) &= 0,
            \label{eq:continuity_bridge_appendix}\\
            \partial_t \theta
            +
            \frac{m_{\mathrm{eff}}}{2}\lVert \mathbf{u}_\psi\rVert^2
            + V_{\mathrm{eff}}
            + Q_{\mathrm{eff}} &= 0,
            \label{eq:hamjac_bridge_appendix}
        \end{align}
        with effective quantum potential
        \begin{align}
            Q_{\mathrm{eff}}
            =
            -\frac{\hbar^2}{2m_{\mathrm{eff}}}
            \frac{\nabla^2 \sqrt{\rho_\psi}}{\sqrt{\rho_\psi}}.
            \label{eq:quantum_potential_appendix}
        \end{align}
        
        \textbf{[ORTHODOX]} Equations~\eqref{eq:continuity_bridge_appendix}--\eqref{eq:quantum_potential_appendix} are the standard Madelung system associated with Eq.~\eqref{eq:schrodinger_bridge_appendix} \cite{Schrodinger1926,Madelung1927}.
    
    \subsection{SST reading of the bridge}
        The conservative SST interpretation is:
        \begin{align}
            \rho_\psi
            &\leftrightarrow
            \text{coarse-grained mode-support density},\\
            \theta
            &\leftrightarrow
            \text{effective circulation phase},\\
            \mathbf{u}_\psi
            &\leftrightarrow
            \text{phase-transport velocity of the low-energy envelope}.
        \end{align}
        
        At this level, the branch-resolved atomic sector of Section~\ref{sec:atomic} is read not as a replacement of the orthodox wavefunction, but as a candidate microscopic underpinning for the effective fields \(\rho_\psi\) and \(\theta\).
    
    \subsection{Two-component branch bookkeeping}
        A minimal two-branch extension consistent with the \(\sigma=\pm1\) structure of the atomic sector is
        \begin{align}
            \Psi
            =
            \begin{pmatrix}
                \psi_{+}\\
                \psi_{-}
            \end{pmatrix},
            \qquad
            \psi_{\pm}
            =
            \sqrt{\rho_{\pm}}\,e^{i\theta_{\pm}},
            \label{eq:two_component_bridge}
        \end{align}
        with total coarse-grained density
        \begin{align}
            \rho_{\mathrm{tot}} = \rho_{+}+\rho_{-}.
        \end{align}
        
        \textbf{[SPECULATIVE]} This two-component structure is the minimal SST bridge to spinor-like bookkeeping. It is not yet a first-principles derivation of fermionic spin statistics.
    
    \subsection{Topological reading of the bridge}
        The most conservative topology statement retained at appendix level is that nontrivial phase maps may encode nontrivial carrier organization. In particular, the effective phase field may carry nontrivial winding or linking data even when the coarse-grained description is written in wave language. This is the appendix-level reason why the canon treats topology and wave structure as compatible rather than as competing descriptions.
    
    \section{Dark-Sector and Galactic Continuum Law}
    \label{sec:appendix_dark_galactic}
    
    \subsection{Purpose and status}
        The main canon already retains the dark-knot taxonomy and the radial swirl pressure law. This appendix collects the next phenomenological layer: the galaxy-scale force law and the simplest rotation-curve ansatz inherited from earlier SST development.
        
        \textbf{[ORTHODOX]} The underlying force law is ordinary Euler balance in steady swirl.
        
        \textbf{[DERIVED]} SST reinterprets that balance as a candidate large-scale acceleration source.
        
        \textbf{[SPECULATIVE]} The explicit galaxy-fit profile below is a phenomenological closure ansatz, not a theorem.

        \textbf{[CRITICAL NOTE]} The canonical \(\rhoF=7.0\times10^{-7}\,\mathrm{kg\,m^{-3}}\) is an effective medium parameter, not an ordinary gravitating baryonic/cosmological rest-mass density. Applying it on galactic scales requires a separate coupling and screening law. Without that law, the galaxy-scale pressure profile is a fit ansatz only and must not be interpreted as a literal dense fluid through which stars experience ordinary hydrodynamic drag.
    
    \subsection{Effective dark acceleration law}
        Starting from the steady radial pressure law of Eq.~\eqref{eq:swirl_pressure_law},
        \begin{align}
            \frac{1}{\rhoF}\frac{dp_{\mathrm{swirl}}}{dr}
            =
            \frac{v_\theta(r)^2}{r},
        \end{align}
        define the effective SST acceleration by
        \begin{align}
            a_{\mathrm{SST}}(r)
            :=
            \frac{1}{\rhoF}\frac{dp_{\mathrm{swirl}}}{dr}.
            \label{eq:appendix_dark_accel}
        \end{align}
        Then
        \begin{align}
            a_{\mathrm{SST}}(r)=\frac{v_\theta(r)^2}{r}.
            \label{eq:appendix_dark_accel_v2r}
        \end{align}
        
        \textbf{[DERIVED]} Equations~\eqref{eq:appendix_dark_accel}--\eqref{eq:appendix_dark_accel_v2r} are immediate consequences of the retained pressure law.
    
    \subsection{Two-component galactic rotation ansatz}
        A retained phenomenological velocity profile is
        \begin{align}
            v_{\mathrm{swirl}}(r)
            =
            \frac{C_{\mathrm{core}}}{\sqrt{1+\bigl(r_c^{\mathrm{gal}}/r\bigr)^2}}
            +
            C_{\mathrm{tail}}
            \left(1-e^{-r/r_c^{\mathrm{gal}}}\right),
            \label{eq:galactic_velocity_ansatz}
        \end{align}
        with corresponding total circular speed
        \begin{align}
            v_{\mathrm{tot}}^2(r)
            =
            v_{\mathrm{bar}}^2(r)
            +
            v_{\mathrm{swirl}}^2(r).
            \label{eq:total_circular_speed_appendix}
        \end{align}
        
        \textbf{[SPECULATIVE]} Equation~\eqref{eq:galactic_velocity_ansatz} is a phenomenological fit ansatz and should be used only at the appendix / research-track level.
    
    \subsection{Relation to dark-knot taxonomy}
        The taxonomy of Eq.~\eqref{eq:dark_knot_taxonomy} classifies low-helicity, low-instability knot sectors as dark or quasi-dark candidates. The present galactic law does \emph{not} assume that all galactic anomalies are produced by localized dark knots alone. Rather, it provides the complementary continuum statement: long-range swirl pressure gradients may produce effective extra acceleration even when the underlying microscopic carriers are topologically sparse or weakly luminous.
    
    \subsection{Use policy}
        The galactic law belongs in the canon only as:
        \begin{enumerate}
            \item an appendix-level phenomenological sector,
            \item a separate astrophysics / dark-sector track,
            \item a fit target for future rotation-curve studies.
        \end{enumerate}
        It is deliberately not promoted in v0.8.18 to the same status as the primitive constant chain, the Swirl-Clock, or the delay-selection sector.
        
% ---------------------------
% RESEARCH TRACKS for future
% ---------------------------
        \newpage
        \newif\ifresearchtrackincluded
        \researchtrackincludedtrue
        \input{SST_CANON-v0.8.20-research-track}
        
% ---------------------------
% Bibliography
% ---------------------------
        \begin{thebibliography}{99}
            
            \bibitem{SST_delay_basic}
            J.~K. Hale and S.~M. Verduyn Lunel,
            \newblock \emph{Introduction to Functional Differential Equations},
            \newblock Springer (1993).
            
            \bibitem{CMS2026Wmass}
            CMS Collaboration,
            \newblock High-precision measurement of the W boson mass with the CMS experiment,
            \newblock Nature \textbf{652} (2026), 321--337,
            \newblock DOI: 10.1038/s41586-026-10168-5.
            
            \bibitem{PDG2024}
            S.~Navas et al. (Particle Data Group),
            \newblock Review of Particle Physics,
            \newblock Phys. Rev. D \textbf{110} (2024), 030001,
            \newblock DOI: 10.1103/PhysRevD.110.030001.
            
            \bibitem{deBlas2022}
            J.~de Blas et al.,
            \newblock Global analysis of electroweak data in the Standard Model,
            \newblock Phys. Rev. D \textbf{106} (2022), 033003,
            \newblock DOI: 10.1103/PhysRevD.106.033003.
            
            \bibitem{CDF2022Wmass}
            CDF Collaboration,
            \newblock High-precision measurement of the W boson mass with the CDF II detector,
            \newblock Science \textbf{376} (2022), 170--176,
            \newblock DOI: 10.1126/science.abk1781.
            
            \bibitem{Batchelor1967}
            G.~K. Batchelor,
            \newblock \emph{An Introduction to Fluid Dynamics},
            \newblock Cambridge University Press (1967).

            \bibitem{Taylor1922}
            G.~I. Taylor,
            \newblock The motion of a sphere in a rotating liquid,
            \newblock \emph{Proceedings of the Royal Society of London A} \textbf{102} (1922), 180--189,
            \newblock DOI: \href{https://doi.org/10.1098/rspa.1922.0079}{10.1098/rspa.1922.0079}.

            \bibitem{Proudman1916}
            J.~Proudman,
            \newblock On the motion of solids in a liquid possessing vorticity,
            \newblock \emph{Proceedings of the Royal Society of London A} \textbf{92} (1916), 408--424,
            \newblock DOI: \href{https://doi.org/10.1098/rspa.1916.0026}{10.1098/rspa.1916.0026}.

            \bibitem{Hill1894}
            M.~J.~M. Hill,
            \newblock On a spherical vortex,
            \newblock \emph{Philosophical Transactions of the Royal Society of London A} \textbf{185} (1894), 213--245,
            \newblock DOI: \href{https://doi.org/10.1098/rsta.1894.0006}{10.1098/rsta.1894.0006}.

            \bibitem{Greenspan1968}
            H.~P. Greenspan,
            \newblock \emph{The Theory of Rotating Fluids},
            \newblock Cambridge University Press (1968),
            \newblock ISBN: 978-0521051477.

            \bibitem{Onsager1949}
            L.~Onsager,
            \newblock Statistical hydrodynamics,
            \newblock \emph{Il Nuovo Cimento} \textbf{6}(S2) (1949), 279--287,
            \newblock DOI: \href{https://doi.org/10.1007/BF02780991}{10.1007/BF02780991}.

            \bibitem{Feynman1955}
            R.~P. Feynman,
            \newblock Application of quantum mechanics to liquid helium,
            \newblock \emph{Progress in Low Temperature Physics} \textbf{1} (1955), 17--53,
            \newblock DOI: \href{https://doi.org/10.1016/S0079-6417(08)60077-3}{10.1016/S0079-6417(08)60077-3}.

            \bibitem{KlecknerIrvine2013}
            D.~Kleckner and W.~T.~M.~Irvine,
            \newblock Creation and dynamics of knotted vortices,
            \newblock \emph{Nature Physics} \textbf{9} (2013), 253--258,
            \newblock DOI: \href{https://doi.org/10.1038/nphys2560}{10.1038/nphys2560}.

            \bibitem{Jacobson2008}
            T.~Jacobson,
            \newblock Einstein-\AE ther gravity: a status report,
            \newblock PoS QG-PH (2007) 020,
            \newblock DOI: 10.22323/1.043.0020.

            \bibitem{BlasPujolasSibiryakov2011}
            D.~Blas, O.~Pujol\`as, and S.~Sibiryakov,
            \newblock Models of non-relativistic quantum gravity: the good, the bad and the healthy,
            \newblock JHEP \textbf{04} (2011) 018,
            \newblock DOI: 10.1007/JHEP04(2011)018.

            \bibitem{Monitor2017}
            B.~P. Abbott et al. (LIGO Scientific Collaboration and Virgo Collaboration),
            \newblock GW170817: Observation of gravitational waves from a binary neutron star inspiral,
            \newblock Phys. Rev. Lett. \textbf{119} (2017) 161101,
            \newblock DOI: 10.1103/PhysRevLett.119.161101.

            \bibitem{Unruh1981}
            W.~G. Unruh,
            \newblock Experimental black-hole evaporation?,
            \newblock Phys. Rev. Lett. \textbf{46} (1981), 1351--1353,
            \newblock DOI: 10.1103/PhysRevLett.46.1351.

            \bibitem{Visser1998}
            M. Visser,
            \newblock Acoustic black holes: horizons, ergospheres and Hawking radiation,
            \newblock Class. Quantum Grav. \textbf{15} (1998), 1767--1791,
            \newblock DOI: 10.1088/0264-9381/15/6/024.

            \bibitem{Painleve1921}
            P. Painlev\'e,
            \newblock La m\'ecanique classique et la th\'eorie de la relativit\'e,
            \newblock C. R. Acad. Sci. (Paris) \textbf{173} (1921), 677--680.

            \bibitem{Gullstrand1922}
            A. Gullstrand,
            \newblock Allgemeine L\"osung des statischen Eink\"orperproblems in der Einsteinschen Gravitationstheorie,
            \newblock Arkiv. Mat. Astron. Fys. \textbf{16}(8) (1922), 1--15.

            \bibitem{Schrodinger1926}
            E. Schr\"odinger,
            \newblock Quantisierung als Eigenwertproblem,
            \newblock Annalen der Physik \textbf{384} (1926), 361--376,
            \newblock DOI: 10.1002/andp.19263840404.

            \bibitem{Madelung1927}
            E. Madelung,
            \newblock Quantentheorie in hydrodynamischer Form,
            \newblock Z. Phys. \textbf{40} (1927), 322--326,
            \newblock DOI: 10.1007/BF01400372.

            \bibitem{Jackson1999}
            J.~D. Jackson,
            \newblock \emph{Classical Electrodynamics}, 3rd ed.,
            \newblock Wiley (1999).

            \bibitem{Planck1901}
            M. Planck,
            \newblock Ueber das Gesetz der Energieverteilung im Normalspectrum,
            \newblock Annalen der Physik \textbf{309} (1901), 553--563,
            \newblock DOI: 10.1002/andp.19013090310.


            \bibitem{Rolfsen1976}
            D.~Rolfsen,
            \newblock \emph{Knots and Links},
            \newblock Publish or Perish (1976); AMS Chelsea reprint (2003).

            \bibitem{Conway1970}
            J.~H.~Conway,
            \newblock An enumeration of knots and links, and some of their algebraic properties,
            \newblock in \emph{Computational Problems in Abstract Algebra}, Pergamon Press (1970), pp.~329--358,
            \newblock DOI: 10.1016/B978-0-08-012975-4.50034-5.
            \bibitem{Wien1896}
            W. Wien,
            \newblock Ueber die Energieverteilung im Emissionsspectrum eines schwarzen K\"orpers,
            \newblock Annalen der Physik \textbf{294} (1896), 662--669,
            \newblock DOI: 10.1002/andp.18962940310.

            \bibitem{Calugareanu1961}
            G. C\u{a}lug\u{a}reanu,
            \newblock Sur les classes d'isotopie des noeuds tridimensionnels et leurs invariants,
            \newblock Czech. Math. J. \textbf{11} (1961), 588--625.

            \bibitem{White1969}
            J.~H. White,
            \newblock Self-linking and the Gauss integral in higher dimensions,
            \newblock Amer. J. Math. \textbf{91} (1969), 693--728,
            \newblock DOI: 10.2307/2373348.

            \bibitem{Saffman1992}
            P.~G. Saffman,
            \newblock \emph{Vortex Dynamics},
            \newblock Cambridge University Press (1992).

            \bibitem{Zhang2005TransverseForce}
            P.-M.~Zhang, L.-M.~Cao, Y.-S.~Duan, and C.-K.~Zhong,
            \newblock Transverse force on a moving vortex with the acoustic geometry,
            \newblock arXiv:hep-th/0501073 (2005).
            \newblock Permalink: \url{https://arxiv.org/abs/hep-th/0501073}.


            \bibitem{KidaTakaoka1994}
            S.~Kida and M.~Takaoka,
            \newblock Vortex reconnection,
            \newblock Annual Review of Fluid Mechanics \textbf{26} (1994), 169--189,
            \newblock DOI: 10.1146/annurev.fl.26.010194.001125.

            \bibitem{KoplikLevine1993}
            J.~Koplik and H.~Levine,
            \newblock Vortex reconnection in superfluid helium,
            \newblock Physical Review Letters \textbf{71} (1993), 1375--1378,
            \newblock DOI: 10.1103/PhysRevLett.71.1375.

            \bibitem{BewleyPaolettiSreenivasanLathrop2008}
            G.~P. Bewley, M.~S. Paoletti, K.~R. Sreenivasan, and D.~P. Lathrop,
            \newblock Characterization of reconnecting vortices in superfluid helium,
            \newblock Proceedings of the National Academy of Sciences \textbf{105} (2008), 13707--13710,
            \newblock DOI: 10.1073/pnas.0806002105.

            \bibitem{Varoquaux2015AndersonPhaseSlippage}
            E.~Varoquaux,
            \newblock Anderson's considerations on the flow of superfluid helium: Some offshoots,
            \newblock Reviews of Modern Physics \textbf{87} (2015), 803--854,
            \newblock DOI: 10.1103/RevModPhys.87.803.

            \bibitem{Moffatt1969}
            H.~K. Moffatt,
            \newblock The degree of knottedness of tangled vortex lines,
            \newblock J. Fluid Mech. \textbf{35} (1969), 117--129,
            \newblock DOI: 10.1017/S0022112069000991.

            \bibitem{Zurek2003}
            W.~H. Zurek,
            \newblock Decoherence, einselection, and the quantum origins of the classical,
            \newblock Rev. Mod. Phys. \textbf{75} (2003), 715--775,
            \newblock DOI: 10.1103/RevModPhys.75.715.

            \bibitem{Weinberg1967}
            S.~Weinberg,
            \newblock A model of leptons,
            \newblock Phys. Rev. Lett. \textbf{19} (1967), 1264--1266,
            \newblock DOI: 10.1103/PhysRevLett.19.1264.

            \bibitem{Gibbons2002}
            G.~W. Gibbons,
            \newblock The maximum tension principle in general relativity,
            \newblock Found. Phys. \textbf{32} (2002), 1891--1901,
            \newblock DOI: 10.1023/A:1022370717626.

            \bibitem{WilczekZee1983}
            F.~Wilczek and A.~Zee,
            \newblock ``Linking Numbers, Spin, and Statistics of Solitons,''
            \newblock \textit{Physical Review Letters} \textbf{51}, 2250--2252 (1983),
            \newblock doi:10.1103/PhysRevLett.51.2250.

            \bibitem{Volovik2003}
            G.~E.~Volovik,
            \newblock \textit{The Universe in a Helium Droplet},
            \newblock International Series of Monographs on Physics, Vol.~117,
            \newblock Clarendon Press, Oxford (2003).

            \bibitem{Autti2016}
            S.~Autti, V.~V.~Dmitriev, J.~T.~M\"akinen, A.~A.~Soldatov,
            G.~E.~Volovik, A.~N.~Yudin, V.~V.~Zavjalov, and V.~B.~Eltsov,
            \newblock ``Observation of Half-Quantum Vortices in Topological Superfluid $^3$He,''
            \newblock \textit{Physical Review Letters} \textbf{117}, 255301 (2016),
            \newblock doi:10.1103/PhysRevLett.117.255301,
            \newblock arXiv:1508.02197.

            \bibitem{ChenTanizaki2023}
            S.~Chen and Y.~Tanizaki,
            \newblock ``Solitonic Symmetry beyond Homotopy:
            Invertibility from Bordism and Noninvertibility from Topological Quantum Field Theory,''
            \newblock \textit{Physical Review Letters} \textbf{131}, 011602 (2023),
            \newblock doi:10.1103/PhysRevLett.131.011602,
            \newblock arXiv:2210.13780.

            \bibitem{Einstein1905}
            A.~Einstein,
            \newblock Zur Elektrodynamik bewegter K{\"o}rper,
            \newblock Annalen der Physik \textbf{322} (1905), 891--921,
            \newblock DOI: 10.1002/andp.19053221004.

            \bibitem{Einstein1915}
            A.~Einstein,
            \newblock Die Feldgleichungen der Gravitation,
            \newblock Sitzungsberichte der K{\"o}niglich Preussischen Akademie der Wissenschaften (1915), 844--847,
            \newblock permalink: https://einsteinpapers.press.princeton.edu/vol6-trans/129.

            \bibitem{BarceloLiberatiVisser2011}
            C.~Barcel{\'o}, S.~Liberati, and M.~Visser,
            \newblock Analogue Gravity,
            \newblock Living Reviews in Relativity \textbf{14} (2011), 3,
            \newblock DOI: 10.12942/lrr-2011-3.

            \bibitem{Jacobson1995}
            T.~Jacobson,
            \newblock Thermodynamics of spacetime: The Einstein equation of state,
            \newblock Phys. Rev. Lett. \textbf{75} (1995), 1260--1263,
            \newblock DOI: 10.1103/PhysRevLett.75.1260.

            \bibitem{Sakharov1967}
            A.~D. Sakharov,
            \newblock Vacuum quantum fluctuations in curved space and the theory of gravitation,
            \newblock Dokl. Akad. Nauk SSSR \textbf{177} (1967), 70--71;
            \newblock Sov. Phys. Dokl. \textbf{12} (1968), 1040--1041.

            \bibitem{JacobsonMattingly2001}
            T.~Jacobson and D.~Mattingly,
            \newblock Gravity with a dynamical preferred frame,
            \newblock Phys. Rev. D \textbf{64} (2001), 024028,
            \newblock DOI: 10.1103/PhysRevD.64.024028.

            \bibitem{FosterJacobson2006}
            B.~Z. Foster and T.~Jacobson,
            \newblock Post-Newtonian parameters and constraints on Einstein-\AE ther theory,
            \newblock Phys. Rev. D \textbf{73} (2006), 064015,
            \newblock DOI: 10.1103/PhysRevD.73.064015.

            \bibitem{Abbott2017GW170817}
            B.~P. Abbott et al. (LIGO Scientific Collaboration and Virgo Collaboration),
            \newblock GW170817: Observation of gravitational waves from a binary neutron star inspiral,
            \newblock Phys. Rev. Lett. \textbf{119} (2017), 161101,
            \newblock DOI: 10.1103/PhysRevLett.119.161101.

            \bibitem{Mohr2025CODATA}
            P.~J. Mohr, D.~B. Newell, B.~N. Taylor, and E.~Tiesinga,
            \newblock CODATA recommended values of the fundamental physical constants: 2022,
            \newblock Rev. Mod. Phys. \textbf{97} (2025), 025002,
            \newblock DOI: 10.1103/RevModPhys.97.025002.

            \bibitem{Planck2018Cosmology}
            N.~Aghanim et al. (Planck Collaboration),
            \newblock Planck 2018 results. VI. Cosmological parameters,
            \newblock Astronomy \& Astrophysics \textbf{641} (2020), A6,
            \newblock DOI: 10.1051/0004-6361/201833910.

            \bibitem{CantarellaKusnerSullivan2002}
            J.~Cantarella, R.~B. Kusner, and J.~M. Sullivan,
            \newblock On the Minimum Ropelength of Knots and Links,
            \newblock \emph{Inventiones Mathematicae} \textbf{150} (2002), 257--286,
            \newblock DOI: \href{https://doi.org/10.1007/s00222-002-0234-y}{10.1007/s00222-002-0234-y},
            \newblock arXiv: \href{https://arxiv.org/abs/math/0103224}{math/0103224}.

            \bibitem{AshtonCantarellaPiatekRawdon2011}
            T.~Ashton, J.~Cantarella, M.~Piatek, and E.~J. Rawdon,
            \newblock Knot Tightening by Constrained Gradient Descent,
            \newblock \emph{Experimental Mathematics} \textbf{20}(1) (2011), 57--90,
            \newblock DOI: \href{https://doi.org/10.1080/10586458.2011.544581}{10.1080/10586458.2011.544581},
            \newblock arXiv: \href{https://arxiv.org/abs/1002.1723}{1002.1723}.

            \bibitem{DenneDiaoSullivan2006}
            E.~Denne, Y.~Diao, and J.~M. Sullivan,
            \newblock Quadrisecants Give New Lower Bounds for the Ropelength of a Knot,
            \newblock \emph{Geometry \& Topology} \textbf{10} (2006), 1--26,
            \newblock DOI: \href{https://doi.org/10.2140/gt.2006.10.1}{10.2140/gt.2006.10.1},
            \newblock arXiv: \href{https://arxiv.org/abs/math/0408026}{math/0408026}.

            \bibitem{BuniyKephart2003}
            R.~V. Buniy and T.~W. Kephart,
            \newblock A Model of Glueballs,
            \newblock \emph{Physics Letters B} \textbf{576} (2003), 127--134,
            \newblock DOI: \href{https://doi.org/10.1016/j.physletb.2003.09.081}{10.1016/j.physletb.2003.09.081},
            \newblock arXiv: \href{https://arxiv.org/abs/hep-ph/0209339}{hep-ph/0209339}.

            \bibitem{LIGOVirgoFermiINTEGRAL2017}
            B.~P. Abbott et al. (LIGO Scientific Collaboration, Virgo Collaboration, Fermi Gamma-ray Burst Monitor, and INTEGRAL),
            \newblock Gravitational Waves and Gamma-rays from a Binary Neutron Star Merger: GW170817 and GRB 170817A,
            \newblock \emph{Astrophysical Journal Letters} \textbf{848}, L13 (2017),
            \newblock DOI: \href{https://doi.org/10.3847/2041-8213/aa920c}{10.3847/2041-8213/aa920c}.

            \bibitem{OostMukohyamaWang2018}
            J.~Oost, S.~Mukohyama, and A.~Wang,
            \newblock Constraints on Einstein-aether theory after GW170817,
            \newblock \emph{Physical Review D} \textbf{97}, 124023 (2018),
            \newblock DOI: \href{https://doi.org/10.1103/PhysRevD.97.124023}{10.1103/PhysRevD.97.124023}.

            \bibitem{ShaoWex2012}
            L.~Shao and N.~Wex,
            \newblock New tests of local Lorentz invariance of gravity with small-eccentricity binary pulsars,
            \newblock \emph{Classical and Quantum Gravity} \textbf{29}, 215018 (2012),
            \newblock DOI: \href{https://doi.org/10.1088/0264-9381/29/21/215018}{10.1088/0264-9381/29/21/215018}.

            \bibitem{ShaoCaballeroKramerWexChampionJessner2013}
            L.~Shao, R.~N. Caballero, M.~Kramer, N.~Wex, D.~J. Champion, and A.~Jessner,
            \newblock A new limit on local Lorentz invariance violation of gravity from solitary pulsars,
            \newblock \emph{Classical and Quantum Gravity} \textbf{30}, 165019 (2013),
            \newblock DOI: \href{https://doi.org/10.1088/0264-9381/30/16/165019}{10.1088/0264-9381/30/16/165019}.

            \bibitem{GuptaHerreroValeaBlasBarausseCornishYagiYunes2021}
            T.~Gupta, M.~Herrero-Valea, D.~Blas, E.~Barausse, N.~Cornish, K.~Yagi, and N.~Yunes,
            \newblock New binary pulsar constraints on Einstein-\ae ther theory after GW170817,
            \newblock \emph{Classical and Quantum Gravity} \textbf{38}, 195003 (2021),
            \newblock DOI: \href{https://doi.org/10.1088/1361-6382/ac1a69}{10.1088/1361-6382/ac1a69}.

            \bibitem{Annala2022ProtectedVortexKnots}
            T.~Annala, R.~Zamora-Zamora, and M.~M{\"o}tt{\"o}nen,
            \newblock Topologically protected vortex knots and links,
            \newblock \emph{Communications Physics} \textbf{5} (2022), 309,
            \newblock DOI: \href{https://doi.org/10.1038/s42005-022-01071-2}{10.1038/s42005-022-01071-2}.

            \bibitem{Kleckner2016SuperfluidKnotsUntie}
            D.~Kleckner, L.~H.~Kauffman, and W.~T.~M.~Irvine,
            \newblock How superfluid vortex knots untie,
            \newblock \emph{Nature Physics} \textbf{12} (2016), 650--655,
            \newblock DOI: \href{https://doi.org/10.1038/nphys3679}{10.1038/nphys3679}.

            \bibitem{Pisanty2018FractionalKnotsLight}
            E.~Pisanty, G.~Jim{\'e}nez, V.~Vicu{\~n}a-Hern{\'a}ndez, A.~Pic{\'o}n, A.~Celi, J.~P.~Torres, and M.~Lewenstein,
            \newblock Knotting fractional-order knots with the polarization state of light,
            \newblock \emph{arXiv:1808.05193} (2018),
            \newblock arXiv: \href{https://arxiv.org/abs/1808.05193}{1808.05193}.

            \bibitem{Ricca1998KnotFluidMechanics}
            R.~L.~Ricca,
            \newblock Applications of knot theory in fluid mechanics,
            \newblock \emph{Banach Center Publications} \textbf{42} (1998), 321--346,
            \newblock Permalink: \href{https://eudml.org/doc/252224}{EuDML:252224}.

        \end{thebibliography}

\end{document}